\documentclass[11pt]{amsart}
\usepackage[a4paper,vmargin=27mm,hmargin=22.45mm]{geometry}
\usepackage[T1]{fontenc}
\usepackage{lmodern,microtype}
\usepackage{amsmath,amssymb,amsthm,mathtools,mathrsfs}
\usepackage{booktabs,longtable,array,enumitem}
\usepackage{float}
\usepackage[hidelinks]{hyperref}
\newtheorem{theorem}{Theorem}[section]
\newtheorem{proposition}[theorem]{Proposition}
\newtheorem{lemma}[theorem]{Lemma}
\newtheorem{corollary}[theorem]{Corollary}
\newtheorem{claim}[theorem]{Claim}
\theoremstyle{definition}
\newtheorem{definition}{Definition}[section]
\newtheorem*{definition*}{Definition}
\newtheorem{example}[theorem]{Example}
\theoremstyle{remark}
\newtheorem{remark}[theorem]{Remark}
\newcommand{\R}{\mathbb R}
\newcommand{\Z}{\mathbb Z}
\newcommand{\N}{\mathbb{Z}_{\ge0}}
\newcommand{\C}{\mathbb C}
\newcommand{\eps}{\varepsilon}
\newcommand{\norm}[1]{\lVert #1\rVert}
\newcommand{\abs}[1]{\lvert #1\rvert}
\newcommand{\ip}[2]{\langle #1,#2\rangle}
\DeclareMathOperator{\Sym}{Sym}
\DeclareMathOperator{\vol}{vol}
\DeclareMathOperator{\rank}{rank}
\DeclareMathOperator{\diag}{diag}
\DeclareMathOperator{\dist}{dist}
\DeclareMathOperator{\supp}{supp}
\DeclareMathOperator{\tr}{tr}
\DeclareMathOperator{\SO}{SO}
\newcommand{\smin}{s_{\min}}
\newcommand{\HS}{\mathrm{HS}}
\newcommand{\op}{\mathrm{op}}
\newcommand{\Id}{I}
\newcommand{\ind}{\mathbf1}
\numberwithin{equation}{section}
\allowdisplaybreaks[1]
\setlength{\emergencystretch}{3em}
\setcounter{tocdepth}{2}
\raggedbottom
\hypersetup{pdftitle={Sharp integrability for the three-dimensional Parsell-Vinogradov system. I. Degrees at least fourteen},pdfauthor={Jaehyun Woo}}
\title[Parsell--Vinogradov integrability I]{Sharp integrability for the three-dimensional Parsell--Vinogradov system\\I. Degrees $k\ge14$}
\author{Jaehyun Woo}
\address{Jaehyun Woo \\ Department of Mathematical Sciences, Seoul National University, Republic of Korea}
\email{wjhwsh1@snu.ac.kr}
\date{28 September 2026}


\makeatletter
\newcommand{\equationalias}[1]{\begingroup\protected@edef\@currentlabel{\p@equation\theequation}\protected@edef\@currentHref{equation.\theHequation}\label{#1}\endgroup}
\makeatother

\newcommand{\partIref}[1]{\hyperref[I:#1]{I.\ref*{I:#1}}}
\newcommand{\partIeqref}[1]{\textup{(\partIref{#1})}}
\newcommand{\partIIref}[1]{\hyperref[II:#1]{II.\ref*{II:#1}}}
\newcommand{\partIIeqref}[1]{\textup{(\partIIref{#1})}}

\begin{document}
\begin{abstract}
We study Tarry's problem for the three-dimensional Parsell--Vinogradov system, namely the integrability, over the polynomial coefficient space, of the oscillatory integral associated with all nonconstant monomials of total degree at most $k$. For every $k\ge14$, the sharp range is $p>\binom{k+3}{4}+3$. The argument uses complete affine jets, a geometric assignment, and coefficient-volume estimates retaining every nonconstant coefficient. Filtered Hermite interpolation treats the regular configurations, while transverse derivative coordinates treat the configurations near surfaces. We include a binary cubic example. The companion paper treats degrees two through thirteen using the shared estimates proved here.
\end{abstract}
\maketitle\enlargethispage{6pt}
\tableofcontents


\section{Introduction}\label{root:introduction}

\subsection{The Parsell--Vinogradov system and the main theorem}\label{intro:main}

For an integer $k\ge2$, let $V_k=\R[x,y,z]_{\le k}/\R$ be the vector space of real polynomials of total degree at most $k$, modulo constants. We equip $V_k$ with Lebesgue measure in the nonconstant monomial coefficients. Its dimension and the exponent in the affine coefficient determinant are
\[
N_k=\binom{k+3}{3}-1,\qquad \chi_k=\binom{k+3}{4}.
\]
The three-dimensional Parsell--Vinogradov system of degree $k$ consists of the equations
\[
\sum_{j=1}^s\mathbf u_j^\alpha=\sum_{j=1}^s\mathbf v_j^\alpha,\qquad \alpha\in\mathbb Z_{\ge0}^3,\quad 1\le|\alpha|\le k,
\]
where $s\ge1$ is an integer, $\mathbf u_j,\mathbf v_j\in\mathbb Z^3$, $|\alpha|=\alpha_1+\alpha_2+\alpha_3$, and $\mathbf u^\alpha=u_1^{\alpha_1}u_2^{\alpha_2}u_3^{\alpha_3}$. We study Tarry's problem for this system: the convergence of its oscillatory integral over the real coefficient space.

For the representative with constant coefficient zero, define
\[
 \mathcal I_k(F)=\int_{[0,1]^3}e^{iF(x,y,z)}\,dx\,dy\,dz.
\]
The absolute value of this integral is independent of the representative.

\begin{theorem}[Parsell--Vinogradov system]\label{root:main}
For every integer $k\ge14$ and every $0<p<\infty$,

\noindent\textup{(i)} \[ \mathcal I_k\in L^p(V_k)\quad\Longleftrightarrow\quad p>\chi_k+3. \]

\noindent\textup{(ii)} For every fixed $a\in C_c^\infty(\R^3)$, the integral $F\mapsto\int a(x)e^{iF(x)}\,dx$ also belongs to $L^p(V_k)$ when $p>\chi_k+3$, with constants depending on the fixed degree, the support, and finitely many smooth seminorms of $a$.
\end{theorem}


\subsection{Earlier work and the three-dimensional problem}\label{intro:earlier}

A discrete counterpart counts integer solutions of the Parsell--Vinogradov system in boxes. Parsell studied these multidimensional mean values in~\cite{Parsell05}. Parsell--Prendiville--Wooley obtained near-optimal bounds by efficient congruencing~\cite{PPW13}. Bourgain--Demeter used decoupling to estimate quadratic and cubic systems in two dimensions~\cite{BD16}, and Bourgain--Demeter--Guo obtained sharp bounds for the planar cubic system~\cite{BDGuo17}. Guo--Zhang proved sharp mean-value bounds, up to arbitrarily small power losses, in every dimension $d\ge2$~\cite{GZ19}. Theorem~\ref{root:main} concerns convergence of the associated integral over the full real coefficient space.

In one variable, Arkhipov--Karatsuba--Chubarikov proved that the convergence exponent is $k(k+1)/2+1$ for every $k\ge2$, using estimates for polynomial oscillatory integrals in terms of their $H$-function \cite[Theorems 1 and 5]{AKC80}. Basu--Guo--Zhang--Zorin-Kranich proved that the corresponding convergence exponent in two variables is $\binom{k+2}{3}+2$ for every $k\ge2$ \cite[Theorem 1.2]{BGZZK}. Mockenhaupt resolved the quadratic case in every dimension at least two in 1996 \cite{Moc96}; see also \cite[Section 1.3]{BGZZK}. The companion paper \cite{PaperII} treats degrees two through thirteen.

In three variables, write $p_{3,k}=\inf\{p>0:\mathcal I_k\in L^p(V_k)\}$. Arkhipov--Karatsuba--Chubarikov proved the upper bound \cite[Theorem 7]{AKC80}, and Ikromov proved the lower bound \cite[Theorem 1]{Ikr97}; see also his addendum \cite{Ikr22}. These bounds are
\[
\chi_k+1\le p_{3,k}\le3\chi_k+3,\qquad k\ge2.
\]
Theorem~\ref{root:main} determines $p_{3,k}=\chi_k+3$ for $k\ge14$.

The planar stationary-set argument uses determinant factorizations whose role becomes more difficult in higher dimension; see \cite[Remark 5.11]{BGZZK}. The proof instead constructs a common geometric assignment and bounds complete coefficient volumes at fixed physical frames. In the surface coefficient maps, the lower normal blocks restore the coefficients lost under differentiation; see Section~\ref{atlas:parent}.

The binary cubic example in Appendix~\ref{pl:section} is included to explain this organization. It recovers the known planar threshold $12$ and introduces the coupling between derivatives of different orders.

\subsection{Proof strategy and organization}\label{intro:strategy}

An affine map $u\mapsto c+Bu$ gives the polynomial $F(c+Bu)-F(c)$, whose coefficients form the complete jet. Changing $c$ mixes its homogeneous parts, whereas a linear substitution preserves each degree. The determinant of this coefficient substitution is $|\det B|^{\chi_k}$. We first choose a family of physical affine maps independently of $F$. When integrating the contribution of a label $j$, we use the coefficient coordinates of its fixed map $T_j$.

We bound the original integral by a weighted sum over the selected labels. Labels are assigned either to regular pieces, where polynomial sampling gives coefficient-volume bounds, or to pieces near a surface, where the implicit graph and its transverse derivatives give those bounds. The assignment is chosen once for all labels of a phase. The coefficient set for a fixed label includes every allowed choice of sample points; a determinant bound for one choice must therefore be combined with a bound for the number of choices.

The regular estimate of Theorem~\ref{reg:regular-tail}, obtained from filtered Hermite interpolation and nilpotent polynomial chains, has exponent $\beta_k$ with $\beta_k-\chi_k-2\ge131/513$ for $k\ge14$; see~\eqref{reg:eq:positive}. The surface estimate of Theorem~\ref{surf:common-tail} holds for $k\ge10$. In the surface case, the normal derivative defining the graph may have any of several orders. Its coefficient description is supplemented by every lower coefficient of the original phase. The local surface construction uses bounded-time comparisons along isotropic directions. The original frame remains fixed throughout the final coefficient disintegration.

For $\chi_k+3<p<\chi_k+3+1/16$, Proposition~\ref{an:unified-incidence} proves the physical incidence bound with no middle-side loss. The limiting condition in~\eqref{an:physical-positivity} is then $p-\chi_k-3>0$. Equation~\eqref{an:larger-p} extends the smooth upper bound to all larger exponents, and Lemma~\ref{an:cube-reduction} gives the cube upper bound. For sufficiently small $\epsilon>0$, Proposition~\ref{an:necessity} constructs sets in the full $N_k$-dimensional coefficient space with measure $\gtrsim_k\epsilon^{-\chi_k-3}$ on which $|\mathcal I_k(F)|\gtrsim_k\epsilon$. The factor $\epsilon^{-\chi_k}$ comes from coefficient dilation; $\epsilon^{-3}$ comes from the plane position and its two slopes. Disjoint sets at dyadic $\epsilon$ give divergence for $p\le\chi_k+3$.

Section~\ref{an:section} gives the analytic reduction, including the conditional passage from physical incidence to integrability. Section~\ref{geo:section} constructs the geometric assignment. Section~\ref{reg:section} treats regular geometry, sampling and coefficient estimates. Section~\ref{atlas:section} develops the coefficient atlases, reconstruction of the original phase and finite slope estimates. Section~\ref{surf:section} proves the surface estimate for $k\ge10$, and Section~\ref{an:assembly} completes the proof and proves sharpness. Appendix~\ref{pl:section} gives the binary cubic example, and Appendix~\ref{app:shared-estimates} collects estimates used in the companion paper \cite{PaperII}. These shared results retain their stated degree ranges. The supplement \cite{PaperSupplement} contains the detailed binary-cubic proof and finite certificate data.

\subsection*{Notation}

For nonnegative quantities, $A\lesssim B$ means $A\le CB$ for some constant $C>0$. Subscripts indicate parameter dependence: $A\lesssim_\theta B$ means $A\le C_\theta B$. We write $A\sim B$ if $A\lesssim B$ and $B\lesssim A$, and use the same convention for subscripts. The symbols $\ll$ and $\gg$ denote sufficiently small and large ratios, respectively. The conditioning parameter satisfies $P\ge2$.


\subsection*{Use of generative artificial intelligence and attribution}
This work was developed through sustained interaction with OpenAI's ChatGPT, beginning with two-variable quadratic and cubic arguments supplied by the author. The author's initial arguments used elementary linear algebra, and their frequency decomposition was motivated by Gan, Maldague, and Oh's \href{https://arxiv.org/abs/2507.09696}{\emph{Sharp local smoothing estimates for curve averages}} (2025, GMO25). The quadratic case was completed solely by the author as an alternative proof of Mockenhaupt's result. For the cubic case, the main ideas were the author's; AI completed the argument by devising an appropriate stopping condition. AI subsequently reformulated that proof in the complete affine jet framework presented in Appendix B (Corollary B.1).

The author recalls the following sequence: motivated by reduction theory, he proposed analyzing symmetries of polynomials; AI then refined the representation-theoretic arguments in response to this strategy and proposed using ideal-adic modules to address higher-degree cases. Some exploratory case-by-case arguments failed or contained gaps, leading to restarts with selected ingredients retained or reformulated. These ideas and proposals describe the development history; not all are necessarily used in the final proof.

In Section 2, AI-generated drafts developed and revised the affine determinant and matrix-logarithm calculations (Lemmas 2.1--2.2), logarithmically regularized minimizing ellipsoids and fixed physical frame catalogues (Lemmas 2.3--2.5), positive localization (Proposition 2.6), and the weighted analytic reduction (Lemmas 2.7--2.9 and 2.12; Propositions 2.10--2.11). In Section 3, AI drafted the rank-four classification (Theorem 3.5), the conditioning, normalization, preparation and packing lemmas (3.1--3.4, 3.6--3.20 and 3.22), and the common null and Gauss assignments (Theorems 3.21 and 3.23). Tests against curved parents, logarithmic cancellation and long flows led to withdrawals and reformulations of earlier routes.

For the regular coefficient estimate (Theorem 4.1), AI assisted with the development and drafting of filtered Hermite interpolation, graded Jordan chains from \(\mathfrak{sl}_2\) representations, and row and determinant bookkeeping. The supporting statements cover row selection and regular decomposition (Lemmas 4.2--4.3), conic-sum and exponent calculations (Lemmas 4.4 and 4.7), and homogeneous, affine and \(Q_t\) sampling (Propositions 4.5--4.6 and 4.8). AI-generated code tested the 22 regular regimes; finite tests supplement the written all-degree argument.

AI drafted and revised the normal filtration and complete moving-graph coefficient atlas (Lemmas 5.1--5.2), derivative-phase lifting (Proposition 5.3), lower parent-coefficient recovery and Haar-volume disintegration (Lemmas 5.4--5.6), and polynomial sampling and slope selection (Lemmas 5.7--5.9). Its code checked determinant factorizations and multiplicities used in finite-degree atlas calculations, including the degree-fourteen boundary case. In Section 6, AI-generated drafts combined wide and vertical ribbon counts, collar nets and complete Gauss coefficient atlases (Propositions 6.1--6.2, 6.4 and 6.6; Lemmas 6.3 and 6.5) into the surface coefficient estimate (Theorem 6.7); later reviews checked Jacobian factors, aspect sums and parent versus derivative-phase scalars.

AI drafted Propositions 7.1--7.2 and assembled Theorem 1.1 using analytic reduction, coefficient incidence and sharpness. It adapted Proposition B.2 and the supporting results C.1--C.12, prepared Appendix A's tables, and assisted with literature searches, notation and exposition. Cited results and classical tools retain their established attribution.

The descriptions above record AI's contributions to the development, drafting, revision, and checking of the manuscript; they do not provide a complete account of the original conception of each result.

Beyond the initial author contributions described above, the author claims no credit for the mathematical developments or proofs presented here. He directed the interactions and requested obstruction tests, corrections, and revisions, and remains solely responsible for the manuscript's mathematical claims, computations, references, and text.

\section{Complete affine jets and the analytic reduction}\label{an:section}

For the analytic estimates in this section, Appendix~\ref{pl:section}, and Section~\ref{an:assembly}, the degree is fixed before scale parameters are chosen. Constants may depend on the fixed degree, prescribed finite derivative orders, the fixed amplitude support and indicated seminorms, and displayed exponent parameters. They are independent of the phase coefficients, frequency scales, physical side ratios, and long geometric scales unless a dependence is stated.

\subsection{Complete affine jets and minimizing frames}\label{an:coefficient-changes}\label{an:jets}\label{an:normalization}

Identify $V_k$ with its zero-constant representatives and write $\|F\|$ for their Euclidean coefficient norm. Here and below, coefficient norms use a fixed basis. For a polynomial $G$ in the variables in use, $[G]$ denotes its class modulo constants, represented by $G-G(0)$. For $a\in C_c^\infty(\mathbb R^3)$ set
\[
E_a(F)=\int_{\mathbb R^3}e^{iF(x)}a(x)\,dx.
\]

\begin{definition}[Frames and complete jets]\label{def:complete-jet}
An affine frame is a map $T(u)=c+Bu$, with $B\in\operatorname{GL}(3,\mathbb R)$. Its complete nonconstant jet is
\[
\begin{gathered}
J_T(F)(u)=[F(c+Bu)]=\ell\cdot u+\sum_{r=2}^k\frac1{r!}T_r[u,\ldots,u],\\
\ell=B^T\nabla F(c),\qquad T_r=D^rF(c)[B,\ldots,B].
\end{gathered}
\]
Here $T_r$ is a symmetric $r$-tensor. These tensor coordinates and the monomial coefficients differ by a fixed linear isomorphism on their symmetric subspaces. For a vector $h$, the notation $\iota_hT_r$ means contraction in one tensor slot.
\end{definition}

Under a change of center by $Bh$ one has the formula
\begin{equation}\label{an:taylor-transport}
T_r\longmapsto\sum_{s=r}^k\frac1{(s-r)!}\iota_h^{s-r}T_s,\qquad 1\le r\le k,
\end{equation}
where $T_1=\ell$. Under a right change of frame by $A$, each tensor becomes $T_r[A,\ldots,A]$. Formula~\eqref{an:taylor-transport} follows by differentiating the finite Taylor expansion of $F$ at $c$.

\begin{lemma}[The affine coefficient determinant]\label{an:character}

\leavevmode\noindent\textup{(i)} In dimension $d\ge1$, put $V_{d,k}=\mathbb R[x_1,\ldots,x_d]_{\le k}/\mathbb R$ and $\chi_{d,k}=\binom{k+d}{d+1}$. For a fixed affine map $T(u)=c+Bu$, its coefficient pullback $\mathcal L_TF=[F\circ T]$ satisfies
\begin{equation}\label{an:character-equation}
|\det\mathcal L_T|=|\det B|^{\chi_{d,k}}.
\end{equation}

\noindent\textup{(ii)} In dimension three, $\chi_{3,k}=\chi_k$. If $v=|\det B|$, $K\ge1$, and $M$ is fixed, then
\begin{equation}\label{an:window-volume}
\operatorname{vol}\left\{F:|B^T\nabla F(c)|\le K,\quad \sum_{r=2}^k\|D^rF(c)[B,\ldots,B]\|_{\mathrm{HS}}\le M\right\}\lesssim_{k,M} K^3v^{-\chi_k}.
\end{equation}
\end{lemma}
\begin{proof}
Order coefficients by degree. Translation is triangular with diagonal entries one. For $B=\operatorname{diag}(b_1,\ldots,b_d)$, pullback on homogeneous polynomials of degree $r$ is diagonal, with entries $\prod_i b_i^{\alpha_i}$ indexed by $|\alpha|=r$. Permutation symmetry shows that each $b_i$ occurs with exponent
\[
\frac r d\binom{r+d-1}{d-1}=\binom{r+d-1}{d}.
\]
Orthogonal pullback has determinant of absolute value one, so singular-value decomposition gives the general case. Multiplying over $1\le r\le k$ gives
\[
\sum_{r=1}^k\binom{r+d-1}{d}=\binom{k+d}{d+1},
\]
which proves~\eqref{an:character-equation}. For~\eqref{an:window-volume}, change variables to $J_T(F)$. Its three linear coefficients range in a ball of radius $K$, and its other coefficients range in a fixed bounded set, by equivalence of the coefficient and tensor norms. The resulting volume is $\lesssim_{k,M}K^3$. The inverse coefficient change contributes $v^{-\chi_k}$.
\end{proof}

We next choose ellipsoids whose higher jet tensors are bounded and which satisfy the translation inequality \eqref{an:translation-stop}. These ellipsoids may depend on $F$. The coefficient integrals in Section~\ref{an:common-reduction} will instead use affine maps chosen independently of $F$.

Choose bounded open sets $X_0\Subset X$ containing $\operatorname{supp}a$ and all fixed enlargements needed below. Fix $R\ge2$, restrict for now to $\|F\|\le2R$, and put
\[
L_R=\log(C_XR),\qquad \lambda=(C_0L_R^2)^{-1}.
\]
The constants $C_X,C_0\ge2$ are fixed. For $c\in X$ and $\det B>0$, consider
\begin{equation}\label{an:energy}
\mathscr E_{F,c}(B)=\|B\|_{\mathrm{HS}}^2+\sum_{r=2}^k\|D^rF(c)[B,\ldots,B]\|_{\mathrm{HS}}^2-2\log\det B+\lambda\operatorname{tr}\bigl((\log(BB^T))^2\bigr).
\end{equation}
This depends only on the ellipsoid $B\mathbb B$, where $\mathbb B$ is the Euclidean unit ball.

\begin{lemma}[A matrix logarithm calculation]\label{an:log-hessian}
For any fixed matrix size $n$, let $A$ be positive definite and $H$ symmetric. For $A(t)=A^{1/2}e^{2tH}A^{1/2}$, in an eigenbasis $A=\operatorname{diag}(e^{a_1},\ldots,e^{a_n})$ one has
\begin{equation}\label{an:log-hessian-equation}
\left.\frac{d^2}{dt^2}\operatorname{tr}\bigl((\log A(t))^2\bigr)\right|_{t=0}=8\sum_iH_{ii}^2+16\sum_{i<j}\theta_{ij}\coth(\theta_{ij})H_{ij}^2\ge8\|H\|_{\mathrm{HS}}^2,
\end{equation}
where $\theta_{ij}=(a_i-a_j)/2$ and $\theta\coth\theta$ has its continuous value $1$ at $0$.
\end{lemma}
\begin{proof}
Apply the spectral Hessian formula in \cite[Theorem 3.3]{LewisSendov} to $\sum_i(\log s_i)^2$. Along the stated curve, $A'(0)=2A^{1/2}HA^{1/2}$ and $A''(0)=4A^{1/2}H^2A^{1/2}$; the chain rule gives \eqref{an:log-hessian-equation}. The first derivative is $4\operatorname{tr}((\log A)H)$. Continuous divided differences give the same formula at repeated eigenvalues, and $\theta\coth\theta\ge1$ gives the lower bound.
\end{proof}

\begin{lemma}[Minimizing ellipsoids and augmented translation injectivity]\label{an:minimizer}

\leavevmode\noindent\textup{(i)} The functional~\eqref{an:energy} has a unique minimizing ellipsoid. Write $B(F,c)$ for a representing frame with positive determinant, understood up to right multiplication by $\operatorname{SO}(3)$. Its positive-definite matrix $BB^T$ depends smoothly on $(F,c)$.

\noindent\textup{(ii)} At a minimizing frame, let $G_r$ be the contraction Gram matrix characterized by $\langle G_ru,u\rangle=\|\iota_uT_r\|_{\mathrm{HS}}^2$. Then
\begin{equation}\label{an:balance}
B^TB+\sum_{r=2}^krG_r+2\lambda\log(B^TB)=I.
\end{equation}

\noindent\textup{(iii)} There are constants $c_*,C_*>0$, independent of $R,F,c$, such that
\begin{equation}\label{an:frame-bounds}
c_*R^{-1}\le s_{\min}(B)\le\|B\|_{\mathrm{op}}\le1,\qquad \sum_{r=2}^k\|T_r\|_{\mathrm{HS}}\le C_*.
\end{equation}

\noindent\textup{(iv)} For $J=[F\circ(c+B\cdot)]$, the augmented translation map satisfies
\begin{equation}\label{an:translation-stop}
\|[u\cdot\nabla J]\|^2+|Bu|^2\ge c_*|u|^2\qquad(u\in\mathbb R^3).
\end{equation}
\end{lemma}
\begin{proof}
Right orthogonal multiplication leaves every term of~\eqref{an:energy} unchanged. For singular values $s_1,s_2,s_3$, the lower bound $\sum_i(s_i^2-2\log s_i)$ makes this energy coercive on the space of positive-definite ellipsoids, so a minimum exists.

For a symmetric matrix $H$, diagonalize $H$ orthogonally. Along $Be^{tH}$, each squared tensor norm and the squared matrix norm is a sum of nonnegative multiples of exponentials in $t$, and the determinant term is affine. Lemma~\ref{an:log-hessian}, applied at each point of this path and using polar decomposition, shows that the last term is strictly convex when $H\ne0$. Every other positive-determinant matrix is $Be^HO$ for a symmetric $H$ and an orthogonal $O$: this follows from the polar decomposition of $B^{-1}B'$. Strict convexity gives uniqueness of the ellipsoid. The Hessian in symmetric frame directions is positive definite; the finite-dimensional implicit function theorem, applied in a local symmetric-coordinate chart on the positive-definite matrices, gives smooth dependence.

For every symmetric $H$, stationarity gives
\[
0=\frac12\left.\frac{d}{dt}\mathscr E_{F,c}(Be^{tH})\right|_{t=0}=\left\langle B^TB+\sum_{r=2}^k rG_r-I+2\lambda\log(B^TB),H\right\rangle.
\]
The logarithm contribution follows from the first derivative formula in Lemma~\ref{an:log-hessian}. Vanishing for all symmetric $H$ proves~\eqref{an:balance}.

Let $e$ be a unit eigenvector of $B^TB$, with eigenvalue $d^2$. Since $G_r$ is positive semidefinite,~\eqref{an:balance} excludes $d>1$. The raw tensors $D^rF(c)$ are bounded by $C_XR$ on $X$. Hence, using $|Be|=d$ and $\|B\|\le1$, one obtains $\|\iota_eT_r\|\le C_XRd$. Taking the quadratic form of~\eqref{an:balance} on $e$ therefore gives
\[
1\le1-2\lambda\log d^2=d^2+\sum_{r=2}^kr\|\iota_eT_r\|^2\le(1+C_XR^2)d^2.
\]
This proves $d\ge c_*R^{-1}$. It also implies $\|\log(B^TB)\|\le C_XL_R$. Taking the trace of~\eqref{an:balance}, and observing that $\operatorname{tr}G_r=\|T_r\|_{\mathrm{HS}}^2$, proves the uniform tensor bound because $\lambda L_R$ is bounded.

All eigenvalues of $B^TB$ are at most one, so $\log(B^TB)$ is negative semidefinite. Thus~\eqref{an:balance} gives
\[
B^TB+\sum_{r=2}^krG_r\succeq I.
\]
The nonconstant coefficients of $u\cdot\nabla J$ are, up to the fixed factorial factors, exactly the contractions $\iota_uT_r$ for $r\ge2$. Equivalence of norms proves~\eqref{an:translation-stop}. Under a bounded right change $A$, replace $u$ by $Au$ in this inequality and apply the bounded invertible tensor maps in the remaining slots. Under a normalized translation $h$, formula~\eqref{an:taylor-transport} changes the higher tensors by $O(|h|)$ on the bounded-jet window. The quadratic form in~\eqref{an:translation-stop} therefore changes by $O(|h|)|u|^2$, so a sufficiently small fixed $|h|$ preserves its lower bound.
\end{proof}

The physical term in~\eqref{an:translation-stop} is necessary. For example, if $F$ depends only on $x_1$, then $[u\cdot\nabla F]=0$ for every $u$ in the span of the other two coordinate vectors, even when the physical frame is invertible.

\begin{lemma}[Local comparison of ellipsoids]\label{an:comparison}

\leavevmode\noindent\textup{(i)} There is a fixed $c_{\mathrm{cmp}}>0$ with the following property. If $B_0$ represents the minimizing ellipsoid at $c$ and $|B_0^{-1}(c'-c)|\le c_{\mathrm{cmp}}\sqrt\lambda$, representatives can be chosen so that
\begin{equation}\label{an:relative-comparison}
\|B_0^{-1}B(F,c')-I\|\lesssim \lambda^{-1/2}|B_0^{-1}(c'-c)|.
\end{equation}

\noindent\textup{(ii)} Fix any finite dilation factor $A_0$. For $\rho=c_0\sqrt\lambda$ with $c_0$ sufficiently small depending on $A_0$, intersecting $A_0$-dilates of $c+\rho B(F,c)\mathbb B$ have relatively comparable matrices and volumes, independently of their side ratios.
\end{lemma}
\begin{proof}
Let $L_HT_r$ denote the sum of the actions of the symmetric matrix $H$ in all $r$ slots. At a critical frame the Hessian $\mathcal H$ satisfies
\[
\mathcal H(H,H)\ge4\sum_{r=2}^k\|L_HT_r\|^2+8\lambda\|H\|^2.
\]
For the tensor term, this is the identity $\frac{d^2}{dt^2}\|e^{tL_H}T_r\|^2|_0=4\|L_HT_r\|^2$, since $L_H$ is self-adjoint; the remaining positive contribution follows from Lemma~\ref{an:log-hessian}. With $B$ fixed,
\[
\left.\frac{d}{dt}D^rF(c+tBh)[B,\ldots,B]\right|_{t=0}=\begin{cases}\iota_hT_{r+1},&2\le r<k,\\0,&r=k.\end{cases}
\]
The mixed center--frame derivative of the energy is consequently
\[
4\sum_{r=2}^{k-1}\langle\iota_hT_{r+1},L_HT_r\rangle.
\]
Choose smooth representing frames along the path of minimizing ellipsoids so that $\dot B=BS$ with $S$ symmetric. Such a choice exists: if $Q=BB^T$, the equation $S=\tfrac12B^{-1}\dot QB^{-T}$ defines a symmetric $S$ and gives $\frac{d}{dt}(BB^T)=\dot Q$. Differentiating criticality and testing it against $S$ gives
\[
\mathcal H(S,S)=-4\sum_{r=2}^{k-1}\langle\iota_hT_{r+1},L_ST_r\rangle.
\]
The bounded-jet estimate and Cauchy--Schwarz give an absolute bound $\lesssim |h|\mathcal H(S,S)^{1/2}$ for the right side. Thus $\|S\|\lesssim \lambda^{-1/2}|h|$.

Write $c(t)=c+tB_0u$ and $A(t)=B_0^{-1}B(F,c(t))$, starting with $A(0)=I$. Then $h(t)=A(t)^{-1}u$ and $\dot A=AS$. As long as $\|A\|,\|A^{-1}\|\le2$, the preceding bound gives $\|\dot A\|+\|\frac{d}{dt}A^{-1}\|\lesssim \lambda^{-1/2}|u|$. Integration and continuity show that these two norms remain at most two for $0\le t\le1$ when $|u|\le c_{\mathrm{cmp}}\sqrt\lambda$. A further integration gives~\eqref{an:relative-comparison}. If two dilated ellipsoids intersect at $x$, apply this result separately from each center to $x$. The displacement from either center is at most $A_0\rho$ in its normalized coordinates. Choosing $A_0c_0$ sufficiently small makes both comparisons applicable. Their composition gives the claimed uniform relative comparison.
\end{proof}

\subsection{Physical localization and a weighted mass bound}\label{an:physical-resource}\label{an:localization}

Fix a sufficiently small constant $\vartheta>0$. A coarse shape is an ordered dyadic triple $D=(d_1,d_2,d_3)$ with $0<d_1\le d_2\le d_3\le1$; write $v_D=d_1d_2d_3$. In each coarse shape, use a multiplicative $(1+\vartheta)$ mesh of singular values. There are $\lesssim\vartheta^{-3}$ fine triples per coarse shape. Retain lengths between $cR^{-1}$ and a fixed upper bound, with $c$ small enough to cover the lower bound in~\eqref{an:frame-bounds}. All choices in the family below are made before $F$ is specified. The labels $j$, frames $T_j$, cutoffs $\phi_j$, and mass weights \eqref{an:mass-weights} are coefficient-independent. For $T_j(u)=c_j+B_ju$, a relative affine neighborhood of $T_i$ with bound $M$ consists of the frames satisfying
\[
\max\{\|B_i^{-1}B_j\|_{\mathrm{op}},\|(B_i^{-1}B_j)^{-1}\|_{\mathrm{op}},|B_i^{-1}(c_j-c_i)|\}\le M.
\]
A bounded right comparison means replacing $T_j$ by $T_j\circ r$, with the homogeneous affine matrices of $r$ and $r^{-1}$ bounded; closeness to the identity will be specified separately.

\begin{lemma}[A finite affine family]\label{an:catalogue}

\leavevmode\noindent\textup{(i)} For a fine diagonal $E=\operatorname{diag}(e_1,e_2,e_3)$ with ordered entries, there is a set of orientations of cardinality
\begin{equation}\label{an:orientation-number}
\lesssim\vartheta^{-3}\prod_{i<j}\frac{e_j}{e_i}=\vartheta^{-3}(e_3/e_1)^2
\end{equation}
that approximates every ellipsoid with these singular values to relative matrix error $C\vartheta$, after right orthogonal adjustment.

\noindent\textup{(ii)} For each resulting frame $B$, the centers are the points of $c\vartheta\rho B\mathbb Z^3$ in a fixed bounded enlargement of $X_0$. Fix $0<r_0<r_1<r_2$ with the dilation ratios required by the covering argument below. The core, plateau, and support ellipsoid of index $j$ are $c_j+r_0\rho B_j\mathbb B$, $c_j+r_1\rho B_j\mathbb B$, and $c_j+r_2\rho B_j\mathbb B$, respectively. Choose once a function $\phi\in C_c^\infty(\mathbb R^3)$ with $0\le\phi\le1$, $\phi=1$ on $r_1\mathbb B$, and $\operatorname{supp}\phi\subset r_2\mathbb B$. Set
\[
\phi_j(x)=\phi\!\left(\rho^{-1}B_j^{-1}(x-c_j)\right).
\]
The full family $\mathcal J_D$ is finite. The number of its indices in any relative affine neighborhood with matrix, inverse, and translation bounds $M\ge2$ is at most
\begin{equation}\label{an:local-packing}
\Lambda M^C,\qquad \Lambda\lesssim \vartheta^{-15}\rho^{-3}.
\end{equation}
\end{lemma}
\begin{proof}
For $0<t<1/10$ define $U_E(t)=\{O\in SO(3):\|E^{-1}(O-I)E\|<t\}$. Thus $U_E(t)U_E(t)\subset U_E(2t+t^2)$ and $U_E(t)^{-1}\subset U_E(t/(1-t))$. The matrix power series for logarithm and exponential commute with conjugation by $E$. They identify $U_E(t)$, up to fixed changes in $t$, with skew-symmetric matrices $H$ satisfying $\|E^{-1}HE\|<t$. For a skew-symmetric $H$, this last condition is equivalent, up to dimension-dependent constants, to
\[
|H_{ij}|\le t\min(e_i/e_j,e_j/e_i)\qquad(i<j).
\]
The Haar density in these logarithm coordinates is bounded above and below in a fixed neighborhood of zero. Thus $\mu(U_E(t))\sim t^3\prod_{i<j}e_i/e_j$, where $\mu$ is Haar measure on $SO(3)$. A maximal family of disjoint left translates of $U_E(c\vartheta)$ has cardinality bounded by the reciprocal volume; maximality and the product and inverse inclusions imply that fixed enlargements cover $SO(3)$. This proves~\eqref{an:orientation-number} and the relative approximation property. Choosing a fixed separated maximal set once for each shape fixes the family independently of coefficients.

To verify~\eqref{an:local-packing}, fix a reference frame $B_*$ and a fine shape. For each orientation label in the relative neighborhood put $C_j=B_*^{-1}B_j$. The entries lie in a box of side $O(M)$ and $\|C_j^{-1}\|\le M$. Divide this box into cubes of side $c\vartheta/M$. Two matrices in the same cube would satisfy $\|C_i^{-1}C_j-I\|<c'\vartheta$, contradicting separation of the orientation set when $c$ is chosen sufficiently small. For each orientation the center lattice has normalized spacing $c\vartheta\rho$. Consequently,
\[
\begin{aligned}
\#\{\text{matrix cubes}\}&\lesssim M^{18}\vartheta^{-9},\\
\#\{\text{centers for one orientation in the translation box}\}&\lesssim M^C(\vartheta\rho)^{-3}.
\end{aligned}
\]
Multiplication by the $\lesssim\vartheta^{-3}$ fine shapes gives~\eqref{an:local-packing}. The physical centers lie in a bounded region and all singular values have a positive lower bound for fixed $R$, so every list is finite.
\end{proof}

\paragraph{Selection and partition of unity.} Fix $\delta_{\rm sel}=C_{\rm sel}\vartheta<1/2$, where $C_{\rm sel}$ accommodates the frame and center approximations below and is chosen before $\vartheta$. Test each index $j$ by
\[
(1-\delta_{\rm sel})^2|B(F,c_j)^Tu|^2\le |B_j^Tu|^2\le(1+\delta_{\rm sel})^2|B(F,c_j)^Tu|^2\qquad(u\in\mathbb R^3).
\]
The selected subfamily $\mathcal I(F)$ is obtained from the indices passing this test by the disjoint-core procedure in the proof.
\begin{proposition}[Localization by a partition of unity]\label{an:resolution}
For every $\|F\|\le2R$ there is a selected subcollection $\mathcal I(F)$ of the fixed family with the following properties.

\noindent\textup{(i)} For each fixed $j$, the set $\{F\in V_k:\|F\|\le2R,\ j\in\mathcal I(F)\}$ is Borel measurable.

\noindent\textup{(ii)} There are functions $a_{F,j}\in C_c^\infty(\mathbb R^3)$ satisfying
\begin{equation}\label{an:positive-reconstruction}
a=\sum_{j\in\mathcal I(F)}a_{F,j},\qquad \sum_{j\in\mathcal I(F)}|a_{F,j}|=|a|,\qquad |a_{F,j}|\le C|a|\phi_j.
\end{equation}

\noindent\textup{(iii)} Every selected frame has uniformly bounded higher jets and satisfies~\eqref{an:translation-stop}.

\noindent\textup{(iv)} Its cutoff supports have bounded overlap, and, with $v_j=|\det B_j|$,
\begin{equation}\label{an:resolution-bounds}
\sum_{j\in\mathcal I(F)}\rho^3v_j\lesssim_X 1,\qquad \|a_{F,j}(c_j+\rho B_j\cdot)\|_{C^m}\lesssim_{m,a} 1
\end{equation}
for every nonnegative integer $m$. No derivative with respect to $F$ is required.
\end{proposition}
\begin{proof}
For each fixed $j$, the set of coefficient vectors passing the displayed test is closed: Lemma~\ref{an:minimizer} gives continuous dependence of $B(F,c_j)B(F,c_j)^T$ on $F$. Near any point of $X_0$, round its minimizing ellipsoid to the fine lengths and orientation mesh, and then choose a nearest point of the associated center lattice. The displacement is $O(\vartheta\rho)$ in normalized coordinates. Lemma~\ref{an:comparison} changes the exact ellipsoid by only $O(\vartheta)$ under this displacement. Thus a fixed sufficiently large $C_{\rm sel}$ ensures that an index passing the test has a supporting ellipsoid covering that point.

Use a fixed ordering of the entire finite family. Skip indices that fail the matrix test, and retain each remaining index if its core is disjoint from the previously retained cores. Intersecting support ellipsoids of indices passing the test have uniformly comparable frames by Lemma~\ref{an:comparison}; fixed dilations of the retained cores therefore cover $X_0$ and have bounded overlap by volume comparison. Choose the fixed plateau ratio $r_1/r_0$ to contain these dilations. Consequently,
\[
S(F,x):=\sum_{j\in\mathcal I(F)}\phi_j(x),\qquad 1\le S(F,x)\lesssim1\quad(x\in X_0),
\]
and the bounded physical region gives $\sum_j\rho^3v_j\lesssim_X1$. Membership of a fixed index in $\mathcal I(F)$ is determined by finitely many Boolean combinations of the closed test sets, so it is Borel measurable.

For every index of the fixed family, define
\[
a_{F,j}(x)=
\begin{cases}
\mathbf1_{\{j\in\mathcal I(F)\}}\,a(x)\phi_j(x)/S(F,x),&x\in X_0,\\
0,&x\notin X_0.
\end{cases}
\]
These functions are Borel measurable in $(F,x)$ and smooth in $x$ for each fixed $F$. Positivity gives~\eqref{an:positive-reconstruction}. Relative-frame comparison, bounded overlap, and $S(F,x)\ge1$ give the scaled derivative bounds in~\eqref{an:resolution-bounds}, using $\rho\|B_j\|\lesssim1$ and $\operatorname{supp}a\Subset X_0$. The jet and injectivity assertions follow from Lemma~\ref{an:minimizer} by bounded right comparison.
\end{proof}

Define fixed nonnegative mass weights by
\begin{equation}\label{an:mass-weights}
\omega_j=\frac{C}{v_j}\int|a(x)|\phi_j(x)\,dx.
\end{equation}
Choose the constant so that $|\int e^{iF}a_{F,j}|\le v_j\omega_j$ for selected indices. Since $\phi_j$ is supported on a set of volume $O(\rho^3v_j)$, one has $0\le\omega_j\lesssim_a\rho^3\le1$.

\begin{lemma}[Weighted mass bound]\label{an:resource}
For $p\ge1$ and each coarse shape,
\begin{equation}\label{an:resource-equation}
\sum_{j\in\mathcal J_D}\omega_j^pv_j^{p-\chi_k}\lesssim_{p,a} \vartheta^{-9}d_1^{p-\chi_k-3}d_2^{p-\chi_k-1}d_3^{p-\chi_k+1}.
\end{equation}
The sum is over the fixed whole family, before selection.
\end{lemma}
\begin{proof}
For a fixed fine frame $B$, bounded overlap of its lattice bumps gives, with $v=|\det B|$,
\[
\sum_{c_j}\omega_j\lesssim_a \vartheta^{-3}v^{-1}.
\]
Use $\omega_j^p\lesssim_{p,a} \omega_j$, multiply by $v^{p-\chi_k}$, and sum over the orientations and fine shapes. Lemma~\ref{an:catalogue} gives an orientation and fine-shape count $\lesssim\vartheta^{-6}(d_3/d_1)^2$, so the sum is
\[
\lesssim_{p,a}\vartheta^{-9}v_D^{p-\chi_k-1}(d_3/d_1)^2,
\]
which is~\eqref{an:resource-equation}.
\end{proof}

\begin{lemma}[The large linear term]\label{an:linear-tail}
Write $\ell_j(F)=B_j^T\nabla F(c_j)$, $T_{r,j}=D^rF(c_j)[B_j,\ldots,B_j]$, and $I_j(F)=\int e^{iF}a_{F,j}$ for selected $j$, and put $I_j=0$ otherwise. For $K_{\mathrm{lin}}$ larger than a fixed constant, let
\[
E_D(F)=\sum_{j\in\mathcal J_D,\ |\ell_j(F)|\le K_{\mathrm{lin}}}I_j(F).
\]
For each fixed integer $m\ge1$,
\begin{equation}\label{an:linear-tail-equation}
\left|E_a(F)-\sum_D E_D(F)\right|\lesssim_{m,a} (1+\rho K_{\mathrm{lin}})^{-m}\qquad(\|F\|\le2R).
\end{equation}
\end{lemma}
\begin{proof}
In the variables $x=c_j+\rho B_jz$, the phase modulo a constant is
\[
\Phi_j(z)=\rho\ell_j\cdot z+\sum_{r=2}^k\frac{\rho^r}{r!}T_{r,j}[z,\ldots,z].
\]
The nonlinear part and its derivatives are $O(\rho^2)$ on the fixed support, while~\eqref{an:resolution-bounds} controls the normalized amplitude. For $e=\ell_j/|\ell_j|$ and $|\ell_j|$ above a fixed constant, $|\partial_e\Phi_j|\ge\rho|\ell_j|/2$. The standard nonstationary-phase estimate therefore gives
\[
|I_j(F)|\lesssim_{m,a} \rho^3v_j(1+\rho|\ell_j|)^{-m}.
\]
Summing over $|\ell_j|>K_{\mathrm{lin}}$ and using~\eqref{an:resolution-bounds} and~\eqref{an:positive-reconstruction} proves the claim.
\end{proof}

\subsection{Reduction to weighted incidence}\label{an:common-reduction}\label{an:incidence-reduction}

The next fixed change of coordinates is useful for the geometric arguments. Write $\mathcal U$ for the six-dimensional group of affine maps with upper unitriangular linear part
\begin{equation}\label{geo:group}
g(a,b,c,d,e,f)(u_1,u_2,u_3)=(u_1+au_2+bu_3+c,\ u_2+du_3+e,\ u_3+f).
\end{equation}
\equationalias{an:upper-affine-group}
The symbols $a,b,c,d,e,f$ in this display are group coordinates, unrelated to the amplitude $a$ elsewhere. Every member of $\mathcal U$ has coefficient determinant one by Lemma~\ref{an:character}.

\begin{lemma}[Triangular affine coordinates]\label{an:rq}

\leavevmode\noindent\textup{(i)} In a coarse dyadic shape family, the family admits a partition into $\lesssim(1+\log(d_3/d_1))^3$ bins. Within each bin there are a fixed positive diagonal matrix $D_\ell$ and coefficient-independent maps $g_j\in\mathcal U$, $A_j\in\operatorname{GL}(3,\mathbb R)$ such that
\begin{equation}\label{an:rq-factorization}
T_j=D_\ell\circ g_j\circ A_j,\qquad \|A_j\|,\|A_j^{-1}\|\lesssim 1.
\end{equation}

\noindent\textup{(ii)} The representatives $T_j^0=D_\ell\circ g_j=T_j\circ A_j^{-1}$ have the same centers and comparable singular values. Bounded jet windows, augmented translation injectivity, and relative packing persist under this replacement.

\noindent\textup{(iii)} Coefficient integration at $T_j^0$ contributes a factor comparable to $v_j^{-\chi_k}$.
\end{lemma}
\begin{proof}
Use RQ factorization $B_j=R_jO_j$, where $R_j$ is upper triangular with positive diagonal and $O_j$ is orthogonal. If $r_{ji}=(R_j)_{ii}$, then $r_{ji}\le\|R_j\|=\|B_j\|$ and $r_{ji}^{-1}\le\|R_j^{-1}\|=\|B_j^{-1}\|$. Each $r_{ji}$ therefore lies between the least and greatest singular values. Put $j$ in the unique dyadic triple bin $\ell_i\le r_{ji}<2\ell_i$, write $D_\ell=\operatorname{diag}(\ell_1,\ell_2,\ell_3)$, and set
\[
\Delta_j=\operatorname{diag}(r_{j1}/\ell_1,r_{j2}/\ell_2,r_{j3}/\ell_3),\quad U_j=D_\ell^{-1}R_j\Delta_j^{-1},\quad A_j=\Delta_jO_j.
\]
Then $U_j$ is upper unitriangular and $B_j=D_\ell U_jA_j$. Define $g_j(u)=U_ju+D_\ell^{-1}c_j$. This proves~\eqref{an:rq-factorization}. There are $O((1+\log(d_3/d_1))^3)$ possible triples, and $\Delta_j,\Delta_j^{-1}$ have uniformly bounded norms.

The jet and translation inequalities persist by the exact bounded frame changes already proved. For packing, if $g_i^{-1}g_j$ lies in a bounded relative affine neighborhood, then
\[
T_i^{-1}T_j=A_i^{-1}g_i^{-1}g_jA_j
\]
lies in a fixed enlarged such neighborhood. Thus the packing estimate applies. Finally, $|\det DT_j^0|=v_j/|\det A_j|$ and $1\le|\det A_j|<8$, so Lemma~\ref{an:character} gives the coefficient factor. The singular values that enter subsequent physical estimates are those of $DT_j^0$, comparable to $d_i$; the entries $\ell_i$ need not be ordered and are not substituted for them.
\end{proof}

At $T_j^0$ the transferred conditions are the bounded-jet and augmented translation inequalities. Subsequent rank tests, velocity tests and branch choices use the fixed presentation~\eqref{an:rq-factorization}.

\begin{definition}[Common assignment]\label{def:assignment}
\leavevmode\noindent\textup{(i)} For each phase $F$, let $\mathcal A_j(F)$ denote the condition that the fixed index $j$ is admissible. A common assignment assigns each admissible index to exactly one piece of a family $\mathscr P(F)$ whose possible labels belong to a fixed countable set. For each fixed original index and piece label, the set of coefficient vectors for which that assignment is made is measurable; the scales are measurable functions of $F$ on the active pieces. The assignment is chosen for all indices at once and is not changed when a particular $j$ is fixed for coefficient integration.

\noindent\textup{(ii)} Unless another measure is specified, $E_j(\cdots)$ is the set of all $J=\mathcal L_{T_j}F$ for which $j$ is assigned to a piece with the indicated scales or type. It includes every such phase and every allowed choice made in constructing that piece.

\noindent\textup{(iii)} Its preimage in the original coefficient variable is $\mathcal L_{T_j}^{-1}E_j(\cdots)$, of volume $v_j^{-\chi_k}\operatorname{vol}E_j(\cdots)$. The frame $T_j$ and weight of $j$ are kept fixed. Lemma~\ref{an:rq} permits the fixed representative $T_j^0$ in place of $T_j$.
\end{definition}
A cover of fixed finite multiplicity gives an assignment by choosing the first admissible piece in a fixed measurable ordering.

\begin{proposition}[Weighted incidence at fixed frames]\label{an:common-incidence}
Let $p>1$, $\eta>0$, and $q=(1+\eta)(p-1)$. Suppose that an assignment has scales $s_I\ge1$ and that, for every $F$,
\begin{equation}\label{an:piece-hypotheses}
\#\{j:j\text{ is assigned to }I\}\le C_0\Lambda s_I,\qquad \sum_{I\in\mathscr P(F)}s_I^{-\eta}\le Q.
\end{equation}
For each fixed original label $j$, change coefficients by $J=\mathcal L_{T_j}F$ and assume
\begin{equation}\label{an:anchor-hypothesis}
\int_{V_k}\sum_{I\ni j}s_I(\mathcal L_{T_j}^{-1}J)^q\mathbf1_{\{j\text{ assigned to }I\}}(\mathcal L_{T_j}^{-1}J)\,dJ\le M.
\end{equation}

\noindent\textup{(i)} Then, for arbitrary fixed numbers $b_j\ge0$,
\begin{equation}\label{an:common-incidence-equation}
\int_{V_k}\left(\sum_j b_jv_j\mathbf1_{\mathcal A_j(F)}\right)^p\,dF\le Q^{p-1}(C_0\Lambda)^{p-1}M\sum_jb_j^pv_j^{p-\chi_k}.
\end{equation}

\noindent\textup{(ii)} The same conclusion holds with a fixed bounded representative $T_j^0$ in~\eqref{an:anchor-hypothesis}, at the cost of a fixed factor. The constants $C_0,Q,M$ may contain explicitly stated physical-side factors; such factors remain in~\eqref{an:common-incidence-equation}.
\end{proposition}
\begin{proof}
Put $w_I=s_I^{\eta(p-1)}$. Two applications of H\"older, first over pieces and then within a piece with count $n_I$, give
\[
\begin{aligned}
\left(\sum_I\sum_{j\in I}b_jv_j\right)^p&\le\left(\sum_Iw_I^{-1/(p-1)}\right)^{p-1}\sum_Iw_In_I^{p-1}\sum_{j\in I}(b_jv_j)^p\\
&\le Q^{p-1}(C_0\Lambda)^{p-1}\sum_j(b_jv_j)^p\sum_{I\ni j}s_I^q\mathbf1_{\{j\text{ assigned to }I\}}.
\end{aligned}
\]
Integrate by Tonelli and use $dF=v_j^{-\chi_k}dJ$ at each fixed original label. Hypothesis~\eqref{an:anchor-hypothesis} gives~\eqref{an:common-incidence-equation}; monotone convergence handles infinite collections. Lemma~\ref{an:rq} gives the bounded-representative version.
\end{proof}

The constant $M$ in \eqref{an:anchor-hypothesis} can be estimated by splitting the assigned coefficient set into dyadic ranges. For regular pieces of length $H\ge1$, a bound $\operatorname{vol}E_j(H)\lesssim H^{-\beta}$ gives $M\lesssim\sum_H H^{q-\beta}$. For surface pieces of area $S\ge1$ and aspect $1\le K\le S$, a bound $\operatorname{vol}E_j(S,K)\lesssim S^{-\sigma}K^{-\tau}$ gives $M\lesssim\sum_{S,K}S^{q-\sigma}K^{-\tau}$. These are bounds in the normalized coefficient variable. The first sum converges if $q<\beta$; the second converges if $q<\sigma$ and $\tau>0$. If the regular bound holds only for $1\le H\le CD_*^{-1}$ and $q>\beta$, the first sum is $\lesssim D_*^{-(q-\beta)}$.

The selected physical labels in Proposition~\ref{an:resolution} need not form a semialgebraic set in $F$. For the selected low-linear labels it is enough to prove
\[
\{j\in\mathcal I(F):|\ell_j(F)|\le K_{\rm lin}\}\subseteq\{j:\mathcal A_j(F)\}
\]
and apply Proposition~\ref{an:common-incidence} on the set on the right. Positivity then bounds the selected sum by the admissibility sum. The geometric arguments below construct the pieces. The coefficient estimates bound their complete coefficient volumes at fixed frames, with polynomial dependence on the complete-jet window and all reciprocal pivots. This polynomial dependence is used in the annular argument in Section~\ref{an:assembly-section}.


\subsection{Physical scales and annular summation}\label{completion:physical}\label{an:physical-summation}

\paragraph{The physical incidence hypothesis.}

Fix a degree $k$, an exponent $p>1$, and $\gamma\ge0$. In each coarse shape $D$, physical reconstruction requires
\begin{equation}\label{an:physical-incidence}
\int_{\mathcal A_R}\left(\sum_{j\in\mathcal J_D}b_jv_j\mathbf1_{\{j\in\mathcal I(F),\ |\ell_j(F)|\le K_{\mathrm{lin}}\}}\right)^p\,dF\lesssim_p K_{\mathrm{lin}}^{A_p}L_R^{B_p}d_2^{-\gamma}\sum_{j\in\mathcal J_D}b_j^pv_j^{p-\chi_k},
\end{equation}
\equationalias{low:4:pf:incidence-input}\equationalias{low:5:qpf:incidence-input}\equationalias{low:6:s36pf:incidence-input}\equationalias{low:6:s36:incidence-bound}
for all fixed $b_j\ge0$, where $\mathcal A_R=\{F:R\le\|F\|\le2R\}$. The physical labels and their positive reconstruction are those of Section~\ref{an:section}. An enlarged admissibility condition is permissible, since every weight is nonnegative.

In the physical estimates write $K=K_{\mathrm{lin}}$ and use the notation of Lemma~\ref{an:linear-tail}. Its estimate uses only physical derivatives of the amplitudes. Fix the polynomial incidence exponents in $K,L_R$ before choosing the integration-by-parts order.

\begin{proposition}[Analytic implication of the physical incidence bound]\label{an:physical-assembly}
Fix $k\ge2$ and $p>1$. Suppose~\eqref{an:physical-incidence} holds for some finite $A_p,B_p$ and $\gamma\ge0$, where $d_2^{-\gamma}$ is the displayed middle-side factor. If
\begin{equation}\label{an:physical-positivity}
\mu_1=p-\chi_k-3>0,\qquad \mu_2=p-\chi_k-1-\gamma>0,\qquad \mu_3=p-\chi_k+1>0,
\end{equation}
the following assertions hold.

\noindent\textup{(i)} $E_a\in L^p(V_k)$ for every $a\in C_c^\infty(\mathbb R^3)$.

\noindent\textup{(ii)} On a fixed compact support its norm is controlled by finitely many derivatives of $a$.
\end{proposition}
\begin{proof}
With the fixed weights $b_j=\omega_j$, positive reconstruction, \eqref{an:physical-incidence}, and \eqref{an:resource-equation} give
\begin{equation}\label{an:shape-estimate}
\|E_D\|_{L^p(\mathcal A_R)}^p\lesssim_{p,a} K_{\mathrm{lin}}^{A_p}L_R^{B_p}d_1^{\mu_1}d_2^{\mu_2}d_3^{\mu_3}.
\end{equation}
\equationalias{low:4:pf:final-shape}\equationalias{low:5:qpf:final-shape}

For a contributing shape, choose a selected label in its low-linear coefficient set. Applying $B_j^{-1}$ in all tensor slots bounds the raw derivatives at its bounded center by $CK_{\mathrm{lin}}d_1^{-1}$ at order one and $Cd_1^{-r}$ at order $r\ge2$. Inverse Taylor expansion gives
\begin{equation}\label{an:annular-shortness}
R\le\|F\|\lesssim \left(K_{\mathrm{lin}}d_1^{-1}+\sum_{r=2}^kd_1^{-r}\right)\lesssim K_{\mathrm{lin}}d_1^{-k}.
\end{equation}
Thus $K_{\mathrm{lin}}=R^\delta$, with $0<\delta<1/2$, forces $d_1\lesssim R^{-(1-\delta)/k}$. Use half of $d_1^{\mu_1/p}$ for this cutoff and sum the remaining positive dyadic powers by Minkowski and \eqref{an:physical-positivity}:
\[
\left\|\sum_D E_D\right\|_{L^p(\mathcal A_R)}\lesssim_{p,a} R^{\delta A_p/p-(1-\delta)\mu_1/(2kp)}L_R^{B_p/p}.
\]
After fixing the incidence constants, choose $\delta$ so that $\delta A_p/p<(1-\delta)\mu_1/(4kp)$. Absorbing the logarithmic factor gives $\|\sum_D E_D\|_{L^p(\mathcal A_R)}\lesssim_{p,a} R^{-\sigma_p}$ for some $\sigma_p>0$.

The annulus has measure $O(R^{N_k})$. Lemma~\ref{an:linear-tail} and $\rho\sim L_R^{-1}$ give
\[
\left\|E_a-\sum_D E_D\right\|_{L^p(\mathcal A_R)}\lesssim_{m,a} R^{N_k/p-\delta m}L_R^m.
\]
Choose $m$ after $\delta$ so that $N_k/p-\delta m<-2\sigma_p$; the remainder is then $O(R^{-\sigma_p})$. Summing the $p$th powers over dyadic annuli and using $|E_a|\le\|a\|_1$ on a bounded coefficient ball proves the claim. The amplitude bounds use only its mass, supremum, and the finitely many derivatives required by this fixed $m$.
\end{proof}

\subsection{Reduction from smooth amplitudes to a cube}\label{completion:cube}\label{an:cube-subsection}

\begin{lemma}[Reduction from smooth amplitudes to a cube]\label{an:cube-reduction}
Fix $d,k\ge1$ and $1\le p<\infty$, with $\chi_{d,k}$ as in Lemma~\ref{an:character}. Suppose that every compactly supported smooth amplitude $a$ satisfies
\[
F\longmapsto E_a(F)=\int_{\mathbb R^d}a(x)e^{iF(x)}\,dx\in L^p(V_{d,k}).
\]
If $p>\chi_{d,k}$, the same conclusion holds for the amplitude $\mathbf1_{[0,1]^d}$.
\end{lemma}
\begin{proof}
Take a nonnegative dyadic smooth partition of $(0,1)$ consisting of one fixed interior bump and, at each endpoint, a rescaled copy of a fixed bump at every scale $2^{-n}$, $n\ge2$. Tensor products give, almost everywhere, a partition of the cube into $a_\alpha\circ T_\alpha^{-1}$, where the $a_\alpha$ belong to a finite list of fixed smooth amplitudes and $T_\alpha(x)=c_\alpha+B_\alpha x$ has diagonal entries of absolute value $1$ or $2^{-n}$.

Physical and coefficient changes of variables, using Lemma~\ref{an:character}, give
\begin{equation}\label{an:atom-scaling}
\|E_{a_\alpha\circ T_\alpha^{-1}}\|_{L^p(V_{d,k})}=|\det B_\alpha|^{1-\chi_{d,k}/p}\|E_{a_\alpha}\|_{L^p(V_{d,k})}.
\end{equation}
With $\epsilon=1-\chi_{d,k}/p>0$, the sum of these norms is finite because
\[
\sum_\alpha|\det B_\alpha|^\epsilon=\left(1+2\sum_{n\ge2}2^{-n\epsilon}\right)^d<\infty.
\]
Minkowski gives convergence in $L^p$. For each fixed phase the oscillatory integrals also sum absolutely, since the nonnegative physical atoms have total integral one. Tonelli identifies their sum with the cube integral.
\end{proof}

\section{Geometric decomposition in three variables}\label{geo:section}

The constant-dependence convention at the start of Section~\ref{an:section} also applies to the geometric and coefficient estimates in Sections~\ref{geo:section}--\ref{surf:section}.

Use the fixed upper-affine presentation of Lemma~\ref{an:rq} and take coefficient integrals in the coordinates of the original family maps. This section classifies the rank-four models and proves the surface assignments of Theorems~\ref{geo:common-allocation} and~\ref{geo:gauss}. Their patches, label counts and multiplicities supply the geometric input to Section~\ref{surf:section}; Section~\ref{reg:section} treats regular pieces.

\subsection{The augmented orbit and exact models}\label{geo:augmented-models}\label{geo:orbit-subsection}

Use the group $\mathcal U$ and its coordinates $(a,b,c,d,e,f)$ from~\eqref{geo:group}.
\begin{definition}[Augmented states]\label{def:augmented-state}
Its infinitesimal generators are $A=y\partial_x$, $C=z\partial_x$, $T_x=\partial_x$, $D=z\partial_y$, $T_y=\partial_y$, and $T_z=\partial_z$. An augmented polynomial is a pair $\widehat F=([F],\ell)$, where $[F]\in V_k$ and $\ell(u)=p+Mu$ is the physical affine map. Every coefficient of $\ell$, including its constant term, is retained. The augmented norm is the product Euclidean coefficient norm, $\|\widehat F\|^2=\|[F]\|^2+|p|^2+\|M\|_{\mathrm{HS}}^2$. The physical entries are parameters, not additional variables in coefficient integration. Define
\[
R_{\widehat F}X=([XF],X\ell),\qquad X\in\operatorname{span}\{A,C,T_x,D,T_y,T_z\}.
\]
The six-column matrix of this map is the augmented orbit differential. Its singular values are ordered decreasingly.

We retain the coefficient window and augmented translation inequality
\begin{equation}\label{geo:translation-stop}
\|\widehat F\|\le P,\qquad \|[h\cdot\nabla F]\|^2+|Mh|^2\ge P^{-2}|h|^2\quad(h\in\mathbb R^3).
\end{equation}

In a phase-only translation condition the physical term $|Mh|^2$ is omitted. The matrix $M$ may be singular in closures of parameter sets, although the physical frames themselves are invertible.
\end{definition}
\begin{definition*}[Pivot conditioning]
A scalar pivot $q$ is conditioned if it satisfies $|q|\ge P^{-C}$ and its stated upper bound. A matrix pivot $L$ is conditioned if, on the specified columns and coefficient space, it satisfies
\[
s_{\min}(L)\ge P^{-C},\qquad \|L\|_{\mathrm{op}}\le P^C.
\]
The exponent $C$ is fixed independently of $P$.
\end{definition*}

In the conditioning estimates below, the constants denoted by $C_P$ and positive $c_P$ satisfy
\[
C_P\le CP^C,\qquad c_P^{-1}\le CP^C.
\]
The same power bound applies to the implicit constants in $O_P$, $\lesssim_P$, $\gtrsim_P$, and $\sim_P$. The constants and exponents in these bounds depend only on the fixed degree and the stated auxiliary data.

Here and below, an algebraic or semialgebraic family has fixed format when its number of variables, defining polynomials, and polynomial degrees are bounded in terms of the fixed degree and the stated auxiliary data. The semialgebraic facts used below are cell decomposition with the external parameters included, curve selection, monotonicity, quantifier elimination, and one-variable semialgebraic growth in fixed format \cite{BPR}. Replacing $P$ by a fixed power accommodates any fixed polynomially conditioned version of these inequalities. The exact rank-four classification uses the closure of the bounded physical parameter set, allowing singular $M$ and retaining the translation coordinate and the augmented translation inequality.

\begin{lemma}[Quantitative semialgebraic conditioning]\label{geo:conditioning}
Suppose a nonempty family $K_P$ of compact parameter sets is defined by a fixed first-order formula over the real numbers, with parameter $P\ge2$, and its coordinates have polynomial bounds in $P$. Assume that $\{(P,x):2\le P\le P_0,\ x\in K_P\}$ is closed for every $P_0<\infty$.

\noindent\textup{(i)} For jointly continuous semialgebraic functions $f,g$ on this family with $f>0$, there are $c,C>0$ such that
\[
\min_{x\in K_P}f(P,x)\ge cP^{-C},\qquad \max_{x\in K_P}|g(P,x)|\le CP^C.
\]

\noindent\textup{(ii)} Fixed-format semialgebraic families have uniformly bounded numbers of connected components, and their zero-dimensional fibers have uniformly bounded cardinalities.

\noindent\textup{(iii)} For a semialgebraic set $A\subset Z\times X$, there is a Borel measurable function $\sigma:\pi_Z(A)\to X$ such that $(z,\sigma(z))\in A$ for every $z\in\pi_Z(A)$.
\end{lemma}
\begin{proof}
Quantifier elimination makes the minimum semialgebraic, and one-variable Puiseux growth gives its power bound at infinity; see \cite[Theorems 2.77 and 3.14]{BPR}. On every bounded $P$ interval the total family is closed and bounded, hence compact, so joint continuity and positivity give a positive minimum there as well. Maxima are treated in the same way. Cell decomposition and triviality applied with the external parameters included give the component, finite-fiber, and selection assertions, with uniform format by quantitative quantifier elimination; see \cite[Theorems 5.6, 5.46, and 14.16]{BPR}. Strict windows are placed inside the stated larger closed windows before taking minima.
\end{proof}

We use the following continuity estimate on the normalized windows under consideration. If $|D(z_0)|\ge P^{-a}$ and $|D(z)-D(z_0)|\le P^b\|z-z_0\|$, then
\[
\|z-z_0\|\le\tfrac12P^{-(a+b)}\quad\Longrightarrow\quad |D(z)|\ge\tfrac12P^{-a}.
\]
The Lipschitz bound here is required on the stated window; it follows from the derivative bounds for the particular determinant or implicit pivot used there.

\begin{lemma}[Paths and variation]\label{geo:paths}

\leavevmode\noindent\textup{(i)} In a connected semialgebraic set of fixed format, two points can be joined by a continuous piecewise analytic semialgebraic path of uniformly bounded format.

\noindent\textup{(ii)} If finitely many continuous semialgebraic functions of fixed format have ranges of lengths at most $B_i$ on the set, the path can be subdivided into a uniformly bounded number of intervals on which all those functions are monotone. Their total variations along the path are therefore $\lesssim B_i$.
\end{lemma}
\begin{proof}
Semialgebraic connectedness implies semialgebraic path connectedness \cite[Theorems 5.22--5.23]{BPR}. Applying the cell-decomposition and curve-selection constructions with both endpoints as parameters gives paths of uniformly bounded format; see \cite[Theorems 3.19, 5.6, and 14.16]{BPR}. Subdivision into analytic pieces and by the signs of the derivatives of the composed functions gives uniformly finitely many monotone intervals. Each contributes variation at most $B_i$.
\end{proof}

\begin{lemma}[Algebraic coefficient maps]\label{geo:coefficient-map}\label{low:4:sf:coefficient-map}\label{low:5:qsf:coefficient-map}\label{low:6:qsf:coefficient-map}
Let $\mathcal T$ be a fixed-format algebraic map between open sets of the same finite dimension. Let $\Omega$ be the union of the selected regular branches on which $|\det D\mathcal T|\ge P^{-C}$.

\noindent\textup{(i)} For every measurable output box $Q$, $|\{u\in\Omega:\mathcal T(u)\in Q\}|\le P^{C'}|Q|$.

\noindent\textup{(ii)} With external parameters $\zeta$ fixed, the same estimate holds for $u\mapsto\mathcal T_\zeta(u)$ when $|\det D_u\mathcal T_\zeta|\ge P^{-C}$ and the selected branches have uniformly bounded algebraic format.
\end{lemma}
\begin{proof}
The nonzero Jacobian makes each selected fiber discrete. Uniform semialgebraic finiteness bounds its cardinality \cite[Theorem 5.46]{BPR}; the area formula \cite[Theorem 1]{low:area} then gives the estimate. For the parameter-dependent assertion, apply this argument to the remaining square map for each fixed parameter; uniform format gives the same bound on the number of preimages. The fixed parameters remain in the joint measure and are subsequently integrated by Fubini.
\end{proof}

\begin{lemma}[The normal columns and the fourth singular value]\label{geo:normal-pivot}
Under \eqref{geo:translation-stop},
\[
\begin{aligned}
\textup{(i)}\quad &s_{\min}(R_{\widehat F}|_{\operatorname{span}\{A,C,T_x\}})\ge P^{-C},\\
\textup{(ii)}\quad &\sigma_4(R_{\widehat F})\ge P^{-C},
\end{aligned}
\]
where the singular values of the full map are ordered decreasingly. In particular, $\operatorname{rank}R_{\widehat F}\ge4$.
\end{lemma}
\begin{proof}
The first restriction is the map
\[
L\longmapsto([LF_x],(Me_x)L),\qquad L\in\operatorname{span}\{1,y,z\}.
\]
If $Me_x\ne0$, its physical part is injective. If $Me_x=0$, the translation inequality implies that $F_x$ is nonconstant. For a nonzero affine $L$, the polynomial $LF_x$ cannot be constant. Thus the restriction is injective in either case. Compactness and Lemma~\ref{geo:conditioning} make this quantitative.

We next exclude two simultaneous augmented kernel fields
\[
Y=\partial_y+L(y,z)\partial_x,\qquad Z=z\partial_y+N(y,z)\partial_x,
\]
where $L,N$ are affine. If they existed, $YF=c_1$ and $ZF=c_0$ for constants $c_0,c_1$, and $Y\ell=Z\ell=0$. Subtracting $zY$ from $Z$ gives
\[
(N-zL)F_x=c_0-c_1z,\qquad (N-zL)Me_x=0.
\]
If $F_x$ is nonconstant and the linear part of $L$ is nonzero, the polynomial on the left of the first identity has degree at least three, whereas the right side has degree at most one. Hence $L$ is constant in this case. If $F_x$ is constant, the translation inequality gives $Me_x\ne0$, so the second identity implies $N=zL$ and again forces $L$ to be constant. A constant $L$ makes $Y$ a nonzero constant translation in the augmented kernel, contradicting \eqref{geo:translation-stop}. If the orbit rank were three, both $T_y$ and $D$ would be combinations of the three independent normal columns and would produce precisely this excluded pair. Therefore the rank is at least four. Positivity of the fourth singular value on the closed stopped window, followed by Lemma~\ref{geo:conditioning}, gives its stated lower bound.
\end{proof}

In the model representation \eqref{geo:classified-model}, the one-variable polynomial $f$ is called the profile and the polynomial $q$ is called the invariant. The cases $\gamma=0$ and $\gamma\ne0$ are called the quadratic-invariant and cubic-invariant branches, respectively.
\begin{theorem}[Exact rank-four models]\label{geo:classification}

\leavevmode\noindent\textup{(i)} If \eqref{geo:translation-stop} holds and $R_{\widehat F}$ has rank four, then $M=0$ and, modulo a constant,
\begin{equation}\label{geo:classified-model}
\begin{split}
F&=f(q)+\lambda\left(y-\frac\gamma2z^2\right)+\mu z,\\
q&=x-\frac r2y^2-syz-ty+\frac{t\gamma-v}{2}z^2+\frac{s\gamma}{3}z^3-wz,\\
r\gamma&=0,\qquad \Delta:=rv-s^2\ne0,\qquad \deg f\ge2.
\end{split}
\end{equation}
If $\gamma=0$, then $\deg f\le\lfloor k/2\rfloor$. If $\gamma\ne0$, then $r=0$, $s\ne0$, and $\deg f\le\lfloor k/3\rfloor$.

\noindent\textup{(ii)} Conversely, each displayed polynomial has injective phase translation differential and, with $M=0$, orbit rank four.

\noindent\textup{(iii)} On the coefficient set defined by \eqref{geo:translation-stop}, all displayed parameters and $|\Delta|^{-1}$ are bounded by $P^C$. Writing $f(t)=\sum_{\nu=0}^k a_\nu t^\nu$, the nonlinear coefficient vector satisfies
\[
\left(\sum_{\nu=2}^k|a_\nu|^2\right)^{1/2}\ge P^{-C}.
\]
\end{theorem}
\begin{proof}
Let $\mathfrak k$ be the augmented kernel. It is a Lie subalgebra: fields in the kernel differentiate $F$ to constants and annihilate $\ell$, so their commutator annihilates both. It has dimension two. The upper-affine algebra is nilpotent, and every two-dimensional nilpotent Lie algebra is abelian. Thus $\mathfrak k$ is abelian.

The normal subspace $\mathfrak n=\operatorname{span}\{A,C,T_x\}$ is an ideal. Lemma~\ref{geo:normal-pivot} shows that projection of $\mathfrak k$ to the quotient by $\mathfrak n$ is injective. That quotient is the Heisenberg algebra with basis $D,T_y,T_z$ and center $T_y$. An abelian two-plane in the Heisenberg algebra contains its center: otherwise its image in the two-dimensional quotient by the center has dimension two, and the nondegenerate commutator form would not vanish. The plane cannot be $\operatorname{span}\{D,T_y\}$, by the excluded-pair argument in Lemma~\ref{geo:normal-pivot}. It therefore has a unique normalized lift
\[
Y=(ry+sz+t)\partial_x+\partial_y,\qquad Z=(uy+vz+w)\partial_x+\gamma z\partial_y+\partial_z.
\]
Computing $[Y,Z]$ gives $(u-s-r\gamma z)\partial_x$. Hence $u=s$ and $r\gamma=0$.

For the polynomial $q$ in \eqref{geo:classified-model}, direct differentiation gives $Yq=Zq=0$. The change $(x,y,z)\mapsto(q,y,z)$ is a polynomial automorphism with polynomial inverse. In these coordinates, $Y=\partial_y$ and $Z=\gamma z\partial_y+\partial_z$. Integrating $YF=\lambda$ and $ZF=\mu$ gives the displayed expression for $F$.

If $\gamma\ne0$, then $r=0$. If $s=0$, $Y=t\partial_x+\partial_y$ would be a constant translation in the augmented kernel; therefore $s\ne0$. The identities $Y\ell=Z\ell=0$ first give $Me_x=Me_y=0$ and then $Me_z=0$. The cubic homogeneous part of $q$ is nonzero, so $\deg f(q)=3\deg f$. The remaining tangential summand has degree at most two. The translation inequality excludes a linear profile, and it follows that $3\deg f\le k$. In this case $\Delta=-s^2\ne0$.

Suppose $\gamma=0$. If $rv-s^2=0$, a nonzero linear combination of $Y,Z$ has constant $\partial_x$ coefficient and gives a prohibited constant translation in the augmented kernel. Thus $\Delta\ne0$. The physical identities $Y\ell=Z\ell=0$ and this nonzero determinant imply $M=0$. The quadratic part of $q$ is nonzero; therefore $2\deg f\le k$. A linear profile would have constant $F_x$ and violate the $x$ translation inequality, so the profile is nonlinear.

Conversely the displayed fields $Y,Z$ annihilate $F$ modulo constants. To prove phase translation injectivity, suppose $[h\cdot\nabla F]=0$ and put $L=h\cdot\nabla q$. This polynomial is independent of $x$. In $h\cdot\nabla F$, the coefficient of the highest power of $x$ is the nonzero leading coefficient of $f'$ times $L$, since the tangential summand is independent of $x$. Hence $L=0$. Its linear $y,z$ coefficients give
\[
rh_y+sh_z=0,\qquad sh_y+(v-t\gamma)h_z=0.
\]
The determinant is $r(v-t\gamma)-s^2=\Delta$, because $r\gamma=0$. Therefore $h_y=h_z=0$, and the constant term gives $h_x=0$. The orbit has rank at most four because $Y,Z$ are independent kernel fields, and at least four by the preceding lemma.

For quantitative bounds, the exact rank-four subset of each closed stopped window is compact. Its kernel varies continuously because the fourth singular value is bounded below and the fifth is zero. The projection to the Heisenberg quotient is injective, and its image avoids the excluded plane. The normalized basis used above therefore has bounded coefficients on that compact set. The normalization $q(x,0,0)=x$ recovers the coefficients of $f$ from the pure $x$ coefficients of $F$; the tangential linear coefficients then recover $\lambda,\mu$. The nonzero quantities $|\Delta|$ and the norm of the nonlinear profile coefficients have positive minima. These constructions and extrema are semialgebraic, so Lemma~\ref{geo:conditioning} supplies polynomial bounds.
\end{proof}

\subsection{Normalization and normal orders}\label{geo:normalization}\label{geo:entry-subsection}

For $q_0=x-yz$, the phase $q_0^m$ has $\partial_x^jF=0$ on $q_0=0$ for every $j<m$, and $\partial_x^mF=m!$. Critical and inflection sheets therefore do not exhaust the normal orders. Moreover, $q_0^2+\varepsilon q_0^m$ has the same critical graph and second normal restriction as $q_0^2$ when $m>2$, but different higher normal coefficients. Even a critical branch must retain all restrictions on the original phase.

A pointwise rate band also supplies no interior margin of the admissible set. For $F=(x-yz)^2+\varepsilon xyz^3$, the critical graph is $c=yz-\varepsilon yz^3/2$, and its null rate satisfies $v(y,0)=3\varepsilon y/2$. A narrow prescribed interval of rate values can therefore cut an arbitrarily thin strip from a full admissible neighborhood. The buffer below is proved for the complete unbanded admissible set.

The second obstruction concerns transport. Put
\[
q_\gamma=x-yz+\gamma z^3/3,\qquad U_\gamma=y-\gamma z^2/2,\qquad F=q_\gamma^2-2\alpha q_\gamma U_\gamma^2.
\]
Its critical graph is $c=yz-\gamma z^3/3+\alpha U_\gamma^2$, with $F_{xx}=2$. Writing a null slope as $d=\gamma z+s(U_\gamma)$ gives the equation $\alpha s^2+s-\alpha\gamma U_\gamma=0$. The branch through $s(0)=0$ has $s'(0)=\alpha\gamma$. Along the characteristic $U_\gamma=0$, the transverse flow Jacobian is consequently $\exp(\alpha\gamma t)$. The identities $(y\partial_x+\gamma z\partial_y+\partial_z)q_\gamma=0$ and $(y\partial_x+\gamma z\partial_y+\partial_z)U_\gamma=0$ show that this affine field preserves both polynomials and hence every normalized phase coefficient along its flow. Coefficient bounds do not imply a polynomial bound for this long transport. Below, characteristic flows are used only for a prescribed polynomially bounded time, with the coefficient defect selected afterwards.

\begin{definition}[Normal sheets]\label{def:normal-sheet}
For $m\ge2$, an order-$m$ sheet is a selected regular graph $x=c(y,z)$ on which the $(m-1)$st normal derivative vanishes and the $m$th does not. The conditioned sheets used below satisfy
\begin{equation}\label{geo:normal-sheet}
\partial_x^{m-1}F(c(y,z),y,z)=0,\qquad \kappa_m(y,z):=\partial_x^mF(c(y,z),y,z),\qquad |\kappa_m|\ge P^{-C},
\end{equation}
We call the sheet critical when $m=2$ and inflection when $m=3$. No vanishing of the lower normal derivatives is required.
\end{definition}
\begin{lemma}[Approximation by rank-four models]\label{geo:rank-entry}
Fix $L\ge1$. There is an exponent $C_L$ such that every label satisfying \eqref{geo:translation-stop} satisfies at least one of the following alternatives.

\textup{(i)} Its fifth singular value satisfies
\[
\sigma_5(R_{\widehat F})\ge P^{-C_L}.
\]
\textup{(ii)} It admits a normalized representative obtained by right composition with a map in $\mathcal U$, whose complete coefficient vector is within $P^{-L}$ of an exact model in \eqref{geo:classified-model}. The normalization satisfies the bounds by powers of $P$ associated with \eqref{geo:translation-stop} and places the phase on an order-$m$ sheet satisfying \eqref{geo:normal-sheet}, where
\[
2\le m\le\lfloor k/2\rfloor.
\]
In case \textup{(ii)}, for every fixed finite order, the graph derivatives through that order are bounded by $P^C$, with $C$ allowed to depend on the order. The normal order, root branch, and chart range over finitely many labels, and the normalization uses at most $P^C$ relative unit neighborhoods. The normalizing map for $(F,j)$ is Borel measurable in $F$ on each such choice; the original index and its weight are retained.
\end{lemma}
\begin{proof}
On the compact set defined by \eqref{geo:translation-stop}, the zero set of the fifth singular value is exactly the rank-four family, by Lemma~\ref{geo:normal-pivot}. On the closed set at the prescribed positive distance from that family, the fifth singular value has a positive minimum. Lemma~\ref{geo:conditioning} makes the minimum an inverse power of $P$. Since $M=0$ on the exact family, small fifth singular value also gives complete augmented proximity, including the physical linear part. The exact model has phase translation injectivity, quantitatively on its compact stopped window; this remains true after a sufficiently small perturbation. If the exact family is empty, the same minimum is positive on the whole stopped window.

For a nonlinear real polynomial $f$, some derivative $f^{(m-1)}$ has a simple real zero: take $m$ equal to the actual degree of $f$, for which that derivative is nonconstant affine. A chosen simple real root and its nonzero next derivative persist on a neighborhood in the coefficient space, including perturbations which increase the degree. The conditioned family of profiles in Theorem~\ref{geo:classification} is compact for each $P$. It is therefore covered by finitely many such persistent root neighborhoods. Equivalently, there is $A(P)<\infty$ such that each profile has some $m,t$ satisfying
\[
f^{(m-1)}(t)=0,\qquad |t|\le A(P),\qquad |f^{(m)}(t)|\ge A(P)^{-1}.
\]
This is a fixed-format quantified semialgebraic condition. Definable choice and semialgebraic growth allow $A(P)\le P^C$.

Translate $x$ to the selected root. The translation is polynomially bounded, so the original coefficient defect remains as small as needed if its exponent was chosen after this bound. The implicit function theorem applied to $\partial_x^{m-1}F=0$, with derivative $\kappa_m$, gives the graph on a neighborhood of radius at least $P^{-C}$ and bounds its derivatives through each prescribed finite order by $P^C$. Normal translation and two normal shears set its position and tangent plane to zero at the reference point. Those transformations are again polynomially bounded. Choose the first permissible $m$ and then the first permissible real root in a fixed ordered branch list. These choices are semialgebraic and hence Borel. A polynomially bounded region of the fixed-dimensional group is covered by $P^C$ relative unit neighborhoods, so the original family packing changes by at most a power of $P$. The resulting normalized frames are the comparison representatives of the original labels.
\end{proof}
\begin{definition}[Tangential charts]\label{def:tangential-charts}
A Gauss, or nonnull, chart has first tangential curvature $r=c_{yy}$ satisfying $P^{-C}\le|r|\le P^C$ and uses $p=c_y$ as a coordinate. A null chart has a selected simple root $d$ with $|d|\le P^C$. Put $B=c_{yz}+dc_{yy}$ and require
\begin{equation}\label{geo:null-equation}
c_{zz}+2dc_{yz}+d^2c_{yy}=0,\qquad |B|\ge P^{-C}.
\end{equation}
For the null equation $Q(d)=c_{zz}+2dc_{yz}+d^2c_{yy}=0$, the root derivative is $|Q^\prime(d)|=2|B|\ge2P^{-C}$. The Gauss chart permits either sign of the Hessian determinant. These tangential charts are independent of the normal order in Definition~\ref{def:normal-sheet}.
\end{definition}
\begin{lemma}[A two-chart shear normalization]\label{geo:two-chart-shear}
Write $H=(h_{ij})$ and $T_d=\left(\begin{smallmatrix}1&d\\0&1\end{smallmatrix}\right)$. On any bounded set of nonsingular symmetric $2$ by $2$ matrices with bounded inverses, a bounded shear can achieve one of the following conditions, with uniform positive lower bounds on the displayed pivots:
\begin{align*}
\mathrm{I}:&\quad (T_d^THT_d)_{12}=0,\qquad |h_{11}|\ge c_P;\\
\mathrm{II}:&\quad (T_d^THT_d)_{22}=0,\qquad |(T_d^THT_d)_{12}|\ge c_P.
\end{align*}
The second chart is used only for indefinite matrices.
\end{lemma}
\begin{proof}
Choose a sufficiently small fixed threshold for $|h_{11}|$. Above it take $d=-h_{12}/h_{11}$. Below it, boundedness of $h_{22}$ and the lower bound for $|\det H|$ imply $h_{12}^2-h_{11}h_{22}\ge c_P>0$ and $|h_{12}|\ge c_P$. The stable root
\[
d=-\frac{h_{22}}{h_{12}+\operatorname{sgn}(h_{12})\sqrt{h_{12}^2-h_{11}h_{22}}}
\]
is bounded, solves $h_{22}+2dh_{12}+d^2h_{11}=0$, and has $|h_{12}+dh_{11}|=\sqrt{h_{12}^2-h_{11}h_{22}}$. This also covers $h_{11}=0$.
\end{proof}

\begin{lemma}[Gauss and real-null alternatives]\label{geo:chart-split}

\leavevmode\noindent\textup{(i)} The sheets in Lemma~\ref{geo:rank-entry} admit a finite exhaustive division into the charts of Definition~\ref{def:tangential-charts}.

\noindent\textup{(ii)} In the null chart, a bounded tangential shear subtracts the selected root, giving $d=0$ and quadratic graph part $ry^2/2+Byz$ at the reference point. These choices require no inverse power of the cubic-invariant parameter $\gamma$.
\end{lemma}
\begin{proof}
The exact invariant graph has Hessian determinant $\Delta=rv-s^2$, also at points on a cubic invariant because $r\gamma=0$. The determinant remains conditioned on a sufficiently small actual neighborhood. Fix a small inverse-polynomial threshold for $|r|$. Above the threshold use $p=c_y$ as a coordinate. In that chart a nearby exact model must have $r\ne0$ if the complete defect is chosen after the threshold; hence $\gamma=0$ and its invariant graph is quadratic.

Below the threshold, apply chart II of Lemma~\ref{geo:two-chart-shear} to the bounded graph Hessian with conditioned inverse. Its stable root is bounded and has $|c_{yz}+dc_{yy}|\ge P^{-C}$; it also covers $r=0$. The tangential shear replacing $y$ by $y+dz$ makes that root zero and gives the displayed quadratic form. The construction depends only on the normal pivot and the two tangential curvature pivots, not on $\gamma$ or a lower normal scalar.
\end{proof}

\begin{lemma}[Pivots at the fixed frame]\label{geo:original-subgroup-pivots}
At every original label in alternative \textup{(ii)} of Lemma~\ref{geo:rank-entry}, write $R_F^{\mathrm{ph}}X=[XF]$ for the phase-only orbit differential and set
\[
\mathfrak h_0=\operatorname{span}\{A,C,T_x,D\}=\operatorname{Lie}\{e=f=0\}.
\]

\noindent\textup{(i)} At the unchanged original frame,
\[
s_{\min}(R_F^{\mathrm{ph}}|_{\mathfrak h_0})\ge P^{-C},\qquad s_{\min}(R_F^{\mathrm{ph}}|_{\operatorname{span}\{A,C,T_x\}})\ge P^{-C}.
\]

\noindent\textup{(ii)} On a Gauss branch, also $s_{\min}(R_F^{\mathrm{ph}}|_{\mathfrak h_{\mathrm G}})\ge P^{-C}$, where $\mathfrak h_{\mathrm G}=\operatorname{span}\{C,T_x,D,T_y\}$. These bounds hold before any phase-dependent endpoint normalization.
\end{lemma}
\begin{proof}
At an exact rank-four model, $M=0$, so the augmented kernel is the phase-only kernel. Its basis $Y,Z$ in Theorem~\ref{geo:classification} has respective $(T_y,T_z)$ components $(1,0)$ and $(0,1)$. Its intersection with $\mathfrak h_0$ is therefore zero. Compactness and semialgebraic conditioning make the phase-only restriction quantitatively injective. Choose the complete phase defect after this bound; then phase-only injectivity persists at the fixed frame. Its three-column normal restriction inherits the same bound. This proves the original $\mathfrak h_0$ assertion.

In the Gauss chart, the original model parameter $r$ is conditioned. A kernel combination $aY+bZ$ belongs to $\mathfrak h_{\mathrm G}$ only if its $T_z$ coefficient $b$ vanishes and its $A$ coefficient $ar$ vanishes; therefore $a=b=0$. Normal translation and normal shears do not change this model curvature, and tangential shearing preserves $c_{yy}$. Selecting the original coefficient defect after the Gauss curvature threshold gives the same quantitative conclusion at the fixed frame. This subgroup is the kernel of the homomorphism $g\mapsto(a,f)$.
\end{proof}

\subsection{Admissible sets and model neighborhoods}\label{geo:activity}\label{geo:activity-subsection}

For each original coefficient vector $F$ and fixed RQ bin, put $F_\ell=F\circ D_\ell$ and retain the fixed physical map $D_\ell$. All graph equations, admissible sets, and common cells below are constructed on this single base orbit. The fixed original labels are $g_j$, and the normalized frames $g_j\circ r_j$, with $r_j,r_j^{-1}$ bounded as above, lie on the finitely many selected root-sheet frame sections of that orbit. Indeed, for every $g\in\mathcal U$,
\[
\partial_x^{m-1}(F_\ell\circ g)=(\partial_x^{m-1}F_\ell)\circ g,
\]
since the derivative of $g$ sends $e_x$ to $e_x$. Thus graph normalization places the labels on the selected sheets of the same phase. Finite choices of normal order, root, and tangential chart are resolved for all labels by the same fixed rules. The original coefficient substitution remains the one through the fixed $g_j$, irrespective of the normalized representative $g_j\circ r_j$.

\begin{definition}[Admissibility]\label{def:admissibility}
\leavevmode\noindent\textup{(i)} For each surface branch, fix the normal order and root in Definition~\ref{def:normal-sheet}, a tangential chart in Definition~\ref{def:tangential-charts}, and the corresponding compact model family of Theorem~\ref{geo:classification} or its stated degree specialization.

\noindent\textup{(ii)} The admissible set retains the complete coefficient window and physical shape comparisons of Section~\ref{an:section}, expressed in the RQ coordinates of Lemma~\ref{an:rq}, the augmented translation inequality~\eqref{geo:translation-stop}, the selected normal and tangential equations with their pivot bounds, and the stated coefficient-distance bound from that model family. In a branch defined using $\partial_x^{m-2}F$, retain every coefficient and physical restriction of $F$. The model defect means the distance between the complete coefficient vectors in the specified normalized frame.

\noindent\textup{(iii)} Fix nested inner and outer admissible sets before coefficient integration. For normalized upper-bound and lower-bound tests, respectively, use
\[
\begin{aligned}
q\le b_{\rm in}\ \text{(inner)},\quad q\le b_{\rm out}\ \text{(outer)},\quad b_{\rm out}-b_{\rm in}&\ge P^{-C},\\
q\ge a_{\rm in}\ \text{(inner)},\quad q\ge a_{\rm out}\ \text{(outer)},\quad a_{\rm in}-a_{\rm out}&\ge P^{-C}.
\end{aligned}
\]
These thresholds are fixed for the selected branch before coefficient integration.

\noindent\textup{(iv)} The normal equation is retained as an equality, separately from these inequality margins. A normalized label is admissible when it satisfies this branch's inner restrictions. A selected regular branch retains its root label and nonzero implicit pivot. Continue this branch across an artificial cell cut whenever the outer restrictions and the pivot permit; such a cut is not a boundary of the admissible domain.
\end{definition}

\begin{definition}[Null coordinates and buffers]\label{def:null-coordinates}
On a null sheet use
\[
Y=\partial_y,\qquad Z=\partial_z+d\partial_y,\qquad \theta=dy-d\,dz,\qquad u=Yd,\qquad v=Zd.
\]
Also write $r=c_{yy}$ and $B=c_{yz}+dc_{yy}$, so $B=c_{yz}$ at a null-normalized point. For an absolutely continuous path $\eta(t)=(y(t),z(t))$ in the selected branch, set
\[
\operatorname{length}_{\theta,dz}(\eta)=\int\bigl(|\dot y(t)-d(\eta(t))\dot z(t)|+|\dot z(t)|\bigr)\,dt.
\]
This is the length used in this section unless another coordinate norm is specified. A horizontal segment has its ordinary length, and a characteristic of $Z$ has length equal to its $z$ increment. A buffer of size $M$ means $\operatorname{length}_{\theta,dz}(\eta)\ge M$ for every such path from the specified set to the complement of the outer admissible set, measured up to its first exit.
\end{definition}

\begin{lemma}[Boxes in model coordinates]\label{geo:model-box}
Fix a polynomial radius $M=P^A$ and a finite derivative order $N$.

\noindent\textup{(i)} There is an inverse-polynomial upper bound for the complete coefficient defect $\delta$ such that an actual order-$m$ null sheet near a conditioned exact model has its selected normal and null roots on the entire model-parameter box of radius $20M$. Write $x_F(s,t),d_F(s,t)$ for those roots and $x_0(s,t),d_0(s,t)$ for the model roots in the same coordinates. They are holomorphic on the corresponding complex model-parameter polydisc and satisfy
\[
\max_{|\alpha|\le N}\left\{|\partial^\alpha(x_F-x_0)|,|\partial^\alpha(d_F-d_0)|\right\}\le P^C(1+M)^C\delta.
\]

\noindent\textup{(ii)} On the real box, the matrix expressing $(\theta,dz)$ in $(ds,dt)$ and its inverse have norms at most $P^C$.

\noindent\textup{(iii)} The outer coefficient and physical inequalities for $F$ in Definition~\ref{def:admissibility} hold throughout the real box. These bounds are uniform when $\gamma$ tends to zero or when $v$ vanishes.
\end{lemma}
\begin{proof}
After centering the selected normal root and applying the normal and null tangent normalizations, the exact model has the form
\[
q=x-\frac r2y^2-Byz+\frac{B\gamma}{3}z^3,\qquad r\gamma=0,\qquad F_0=f(q)+\lambda(y-\gamma z^2/2)+\mu z,
\]
where $f^{(m-1)}(0)=0$, $|f^{(m)}(0)|\ge P^{-C}$, and $|B|\ge P^{-C}$. The commuting affine fields
\[
Y_0=(ry+Bz)\partial_x+\partial_y,\qquad Z_0=By\partial_x+\gamma z\partial_y+\partial_z
\]
preserve $q$ and differentiate the tangential summand to constants. Their joint affine flow $h_0(s,t)$, as well as its inverse, is polynomial of degree at most three in $(s,t)$. This follows by exponentiating their strictly upper-triangular $4\times4$ homogeneous-coordinate matrices. Thus $F_0\circ h_0(s,t)=F_0$ modulo constants for all $(s,t)$.

On the complex parameter polydisc of radius $20M$, polynomial pullback bounds $(F-F_0)\circ h_0(s,t)$ and its prescribed derivatives by $P^C(1+M)^C\delta$, where $\delta$ is the coefficient distance $\|[F]-[F_0]\|$ in the normalized frame. The normal and slope equations have the conditioned model derivatives $f^{(m)}(0)$ and $2B$. The quantitative holomorphic implicit function theorem, on fixed inverse-polynomial root discs, therefore gives the unique roots and their correction bounds when $\delta$ is chosen after $M$; Cauchy estimates on nested polydiscs give the prescribed derivative bounds.

On the exact critical orbit, $y=s+\gamma t^2/2$ and $z=t$. The coframe of the perturbed phase in these parameters is
\[
\theta=ds+(\gamma t-d_F)\,dt,\qquad dz=dt.
\]
The preceding estimates bound its error and its derivatives by $P^C(1+M)^C\delta$. Choose $\delta$ so this error is smaller than the prescribed margin. Normal root corrections and tangent shears give relative affine maps with polynomial bounds and polynomial inverse bounds. All nonconstant coefficients of the transformed phase $F$ remain in the outer set of Definition~\ref{def:admissibility}. The physical center and matrix are transformed by $\ell\circ h$, so their norms and their prescribed dyadic shape comparisons change by at most the corresponding polynomial factors. The outer physical windows were enlarged by these factors before the defect was selected. The only inversions are the conditioned normal and mixed-curvature pivots.
\end{proof}

\begin{lemma}[An interior subset of the admissible set]\label{geo:core}
Fix a local width $m_0=P^a\ge1$ and a larger radius $M$ measured by the coframe in Definition~\ref{def:null-coordinates} before choosing the inner model defect.

\noindent\textup{(i)} On each real-null branch, the normalized original labels from alternative \textup{(ii)} of Lemma~\ref{geo:rank-entry} belong to a fixed-format semialgebraic core $A$ in an outer admissible set $E$ retaining all conditions on $F$ in Definition~\ref{def:admissibility}. Every path from $\overline A$ to $E^c$ has coframe length at least $M$.

\noindent\textup{(ii)} On the local boxes in $E$ needed below,
\begin{equation}\label{geo:small-rates}
|u|+|Yv|+|Zv|\le P^C\delta,\qquad |v|\le P^C,
\end{equation}
where $\delta$ can be selected as a sufficiently small inverse power after $m_0,M$ are fixed.

\noindent\textup{(iii)} Original labels and their weights are retained, and local relative-frame packing changes by only a power of $P$. No rate band is imposed.
\end{lemma}
\begin{proof}
At every exact null model, the selected slope is $d=\gamma z$, so $u=0$, $v=\gamma$, and $Yv=Zv=0$. Choose first an outer complete coefficient defect small enough for \eqref{geo:small-rates} on the prescribed local operations. Choose a model-box radius larger by a polynomial factor than $M$. Lemma~\ref{geo:model-box} supplies an inner defect which keeps the actual complete phase, root pivots, physical data, and model defect inside the strict outer restrictions throughout that larger box.

Define membership in $A$ by the following quantified test: there exists a conditioned reference model and selected root branches for which the full outer restrictions, with a fixed positive part of the gap between the inner and outer bounds in Definition~\ref{def:admissibility} left unused, hold throughout the specified polynomial model box. The reference-model conditions, affine flow, root equations, and full admissibility inequalities are semialgebraic of fixed format. Quantifier elimination \cite[Theorems 2.77 and 14.16]{BPR} therefore makes this a fixed-format semialgebraic condition. The coframe comparison on the larger box proves the buffer. Any path of length less than $M$ starting in the smaller model box remains in the larger box, since in model coordinates each displacement is bounded by a polynomial constant times that length and the box was chosen with the corresponding extra margin. The larger box lies in $E$. The same argument applies to the closure because this unused margin is strict.

Every original inner label passes the test by Lemma~\ref{geo:model-box}. If a model or root is selected, definable choice with fixed ties gives a Borel selection. The model certifies local admissibility, with coefficient integration in the original coordinates. All initial normalizations lie in a polynomially bounded region of $\mathcal U$, which is covered by $P^C$ relative unit neighborhoods. Thus \eqref{an:local-packing} gives at most $P^C\Lambda$ original labels in each such neighborhood. For each phase, the rule is chosen for all labels at once and is kept when a label is fixed for coefficient integration.
\end{proof}

Use the branch convention in Definition~\ref{def:admissibility}. The nonzero implicit derivative, bounded below by $P^{-C}$ on the retained branch, gives unique continuation through artificial cell cuts, and the polynomial degree bounds the number of retained normal and null branches. The local boxes and paths therefore retain the buffer of Lemma~\ref{geo:core}.

\begin{definition}[Prepared lists]\label{def:preparation}
Preparation refers to a specified finite list of scalar functions, derivative orders, coordinates, scales and amplitude bounds. A good point satisfies all the stated derivative inequalities simultaneously. Absolute preparation uses the stated amplitude bounds; relative preparation uses the local nonzero amplitude as in Lemma~\ref{geo:relative-preparation}; characteristic preparation uses derivatives along the selected trajectory. Each application specifies the retained measure and any interval or cell-width requirement. In the later coefficient estimates, this preparation is applied jointly to the finite lists used for scalar widths, curvature widths, seed rounding, and reciprocal collar Jacobians.
\end{definition}
\begin{lemma}[Derivative bounds outside a small exceptional set]\label{geo:preparation}
Let $\mathcal F$ be a finite list of smooth scalar functions on a rectangle $R$ of side lengths $L_1,L_2$, with bounds $|f|\le B_f$ and fixed-format algebraic derivative expressions on their selected regular branches.

\noindent\textup{(i)} For every fixed derivative order $N$ and $\varepsilon>0$, there is a common set $G\subset R$ such that
\[
|R\setminus G|\le\varepsilon L_1L_2,\qquad |\partial^\alpha f(x)|\le C B_f L_1^{-\alpha_1}L_2^{-\alpha_2}\quad(x\in G,\ f\in\mathcal F,\ |\alpha|\le N),
\]
where $\alpha\in\mathbb{Z}_{\ge0}^2$ and $C$ depends only on the format, order, and $\varepsilon$.

\noindent\textup{(ii)} The statement also holds in a monotone flux parameter when every differentiated expression is algebraic in the original seed parameter.

\noindent\textup{(iii)} If $|\partial_\ell f|\le M_\ell$ for $\ell=1,2$, the same common set satisfies
\[
|\partial^\alpha\partial_\ell f|\lesssim M_\ell L_1^{-\alpha_1}L_2^{-\alpha_2},\qquad \ell=1,2,
\]
for each prescribed multi-index $\alpha=(\alpha_1,\alpha_2)$, with the same dependence of the constants.
\end{lemma}
\begin{proof}
In one variable, suppose $|f|\le B$ on an interval of length $L$. Each signed superlevel set $\{\pm f^{(j)}>CBL^{-j}\}$ has uniformly finitely many interval components. On one of length $\ell$, the divided-difference formula \cite[equations (44) and (52)]{deBoor}, applied to equally spaced nodes, gives $CBL^{-j}\ell^j\lesssim_j B$. Hence $\ell\lesssim_j C^{-1/j}L$. Choosing $C$ large makes the total exceptional length arbitrarily small, simultaneously for the finite list of functions and orders.

For two variables, first apply this argument on the horizontal sections and use Fubini. The simultaneous good set is semialgebraic of bounded format. On each vertical section it has a bounded number of interval components. Discard components shorter than a sufficiently small fixed fraction of the vertical side length; their total length is correspondingly small. On each remaining interval, every previously prepared horizontal derivative is bounded, so the one-dimensional argument applied to that derivative bounds its vertical derivatives with the desired scales. Sum the finitely many discarded proportions and repeat through the finite derivative list. This gives the claimed area bound without requiring the final good set to be a rectangle. For the last assertion, include $\partial_1f$ and $\partial_2f$ in the finite scalar list with their respective amplitude bounds $M_1,M_2$.

If $s$ is a monotone flux parameter, the change from the original algebraic parameter $\zeta$ has nonvanishing algebraic derivative. Hence each fixed $s$ derivative is an algebraic expression in $\zeta$, with conditioned denominator on the selected branch. A monotone change of parameter preserves the number of interval components. The same finite-difference proof is performed in the actual $s$ interval, so algebraicity of the primitive defining $s$ is unnecessary.
\end{proof}

\subsection{Null certificates and common assignment}\label{geo:common-assignment}\label{geo:null-subsection}

The local projection cutoff is denoted by $\tau_0$, and the normal order remains $m$. Fix the conditioning powers in the local comparison and moderate-corridor estimates, and let $C_{\mathrm{loc}}$ dominate the finitely many exponents required for their lower width bounds. Choose a fixed $b>0$ for the small local comparisons and then choose $a>b+C_{\mathrm{loc}}+1$. Set $m_0=P^a$ and $\tau_0=P^{-b}m_0$. In particular $P^{-C_{\mathrm{loc}}}\tau_0\ge1$. These local scales are fixed before the model-box radius and the coefficient defect in Lemma~\ref{geo:core}; choose its intrinsic buffer radius $M>\tau_0+6m_0+1$, increasing the fixed polynomial margin for the local boxes if necessary. Thus every lower-width requirement $P^{-C}\tau_0\ge1$ below refers to one of these fixed comparison exponents.

\begin{lemma}[Bounded-time characteristic comparison]\label{geo:local-flow}

\leavevmode\noindent\textup{(i)} A characteristic of $Z$ lying in a prescribed outer model box for $|t|\le2m_0$ has transverse Jacobian and inverse bounded by powers of $P$, independently of any eventual long scale.

\noindent\textup{(ii)} A point with buffer greater than $4m_0$ has such a characteristic for this interval, and its remaining buffer decreases by at most $2m_0$.
\end{lemma}
\begin{proof}
For the characteristic $(y(t;y_0),z_0+t)$, its transverse Jacobian satisfies $J'=uJ$, $J(0)=1$. Thus $J=\exp(\int_0^tu)$, and choosing $2P^C\delta m_0\le\log2$ in \eqref{geo:small-rates} gives $1/2\le J\le2$. Its intrinsic length is $|t|$; the buffer, bounded $v$ and conditioned root pivots give continuation for $|t|\le2m_0$ and remaining buffer at least the original buffer minus $2m_0$. Relative affine comparison follows from the local model box.
\end{proof}

\begin{definition}[Seeds and collars]\label{def:seeds-collars}
A seed is a parametrized arc $\gamma=(a,b)$. Its signed flux is $\theta(\gamma')$; on an arc of nonzero flux, orient and parameterize it by $s$ so that $\theta(\gamma_s)=1$. For a graph $y=\Gamma(z)$, its residual and flux length are $\alpha=\Gamma'-d(\Gamma,z)$ and $\int|\alpha|\,dz$. Write $\gamma(s)=(a(s),b(s))$, $D(s)=d(a(s),b(s))$, and $V(s)=v(a(s),b(s))$. In the collar formulas below, put $\nu=b'(s)=z_s$. The unit-flux identity is $a^\prime(s)-D(s)b^\prime(s)=1$. The parabolic collar is the image of
\begin{equation}\label{geo:collar-map}
\Psi(s,t)=\left(a(s)+D(s)t+\frac12V(s)t^2, b(s)+t\right),\qquad |t|<c_Pm_0,
\end{equation}
A characteristic tube instead uses the actual $Z$-flow from the seed.
\end{definition}
Seeds in $A$ satisfy the buffer and rate hypotheses of the following lemma.

\begin{lemma}[Parabolic collars]\label{geo:collar}
Let $\gamma$ be a fixed-format algebraic seed arc in $E$, parameterized by flux with $|z_s|\le1$. Suppose it has buffer greater than $4m_0$ and \eqref{geo:small-rates} holds on the outer model boxes retaining the conditions on $F$ in Definition~\ref{def:admissibility} for its local characteristic and horizontal comparisons.

\noindent\textup{(i)} For sufficiently small inverse-polynomial $c_P>0$, the collar \eqref{geo:collar-map} lies in the full outer admissible set.  The matrix expressing $\Psi^*(\theta,dz)$ in $(ds,dt)$ and its inverse have norms at most $P^C$.

\noindent\textup{(ii)} Each image point has a uniformly bounded number of preimages.

Writing $\ell_\gamma$ for the seed's flux length, the collar area satisfies
\[
P^{-C}m_0\ell_\gamma\le\operatorname{area}(\operatorname{im}\Psi)\le P^Cm_0\ell_\gamma.
\]

\noindent\textup{(iii)} In the original algebraic seed parameter, the collar map is semialgebraic of fixed format.
\end{lemma}
\begin{proof}
Write $\nu=b'$. The seed identities are
\[
a'=1+D\nu,\qquad D'=u+\nu V,\qquad V'=Yv+\nu Zv.
\]
Direct differentiation gives the coordinate Jacobian
\begin{equation}\label{geo:collar-jacobian}
\det D\Psi=a'+D't+\tfrac12V't^2-\nu(D+Vt)=1+ut+\tfrac12V't^2.
\end{equation}
By \eqref{geo:small-rates}, it lies between $1/2$ and $3/2$ after the defect is chosen sufficiently small.

Integrating $d'=v$ and $v'=Zv$ gives the actual characteristic slope $D+Vt+O_P(\delta t^2)$ and position $a+Dt+Vt^2/2+O_P(\delta|t|^3)$. Lemma~\ref{geo:local-flow} retains this trajectory. The polynomial collar differs by a horizontal displacement $O_P(\delta m_0^3)$; choose $c_P\le1$ and then $\delta$ so this is less than one. The total intrinsic displacement is at most $m_0+1<4m_0$, so the collar and correction segment remain in the certified outer boxes with full parent admissibility.

Comparison with the actual characteristic and the bound for $Yd$ give
\[
d(\Psi(s,t))=D+Vt+O_P(\delta t^2+\delta^2|t|^3).
\]
Consequently the $dt$ coefficient $C_t=D+Vt-d(\Psi(s,t))$ in $\Psi^*\theta$ is small, and
\[
\begin{pmatrix}\Psi^*\theta\\\Psi^*dz\end{pmatrix}=\begin{pmatrix}\det D\Psi+\nu C_t&C_t\\\nu&1\end{pmatrix}\begin{pmatrix}ds\\dt\end{pmatrix}.
\]
These identities bound the coframe matrix and its inverse by powers of $P$, uniformly in the total seed length.

In the original algebraic parameter, all seed coordinates and $D,V$ are algebraic of fixed format. The map has nonvanishing Jacobian, so the finite-fiber and area-formula argument in Lemma~\ref{geo:coefficient-map}, together with \eqref{geo:collar-jacobian}, gives both area comparisons.
\end{proof}

\begin{definition}[Corridors, ribbons and scales]\label{def:ribbons}
A corridor is a vertical cell $\Gamma(z)-w(z)<y<\Gamma(z)+w(z)$. A horizontal ribbon about its midpoint seed is the image of $(t,z)\mapsto(\Gamma(z)+t,z)$ on a specified rectangle; it is contained when its transverse half-width is at most $w(z)$. We call it vertically long when the $z$ length is the larger side, and wide when the horizontal width is the larger side.

For a moderate corridor with $w\le w(z)\le2w$, $z$ length $H$, and residual bound $|\alpha|\le2h$, define the geometric lengths by
\begin{equation}\label{geo:assigned-scales}
L_{\rm geom}=\max\{w,H,hH\},\qquad W_{\rm geom}=wH/L_{\rm geom}.
\end{equation}
Choose the assigned dyadic scales by the stated polynomial conditioning and rounding, so that $L\sim_P L_{\rm geom}$ and $W\sim_P W_{\rm geom}$, and put $S=LW$ and $K=L/W$. Thus $S\sim_P wH$. Coordinate side lengths are distinguished from these scales; for a vertically long ribbon they are $L_0=H$, $W_0=w$, with $K_0=H/w$.

For a parametrization $\Psi(s,t)$, a bounded coframe means
\[
\begin{pmatrix}\Psi^*\theta\\\Psi^*dz\end{pmatrix}=C_\Psi(s,t)\begin{pmatrix}ds\\dt\end{pmatrix},\qquad \|C_\Psi\|_{\rm op}+\|C_\Psi^{-1}\|_{\rm op}\le P^C.
\]
\end{definition}
\begin{lemma}[Narrow seeds tangent to the null direction]\label{geo:narrow-null}
Suppose a seed in the buffered set $A$ is given by $y=\Gamma(z)$ on an interval of length $H$, with $|\alpha|\le1$.

\noindent\textup{(i)} If all assigned representatives lie within horizontal distance $P^Cm_0$ of the seed, the horizontal ribbon of Definition~\ref{def:ribbons} with $|t|<c_Pm_0$ is contained in $E$, has bounded coframe, and has area $\sim_P Hm_0$.

\noindent\textup{(ii)} For $H\gtrsim_Pm_0$ its vertical length is at least its horizontal width, up to powers of $P$.
\end{lemma}
\begin{proof}
A horizontal ruling has intrinsic length $|t|$, so the buffer gives containment. The coframe on the displayed map is $dt+(\Gamma'-d(\Gamma+t,z))\,dz,dz$. Its first coefficient is bounded by $1+P^C|t|$, using $Yd=u$, and its determinant has absolute value one. The ordinary area Jacobian also equals one. Thus both matrix bounds and the area claim follow. The initial horizontal bridge from any assigned representative has fixed polynomial length. 
\end{proof}

\begin{lemma}[Short-projection caps]\label{geo:caps}
Let a connected fixed-format corridor in $A$ have $z$ projection of length less than the fixed cutoff $\tau_0$.

\noindent\textup{(i)} There is a section $z=z_0$ in that interval onto which the actual characteristic projection is a horizontal interval of length $\ell$. The corridor admits a parabolic cap whose horizontal seed contains that interval and has length comparable to $\max(m_0,\ell)$. Its scales satisfy $L\sim_P\max(m_0,\ell)$ and $W\sim_Pm_0$, with $L\ge W$. Every characteristic used has time length at most $\tau_0$.

\noindent\textup{(ii)} All original labels of the corridor are assigned to this cap; their number is at most $P^C\Lambda LW$, and the collar area is at least $P^{-C}LW$.

\noindent\textup{(iii)} The assignment is Borel in the coefficients.
\end{lemma}
\begin{proof}
Let $z_0$ be the midpoint of the corridor's $z$ projection. For the fixed phase and selected root branch, let $\pi_{z_0}(y,z)$ be the $y$ coordinate reached by the actual characteristic through $(y,z)$ at height $z_0$, after time $z_0-z$. Lemma~\ref{geo:local-flow} applies during the whole time interval. The image is an interval $I$ by connectedness and continuity. Every point of $I$ is reached from $A$ along a characteristic of length at most $\tau_0$, so the model-box test in Lemma~\ref{geo:core} retains all restrictions on the original phase, root pivots, and \eqref{geo:small-rates}, with buffer at least $M-\tau_0$. If $|I|<m_0$, enlarge it horizontally to length $m_0$; each added point is at intrinsic distance at most $m_0$ from $I$. The enlarged seed remains in the same certified outer boxes and has buffer at least $M-\tau_0-m_0>4m_0$. A horizontal interval has fixed algebraic format, unit flux, and zero vertical speed, so Lemma~\ref{geo:collar} gives a contained algebraic collar without requiring this seed to lie in $A$. Its area is at least $P^{-C}LW$; this collar is used for the coefficient-volume estimate.

For count only, use the actual local flow tube from the enlarged interval for $|t|\le2\tau_0$. It contains all assigned points and has parameter coframe $J\,da,dt$, with $J$ and inverse polynomially bounded. Cover its parameter rectangle by $O_P(\ell m_0+m_0^2+1)$ unit boxes. Their relative normalized frames are polynomially bounded by the local model-box comparison. Each box contains at most $P^C\Lambda$ original labels by family packing and the normalization multiplicity. The resulting bound is $P^C\Lambda LW$.

Selected regular roots and bounded-time flows depend continuously on their conditioned parameters. Parameter-compatible semialgebraic cells provide a countable family of corridor points that is dense in each corridor and whose members are Borel functions of the coefficient vector, with the bin and root-branch labels fixed. The infimum and supremum of $\pi_{z_0}$ over this family are Borel functions of the same parameters, and they equal the endpoints of $I$. Fixed enlargement and rounding rules consequently give Borel cap labels. 
\end{proof}

\begin{lemma}[Fixed-frame ranges for transverse collars]\label{geo:collar-endpoints}
Let a corridor have midpoint seed $y=\Gamma(z)$, horizontal width comparable to $w$, and residual $\alpha$ of constant nonzero sign. Suppose either $w\lesssim_Pm_0$ and $|\alpha|>1$, or its $z$ length $H\ge \tau_0$ and $|\alpha|\ge\max(2,C_Pw/m_0)$. Set $\ell=\int|\alpha|\,dz$ and $L\sim_P\max(m_0,\ell)$.

\noindent\textup{(i)} The seed has a flux parametrization satisfying Lemma~\ref{geo:collar}. Relative to the fixed original frame of each assigned label, all certificate endpoints have ranges
\begin{equation}\label{geo:collar-ranges}
|f|\le P^CL,\qquad |d|\le P^CL,\qquad |y|+|y-df|\le P^CL^2.
\end{equation}

\noindent\textup{(ii)} The candidate $(\eta,f)$ area, where $\eta=y-df$, divided by the collar area where the bounds hold is at most $P^CL^2$.

\noindent\textup{(iii)} If the endpoint shear is held fixed in the coefficient map and integrated separately, its integration interval has length at most $P^CL$.
\end{lemma}
\begin{proof}
Orient the seed so that its signed flux becomes positive and parameterize it by $ds=\theta(\gamma')\,dz$. Then $\theta(\gamma_s)=1$ and $|z_s|=|\alpha|^{-1}\le1$. The original core supplies the local bounds and buffer at every seed point, so Lemma~\ref{geo:collar} applies over the complete flux interval. In the high-residual case, $\ell\ge C_PwH/m_0\gtrsim_Pw$ and the original corridor area $wH$ is at most $P^Cm_0\ell$. In the narrow case the initial horizontal bridge has fixed polynomial length.

An original assigned representative first follows the contained horizontal segment to the seed. Its length is $O_P(w)\le P^CL$ and it changes $d$ by at most $P^CL$. Along the seed, with $\nu=z_s$, the identities $D'=u+\nu V$, $b'=\nu$, and $a'=1+D\nu$ imply successively that $D,b$ range over intervals of length $P^CL$, and $a$ over an interval of length $P^CL^2$, relative to the fixed frame. The initial fixed normalization and the bounded transverse collar width change these estimates only by powers of $P$. This proves \eqref{geo:collar-ranges}. The resulting loss is bounded by $\frac{\text{candidate }(\eta,f)\text{-area}}{\text{good collar area}}\le\frac{P^CL^3}{P^{-C}L}\lesssim_P L^2$. The separate shear range is the second bound in \eqref{geo:collar-ranges}. For a cap, Lemma~\ref{geo:caps} supplies the buffered horizontal seed; the initial characteristic comparison has bounded time and the same estimates apply.
\end{proof}

\begin{lemma}[Prepared collar seed jets]\label{geo:collar-jets}

\leavevmode\noindent\textup{(i)} On a subset of a collar of comparable area, all prescribed finite seed derivatives can be prepared jointly. In the actual normalized frame at an evaluation point, its centered tangential germ can be written
\[
Y=A(s)+D_*(s)t+\tfrac12V_*(s)t^2,\qquad Z=B_*(s)+t,\qquad A(0)=B_*(0)=0,
\]
with the convention $A_j=A^{(j)}(0)$, $(B_*)_j=B_*^{(j)}(0)$, $(D_*)_j=D_*^{(j)}(0)$, and $(V_*)_j=V_*^{(j)}(0)$ for raw seed derivatives, without factorial division. On one common prepared subset, retain the unrounded seed derivatives through order $2k+1$ and the full scalar jets through total order $2k$, as in Lemma~\ref{geo:preparation}. For the seed indices in these finite ranges, the bounds are
\begin{equation}\label{geo:collar-seed-widths}
\begin{aligned}
|A_1|&\le P^C,& |A_j|&\le P^CL^{2-j}\quad(j\ge2),\\
|(B_*)_j|+|(D_*)_j|&\le P^CL^{1-j}\quad(j\ge1),& |(V_*)_j|&\le P^CL^{-j}\quad(j\ge0).
\end{aligned}
\end{equation}

\noindent\textup{(ii)} The value $D_*(0)$ is inverse-polynomially small after the model defect $\delta$ is chosen. On the fixed polynomially bounded transverse window, the linear part satisfies
\[
\begin{pmatrix}A_1&D_*(0)\\(B_*)_1&1\end{pmatrix}=\begin{pmatrix}1&0\\\nu&1\end{pmatrix}+E,\qquad \|E\|_{\rm op}\le P^C\delta,\qquad |\nu|\le1.
\]
Only $A(0)$ and $B_*(0)$ have been set to zero; $D_*(0)$ and $V_*(0)$ retain their actual values.

\noindent\textup{(iii)} The Jacobian after recentering $\mathcal J=\det D\Psi$ remains \eqref{geo:collar-jacobian}. On the same prepared set, for every prescribed finite order $q$,
\begin{equation}\label{geo:reciprocal-collar-jacobian}
\left|\partial_s^i\partial_t^j\mathcal J^{-1}\right|\le P^CL^{-i}\qquad(i+j\le q).
\end{equation}
\end{lemma}
\begin{proof}
On the actual collar before endpoint recentering, apply Lemma~\ref{geo:preparation} simultaneously to the unrounded algebraic seed data through order $2k+1$, the scalar lists specified in the coefficient estimate through total order $2k$, and $\mathcal J^{-1}$ with all derivatives required by the prescribed order $q$ and the finite chain rule. The seed ranges follow from Lemma~\ref{geo:collar-endpoints}, and $V$ is bounded by the admissibility conditions. Lemma~\ref{geo:collar} gives $1/2\le\mathcal J\le3/2$. Every differentiated reciprocal-Jacobian expression is algebraic in the original seed parameter and $t$, with the nonvanishing flux and Jacobian denominators retained. Preparation on the longitudinal scale $L$ and the fixed polynomial transverse scale gives \eqref{geo:reciprocal-collar-jacobian}; the transverse scale factors are absorbed only into $P^C$. Subtracting a frozen endpoint shear in the output has determinant one, and translating the parameter origin replaces $\mathcal J$ only by its translate. Thus endpoint recentering preserves the actual Jacobian and these prepared estimates. At a point with transverse coordinate $t_0$, write $t=t_0+\tau$ and subtract its actual null slope $d_0$ from the horizontal coordinate. The new seed functions, before subtracting their constant terms where appropriate, are
\[
A=a+Dt_0+\tfrac12Vt_0^2-d_0b,\qquad B_*=b,\qquad D_*=D+Vt_0-d_0,\qquad V_*=V.
\]
Here $|t_0|\le P^C$ and $|d_0|\le P^CL$. The bounds for derivatives of order at least two, and those for $B_*,D_*,V_*$, follow from the displayed formulas and preparation. The stronger bound for $A_1$ follows from the bounded coframe in Lemma~\ref{geo:collar}, since $d_0$ is the actual slope at the evaluation point. The characteristic comparison in that lemma shows $D+Vt_0-d_0=O_P(\delta t_0^2+\delta^2|t_0|^3)$ and gives the stated linear part. The coefficient atlas and its finite seed rounding are proved separately in Section~\ref{surf:section}; this lemma supplies their geometric seed ranges.
\end{proof}

\begin{lemma}[Width and residual bands; corridor packing]\label{geo:bands-packing}
Let \(A\subset E\) be fixed-format semialgebraic sets on a selected null branch, with intrinsic buffer at least \(1\) from \(\overline A\) to the complement of \(E\), \(|u|\le P^C\), and \(|y|+|z|+|d|\le T\), where \(T\ge2\). Assume the local model boxes give intrinsic neighborhoods of an inverse-polynomial radius, with polynomial upper and lower area bounds and polynomial relative-frame bounds on fixed enlargements. Assume the original-label count is at most \(\Lambda M^C\) in each relative-frame neighborhood with bound \(M\ge2\), as in~\eqref{an:local-packing}; include any prescribed normalization or relocation multiplicity in \(\Lambda\). Fix a polynomial cutoff \(\tau\ge1\) and a fixed exponent \(C_0>0\), and put
\[
w_*=P^{-C_0}\tau.
\]
The constants below satisfy
\[
0<c_P\le1/4,\qquad C_P\ge2,\qquad c_P^{-1}+C_P\le P^C.
\]

\noindent\textup{(i)} There is a parameter-compatible Borel partition into at most \(P^C(1+\log T)^C\) graph corridors, arcs and points. For a corridor with midpoint \(\Gamma(z)\) and half-width \(w(z)\), use the residual \(\alpha\) of Definition~\ref{def:seeds-collars}. Its width is either at most \(w_*\), with \(|\alpha|\le1\) or signed \(|\alpha|>1\), or belongs to a dyadic band
\[
w_0\le w(z)\le2w_0,\qquad w_0\ge w_*.
\]
In each latter band, the residual is in one low range, a signed dyadic intermediate range, or a signed high range, with cutoffs
\[
h_{\mathrm{low}}=c_Pw_0/T,\qquad h_{\mathrm{high}}=\max(2,C_Pw_0/\tau).
\]

\noindent\textup{(ii)} For any such corridor \(C\) of projection length \(H\) and half-width at most \(2w_0\), including \(w_0=w_*\) in the small-width case, put \(\ell_{\mathrm{int}}=H+\int|\alpha|\,dz\). This is the intrinsic midpoint length and is at most \(P^CT^2\). If \(\operatorname{Per}_{\theta}(C)\) denotes its intrinsic boundary length and \(|C|\) its ordinary area, then
\[
|C|+\operatorname{Per}_{\theta}(C)\le P^C(\ell_{\mathrm{int}}+w_0+Hw_0),\qquad
\#\{\text{assigned labels in }C\}\le P^C\Lambda\bigl(1+|C|+\operatorname{Per}_{\theta}(C)\bigr).
\]
For an arc of intrinsic length \(\ell_{\mathrm{int}}\), the count is at most \(P^C\Lambda(1+\ell_{\mathrm{int}})\). Labels are counted with their original multiplicities.
\end{lemma}
\begin{proof}
Take a parameter-compatible vertical cell decomposition \cite[Theorems 5.6 and 14.16]{BPR}, retaining the selected root branches and fixed ties. Merge all widths below \(w_*\), and split their midpoint arcs at \(\alpha=\pm1\). Above \(w_*\), use dyadic widths, the displayed residual cutoffs, and signed dyadic intermediate bands. There are \(O_P(1+\log T)\) choices of each dyadic parameter, since \(w_0\ge P^{-C_0}\tau\) and \(h_{\mathrm{high}}/h_{\mathrm{low}}\le P^CT/\tau\). Each selected set has fixed format and hence boundedly many interval components. Zero residuals and arbitrarily large residuals require no further bands. Lower-dimensional cells are graph arcs, horizontal intervals and points.

On a side curve \(\Gamma\pm w\), the flux density is
\[
\alpha\pm w'+d(\Gamma,z)-d(\Gamma\pm w,z).
\]
The last term is at most \(P^Cw_0\). Fixed format and the width range give \(\operatorname{Var}(w)\lesssim_Pw_0\), including when the width tends to zero. The sides, end segments and area therefore have the stated bound. Also \(\ell_{\mathrm{int}}\le\operatorname{Var}(\Gamma)+(1+T)H\le P^CT^2\), by the bounded number of monotonicity intervals.

Choose a sufficiently small inverse-polynomial intrinsic radius in the local model boxes. A maximal disjoint family of smaller neighborhoods centered at labels away from the boundary has at most \(P^C|C|\) members, by their area lower bound; fixed enlargements cover these labels. Labels near the boundary lie in enlarged neighborhoods of an intrinsic mesh on its boundedly many arcs, with at most \(P^C(1+\operatorname{Per}_{\theta}(C))\) points. Each enlarged neighborhood has polynomial relative-frame bounds and hence at most \(P^C\Lambda\) original labels. The same mesh proves the arc estimate. All these neighborhoods are taken in the original outer domain \(E\); no buffer within a width or residual band is assumed.
\end{proof}

\begin{lemma}[Moderate algebraic corridors]\label{geo:moderate}
Let a vertical cell of $A$ have midpoint $\Gamma(z)$ and half-width $w(z)$ satisfying $w\le w(z)\le2w$ on an interval of length $H$, where $w,H\ge P^{-C}\tau_0\ge1$. Suppose $|\alpha|\le2h$ and $h\le P^Cw/\tau_0$.

\noindent\textup{(a)} The whole cell is assigned to its contained horizontal ribbon
\begin{equation}\label{geo:moderate-ribbon}
(y,z)=(\Gamma(z)+\xi,z),\qquad |\xi|<c_Pw.
\end{equation}
It has area $\sim_P wH\sim_P S$ and the assigned scales \eqref{geo:assigned-scales}. Then $L\ge W\ge1$, the number of assigned labels is at most $P^C\Lambda S$, and each assigned fixed frame reaches the ribbon along a contained path of length $O_P(L)$.

\noindent\textup{(b)} Set $e_{\rm rel}(\xi,z)=\Gamma^\prime(z)-d(\Gamma(z)+\xi,z)$. The endpoint area ratio below is the candidate endpoint area divided by the lower bound for the common good endpoint area, as specified in Definition~\ref{def:point-factor}. For the prescribed orders $j\ge2$, the following bounds hold on the common prepared subset.

\textup{(i)} If $H\ge w$, the coordinate scales are $L_0=H$, $W_0=w$, and $K_0=H/w$, and
\[
\begin{aligned}
|\Gamma^{(j)}|&\le P^CH^{3-j},&|e_{\rm rel}|&\le P^CH,\\
\text{endpoint area ratio}&\le P^CH^2K_0.
\end{aligned}
\]
\textup{(ii)} If $w\ge H$, the absolute seed bound $D_*$ and the higher seed derivatives satisfy
\[
\begin{aligned}
D_*&\le P^CwH,&|\Gamma^{(j)}|&\le P^CD_*H^{1-j},\\
|e_{\rm rel}|&\le P^Cw,&\text{endpoint area ratio}&\le P^CH^2.
\end{aligned}
\]
\end{lemma}
\begin{proof}
The ribbon is contained in the original cell by its width lower bound. Its ordinary area element is $d\xi\,dz$. Its coframe discrepancy obeys
\[
|\Gamma'-d(\Gamma+\xi,z)|\le |\alpha|+\left|\int_0^\xi u(\Gamma+t,z)\,dt\right|\le P^C(h+w)\le P^Cw.
\]
For $H\ge w$ this is an allowed $P^CH$ discrepancy; for $w\ge H$ it is an allowed $P^Cw$ discrepancy. All full admissibility conditions hold because the ribbon lies in $A\subset E$.

The maximum in \eqref{geo:assigned-scales} is at least $\sqrt{wH}$. Each of its three terms is at most $P^CS/\tau_0$, using the lower bounds for $w,H$ and the upper bound for $h$. Thus $L\ge W$ and $W\ge P^{-C}\tau_0$. The fixed local width was chosen large enough that the latter quantity is at least one. From an assigned anchor follow a horizontal bridge to the midpoint, follow the midpoint, and then follow another horizontal bridge into the ribbon. Their combined length is at most $P^C(w+H+hH)\le P^CL$. They lie inside the original cell. 

Lemma~\ref{geo:bands-packing} bounds the cell's area and intrinsic boundary length by \(P^C(H+hH+w+wH)\le P^CS\), since \(w,H\ge1\) and \(h\le P^Cw/\tau_0\). Its packing estimate gives the required \(P^C\Lambda S\) original labels.

For the vertical endpoint bounds, use the unchanged fixed frame as origin and absorb its bounded initial comparison into $P^C$. Along the horizontal bridge, $d$ changes by $O_P(w)$. Along the midpoint, $d'=v+\alpha u$, so $|d'|\le P^C(1+h)\le P^CH$. Consequently
\[
|z|\le P^CH,\qquad |d|\le P^CH^2,\qquad |y|+|y-dz|\le P^CH^3.
\]
Consequently, $\frac{\text{candidate }(\eta,f)\text{-area}}{\text{good ribbon area}}\lesssim_P\frac{H^4}{wH}=H^2K_0$. The algebraic scalar $\Gamma'$ has magnitude at most $P^CH^2$. Lemma~\ref{geo:preparation} therefore gives $|\Gamma^{(j)}|\le P^CH^{3-j}$ at a positive proportion of common endpoints. Subtracting the actual endpoint slope leaves the already bounded relative discrepancy.

For the wide bounds, put $D_0=w+hH\le P^CwH$. The same paths give $|d|\le P^CD_0$, $|z|\le P^CH$, and $|y|+|y-dz|\le P^CD_0H$. Thus $\frac{\text{candidate }(\eta,f)\text{-area}}{\text{good ribbon area}}\lesssim_P\frac{D_0H^2}{wH}=\frac{D_0H}{w}\lesssim_P H^2$. Preparation of $\Gamma'$ gives $|\Gamma^{(j)}|\le P^CD_0H^{1-j}$. For the fixed coefficient set one may replace $D_0$ by the larger parameter-independent bound $D_*=P^CwH$. All these assertions hold for every fixed frame assigned to the whole cell, including anchors outside the retained central ribbon. No long characteristic transport is used.
\end{proof}

\begin{definition}[Geometric certificates]\label{def:geometric-certificate}
A geometric certificate specifies the contained patch and simultaneous derivative bounds used for the coefficient-volume estimate of an assigned piece. It consists of a specified contained patch of fixed algebraic format, its scale and endpoint data, and a common good evaluation set with the prescribed lower area bound. It retains the original phase's complete coefficient and physical admissibility, the pivots of the selected chart, and the scalar lists and derivative orders specified in its coefficient estimate, with all their bounds holding on that common set. Any additional rate window is specified in the application.
\end{definition}
\begin{theorem}[Common null assignment]\label{geo:common-allocation}

\leavevmode\noindent\textup{(i)} The original normalized labels in $A$ have a Borel partition $\{\Pi\}$, chosen for all labels of each phase before a label is fixed for coefficient integration. Let $T$ bound the absolute upper-affine coordinates and their polynomially conditioned normalizations. The number of pieces and, for each piece $\Pi$, its dyadic scales and assigned-label count satisfy
\[
\begin{gathered}
L\ge W\ge1,\qquad S=LW,\qquad K=L/W,\\
\#\{\Pi\}\le P^C(1+\log T)^C,\qquad \#\{\text{labels in }\Pi\}\le P^C\Lambda S.
\end{gathered}
\]

\noindent\textup{(ii)} Each piece has a geometric certificate in the sense of Definition~\ref{def:geometric-certificate}, whose patch lies in $E$. Each assigned fixed frame reaches this patch along a path of length at most $P^CL$ on which the full admissibility conditions hold. The seed derivative bounds and fixed-frame endpoint bounds proved above hold on a common prepared subset. The patch of each certificate is one of the following: a moderate horizontal ribbon from Lemma~\ref{geo:moderate}; a fixed-width horizontal ribbon whose vertical length is larger; a transverse parabolic collar whose seed is parametrized by $\theta(\gamma_s)=1$ and $|z_s|\le1$; or a cap from Lemma~\ref{geo:caps}.

\noindent\textup{(iii)} For every $\eta>0$,
\begin{equation}\label{geo:inverse-area-weight}
\sum_{\Pi}S_\Pi^{-\eta}\le P^C(1+\log T)^C.
\end{equation}
The permitted physical coordinate range gives $\log T\lesssim (1+\log P+\log d_1^{-1})$.
\end{theorem}
\begin{proof}
\par\medskip\noindent\textit{Step 1. Assignment of common cells.}\par\nobreak\smallskip
Apply Lemma~\ref{geo:bands-packing} to the full-parent core with \(\tau=\tau_0\), using the local comparisons of Lemmas~\ref{geo:model-box}--\ref{geo:core}. Write \(w\) for its width-band scale \(w_0\). Work on the selected root branches, retain continuation through artificial cuts, and assign each original label to the first cell in the fixed ordering that contains its normalized representative and satisfies the selected branch conditions. Lower-dimensional arcs use the small-width construction below; isolated points use caps.

\par\medskip\noindent\textit{Step 2. Small-width and short-projection components.}\par\nobreak\smallskip
For a small-width component with \(|\alpha|\le1\), use Lemma~\ref{geo:caps} if its projection length is less than \(\tau_0\); otherwise choose \(c_Pm_0\le\tau_0\), and the original buffer and Lemma~\ref{geo:narrow-null} give a vertically long horizontal ribbon of that width. For signed \(|\alpha|>1\), use the flux-parametrized parabolic collar of Lemma~\ref{geo:collar}, or a cap if its seed length is below the local width. The original labels reach these certificates by horizontal bridges of polynomial length. Their certificate areas dominate \(1+\ell_{\mathrm{int}}+w_*+Hw_*\), up to powers of \(P\), so Lemma~\ref{geo:bands-packing} gives the required counts.

In every remaining width and residual component whose projection length is less than \(\tau_0\), use Lemma~\ref{geo:caps}. Its containment uses the original full-core buffer.

\par\medskip\noindent\textit{Step 3. High-residual components.}\par\nobreak\smallskip
For a high-residual component with $H\ge \tau_0$, let $\ell=\int|\alpha|\,dz$. Then $\ell\gtrsim_PwH/\tau_0$ and $\ell\gtrsim_PH$. Its flux seed has $|z_s|\le1$. Assign the whole component to the fixed-width parabolic collar with scales comparable to $\max(\tau_0,\ell),\tau_0$. Its area dominates both the cell area and its intrinsic boundary length. The initial horizontal bridge has length $O(w)\le P^C\ell$. Lemmas~\ref{geo:collar}--\ref{geo:collar-jets} give the certificate, path bound, and endpoint ranges, and Lemma~\ref{geo:bands-packing} gives the packing bound.

\par\medskip\noindent\textit{Step 4. Moderate components.}\par\nobreak\smallskip
Every remaining component satisfies Lemma~\ref{geo:moderate}; for the merged low-residual case use its upper bound as $h$. Assign it directly to its actual algebraic horizontal ribbon. Round its geometric scales dyadically. The coordinate aspect of this ribbon need not equal its assigned aspect $K$; both are retained in their respective estimates.

\par\medskip\noindent\textit{Step 5. Counts and measurability.}\par\nobreak\smallskip
Lemma~\ref{geo:bands-packing} gives the piece count and the bound \(P^CT^2\) for midpoint intrinsic lengths, so the selected scales have only permitted logarithms. Since every \(S\ge1\), the count implies \eqref{geo:inverse-area-weight}. Enlarge and prepare all finite derivative lists jointly by Lemma~\ref{geo:preparation}.

All choices except flux primitives and the local cap endpoints are semialgebraic. Those remaining quantities are Borel by monotone integration and the endpoint argument in Lemma~\ref{geo:caps}. Fixed ties therefore give a common Borel partition in the original coefficient parameters. No step selects an anchor before making this partition. Finally, in a fixed RQ bin, bounded physical center and matrix entries, together with the inverse diagonal matrix bound, put the original group coordinates in a box of radius $P^Cd_1^{-1}$. The polynomially bounded normalization multiplies this bound by a further fixed power. This gives the asserted estimate for $\log T$ and completes the proof.
\end{proof}


\subsection{Gauss assignment}\label{geo:gauss-subsection}

\begin{lemma}[Buffered planar strips]\label{geo:planar-strips}
Let $A\subset E\subset\mathbb R^2$ be semialgebraic sets of fixed format, with $E$ open and bounded, and suppose the closure of $A$ has Euclidean distance at least $\delta_P=P^{-C}$ from $E^c$. Suppose $E$ lies in a box of radius $T$, and each square of side $\delta_P$ contains at most $P^C\Lambda$ indices. Put $T_+=\max\{1,T\}$.

\noindent\textup{(i)} There is a measurable assignment of the labels in $A$ to at most $P^C(1+\log T_+)^C$ pieces with scales $L\ge W\ge1$, up to powers of $P$ in the comparison with actual lengths. Each piece $\Pi$ satisfies
\[
\#\{\text{assigned labels in }\Pi\}\le P^C\Lambda LW,\qquad \operatorname{diam}\Pi\le P^CL.
\]

\noindent\textup{(ii)} It has an associated strip contained in $E$ whose area, after derivative preparation, is at least $P^{-C}LW$. After choosing from a fixed finite collection of coordinate directions, the prepared strip is parameterized by
\[
(s,\Gamma(s)+t),\qquad |\Gamma^{(j)}|\le P^CL^{1-j}\quad(1\le j\le N),
\]
with $s$ range comparable to $L$, $t$ range comparable to $W$, where $N$ is any fixed prescribed order. For a chosen point $s_0$ of the parameter interval, $|\Gamma(s)-\Gamma(s_0)|\le P^CL$.

\noindent\textup{(iii)} The candidate endpoint area relative to any assigned point is at most $P^CL^2$.
\end{lemma}
\begin{proof}
By quantifier elimination \cite[Theorem 2.77]{BPR}, the distance function $\rho(x)=\operatorname{dist}(x,E^c)$ is semialgebraic: its graph is expressed by a minimum over the closed complement, with an a priori bounded search region for points in the box under consideration. Divide $A$ according to $W_0\le\rho<2W_0$, with dyadic $W_0\ge c\delta_P$. Work with the corresponding components of the enlarged bands $W_0/2<\rho<4W_0$, and resolve overlaps by a fixed order. Their number at each scale is bounded by fixed-format cell decomposition. Let $C$ be such a component containing assigned labels, and put $L_0=\max(W_0,\operatorname{diam}C)$. There are only $O_P(1+\log T_+)$ distance scales.

The component has area $\lesssim L_0W_0$. Indeed it lies in a ball of radius $CL_0$, and every point is within $4W_0$ of the boundary of $E$. The portion of that boundary in a slightly larger ball is a union of boundedly many semialgebraic arcs and points. Split the arcs into pieces monotone in each coordinate. Each piece has length bounded by the sum of the variations of its coordinates, hence by $CL_0$. Its $4W_0$ neighborhood has area $\lesssim L_0W_0+W_0^2$, by covering an arclength mesh with squares of side $W_0$. Sum the bounded number of arcs. The same argument controls a local-width enlargement of $C$. Cover that enlargement by squares of side $\delta_P$; its square count is at most $P^CL_0W_0$. Local packing proves the packing bound.

If $\operatorname{diam}C\le CW_0$, a ball of radius $cW_0$ around any assigned point lies in $E$ and supplies the strip. Otherwise join two points of $C$ at distance comparable to $L_0$ by the path in Lemma~\ref{geo:paths}. Its coordinate variation is $\lesssim L_0$. Split it into boundedly many graph arcs, assigning each arc a coordinate direction in which its tangent projection has magnitude at least $1/2$. At least one such graph arc has longitudinal projection of length $\gtrsim L_0$; its slope is bounded. Every point of the arc has distance at least $W_0/2$ from $E^c$, so its transverse enlargement by $cW_0$ is contained in $E$. Lemma~\ref{geo:preparation}, applied to the centerline, gives all specified derivative bounds on a subset of comparable longitudinal measure. That good set has a bounded number of interval components, so one may retain an interval of comparable length after decreasing constants. Its strip has area $\gtrsim L_0W_0$.

All points of $C$ and of the contained strip just constructed lie in a ball of radius $CL_0$. Thus relative to every assigned anchor the possible endpoint region has area $\lesssim L_0^2$. Replace lengths below one by one, and round larger lengths dyadically. Since $W_0\ge c\delta_P$, this changes all comparisons only by powers of $P$. Cell and path choices made jointly in the coefficient parameters, finite derivative tests, and fixed tie rules give a measurable assignment. The number of components, not the number of indices, controls its multiplicity.
\end{proof}

\begin{theorem}[Gauss assignment]\label{geo:gauss}

\leavevmode\noindent\textup{(i)} For every normal order in Lemma~\ref{geo:rank-entry}, the curvature branch of Lemma~\ref{geo:chart-split} admits an assignment to at most $P^C(1+\log\max\{1,T\})^C$ pieces in the coordinates $(p,z)=(c_y,z)$ of the selected graph. Each piece has $L\ge W\ge1$, $S=LW$, $K=L/W$, and at most $P^C\Lambda S$ assigned labels.

\noindent\textup{(ii)} Its associated strip lies in the admissible set and has area at least $P^{-C}S$ after preparation. The centerline derivatives are bounded by $P^CL^{1-j}$ for $1\le j\le N$, for every prescribed finite order $N$; the order-zero bound is understood for the centerline minus its value at a chosen parameter point.

\noindent\textup{(iii)} Relative to each assigned fixed frame, the candidate endpoint set has area at most $P^CL^2$. The assertion includes both elliptic and hyperbolic graph Hessians and retains all admissibility conditions of the parent phase.
\end{theorem}
\begin{proof}
\par\medskip\noindent\textit{Step 1. Root branches.}\par\nobreak\smallskip

Put $G=\partial_x^{m-1}F$, $\kappa=\partial_x^mF|_c$, $r_0=c_{yy}$, and $p=c_y$. At fixed $(p,z)$ solve
\begin{equation}\label{geo:gauss-equations}
G(x,y,z)=0,\qquad G_y(x,y,z)+pG_x(x,y,z)=0
\end{equation}
for $(x,y)$. At a regular graph point, with $p,z$ fixed, $\det\frac{\partial(G,G_y+pG_x)}{\partial(x,y)}=-\kappa^2r_0$. The first equation has degree at most $D=k-m+1$, and the second degree at most $D-1$. The isolated-zero form of B\'ezout's theorem \cite[Corollary 1.3]{low:roots} gives at most $D(D-1)$ regular solutions. A positive-dimensional common component cannot contain a point with this nonzero Jacobian. Thus the number of retained root branches is bounded without omitting any regular branch.

\par\medskip\noindent\textit{Step 2. The frame and its body differential.}\par\nobreak\smallskip

Write $y=Y(p,z)$ for the inverse of $p=c_y(y,z)$ on the selected branch. In the formulas below, $c$, its derivatives, and $r_0=c_{yy}$ are evaluated at $(Y(p,z),z)$; their dependence on the fixed phase $F$ is suppressed. Set
\[
d=-c_{yz}/r_0,\qquad q_0=c_z(Y(p,z),z),\qquad s_0=\det D^2c/r_0.
\]
The normalized affine frame is
\begin{equation}\label{geo:gauss-frame}
h_{p,z}(x,u,v)=\bigl(x+c+pu+(q_0+dp)v, Y+u+dv, z+v\bigr).
\end{equation}
Implicit differentiation gives $Y_p=r_0^{-1}$ and $Y_z=d$. Using homogeneous affine matrices for $h=h_{p,z}$, its body differential is
\begin{equation}\label{geo:gauss-body}
h^{-1}dh=A\,dp+C\,s_0\,dz+D(d_p\,dp+d_z\,dz)+T_y r_0^{-1}\,dp+T_z\,dz.
\end{equation}
In the group coordinates \eqref{geo:group} one has $a=p$, $b=q_0+dp$, $e=Y$, and $f=z$. The body $C$ component is $db-a\,dd=dq_0+d\,dp=s_0\,dz$, since $(q_0)_p=-d$. The body $T_y$ component is $dY-d\,dz=r_0^{-1}\,dp$, and the body $T_x$ component is $dc-p\,dY-q_0\,dz=0$. This proves \eqref{geo:gauss-body}.

In the normalized tangential variables $(u,v)$ the mixed graph curvature is zero. Differentiating $d$ gives
\[
d_p=-c_{uuv}/r_0^2,\qquad d_z=-c_{uvv}/r_0.
\]
The coefficients of the original phase and $\kappa^{-1}$ bound these graph derivatives by repeated implicit differentiation, so every body coefficient in \eqref{geo:gauss-body} is bounded by $P^C$. Integrating an upper-unipotent matrix equation uses iterated integrals only through order three. Relative frames along paths of polynomially bounded length consequently have polynomial bounds independent of an absolute shear coordinate. The quadratic reference models and the same order of fixing the model box and then choosing the coefficient error as in Lemma~\ref{geo:model-box} provide a full-parent buffer of a prescribed local polynomial radius in the $(p,z)$ plane. The exact model is quadratic because $r_0$ is conditioned; no sign condition is imposed on $\det D^2c$.

\par\medskip\noindent\textit{Step 3. Covariance and scalar bounds.}\par\nobreak\smallskip

Here and below, $[m]G$ denotes the coefficient of the monomial $m$ in $G$; for Taylor series, coefficients are taken at the stated evaluation point. The map $\pi_{\mathrm G}:\mathcal U\to\mathbb R^2$, $\pi_{\mathrm G}(g)=(a,f)$, is a homomorphism. Under the upper-affine change of fixed frame $g\in\mathcal U$, the phase becomes $F\circ g$. Its normal graph is the transformed graph of $F$, and the Gauss variables change by $(p,z)\mapsto(p-a,z-f)$. The intrinsic scalar
\[
\Phi_r(p,z)=\partial_x^rF(c(Y(p,z),z),Y(p,z),z)
\]
therefore translates in these coordinates, with an immaterial additive constant when $r=0$. For $r\ge1$ it is an actual coefficient of the normalized parent phase and is bounded by $P^C$. For $r=0$, $\partial_p\Phi_0=r_0^{-1}[u](F\circ h_{p,z})$, $\partial_z\Phi_0=[v](F\circ h_{p,z})$. Hence $|\nabla\Phi_0|\le P^C$. This remains valid when $F_x|_c$ is nonzero: the tangential frame in \eqref{geo:gauss-frame} includes the corresponding $F_x$ terms.

\par\medskip\noindent\textit{Step 4. Cuts, strips, and packing.}\par\nobreak\smallskip

We now justify applying Lemma~\ref{geo:planar-strips} to all branches. On its regular locus, the algebraic solution set of \eqref{geo:gauss-equations} projects locally diffeomorphically to the $(p,z)$ plane and has uniformly finite fibers. Parameter-compatible cell decomposition \cite[Theorems 5.6 and 14.16]{BPR} gives finitely many single-valued root charts over planar cells.

For each artificial cut, consider only the portions within a fixed sufficiently small local distance of the inner admissible set. The inner-to-outer admissibility margin and the implicit function theorem continue each relevant regular branch across these portions. There are at most $D(D-1)$ branches. Split the resulting fixed-format arcs into boundedly many bounded-slope graph pieces, using either $p$ or $z$ as the longitudinal coordinate; use local squares at endpoints. Assign nearby labels to the connected cut components, resolving ties by a fixed rule. The area-plus-boundary argument in Lemma~\ref{geo:planar-strips}, with this fixed transverse width, bounds the number of labels in each resulting neighborhood. If a component is long, one graph subarc has comparable longitudinal range, and its transverse enlargement supplies a contained strip. Use a local square for a shorter component. The component diameter gives the candidate endpoint bound $P^CL^2$, and apply finite derivative preparation on this strip. The number of branches, cut components, and conditioned local choices is bounded by a power of $P$; none is chosen separately for each original label.

The remaining inner admissible set in each root chart is at inverse-polynomial distance from the cuts and the outer boundary, so Lemma~\ref{geo:planar-strips} applies there.

The local frame bound and finite branch count transfer original family packing to the intrinsic squares: each such square corresponds to a polynomially bounded relative-frame neighborhood and contains at most $P^C\Lambda$ original labels. Lemma~\ref{geo:planar-strips} supplies the assignment, count, contained strips, and prepared centerline bounds. Every assigned original label differs by only a bounded initial comparison from a point in the assigned intrinsic component. Translation covariance under $\pi_{\mathrm G}$ makes the relative endpoint candidate area $P^CL^2$ at that unchanged fixed frame. The endpoint-area loss is $\frac{\text{candidate endpoint area}}{\text{good strip area}}\lesssim_P\frac{L^2}{LW}=K$. The selected root charts, cut rules, and planar assignments are made in the coefficient parameters before choosing any anchor, so the assignment is common and Borel. Its piece count and scale ranges have precisely the permitted polynomial and logarithmic dependence.
\end{proof}

Combining Lemma~\ref{geo:rank-entry}, Theorem~\ref{geo:common-allocation}, and Theorem~\ref{geo:gauss} accounts for every label in alternative \textup{(ii)} of Lemma~\ref{geo:rank-entry} and every permitted normal order. The complementary quantitative region defined by the fifth singular value is assigned to one-dimensional pieces in the regular preparation of Section~\ref{reg:section}. In all regions the assignment precedes the fixed-anchor coefficient integral. The present section proves the existence and geometry of the pieces; the complete coefficient-volume estimates for their actual certificates are established in Sections~\ref{reg:section} and~\ref{surf:section}.


\section{Regular coefficient estimates}\label{reg:geometry-section}\label{reg:section}

With the augmented admissibility of Section~\ref{geo:section} on $V_k$, we construct regular pieces, fixed original-coordinate sampling nets and interpolation maps. The uniform coefficient estimate is stated after the setup and proved in Subsection~\ref{reg:uniform}.

\subsection{Regular pieces and their reference frames}\label{reg:common-pieces}\label{reg:preparation}
Use the augmented states and generators of Definition~\ref{def:augmented-state}, writing the physical map as $\ell(x)=p+Bx$. Impose the complete coefficient and physical windows of the analytic reduction and their prescribed enlargements, with
\begin{equation}\label{reg:stop}
\|[h\cdot\nabla J]\|^2+|Bh|^2\ge\kappa_0^2|h|^2\qquad(h\in\mathbb R^3).
\end{equation}
The reciprocal of $\kappa_0$, all fixed window bounds, and all reciprocal quantitative rank bounds are bounded by powers of $P\ge2$. The smallest original physical singular value is $d_1>0$, and the number of original indices in relative affine neighborhoods satisfies \eqref{an:local-packing} with parameter $\Lambda$. All parameter sets and branch selections have bounded semialgebraic format for fixed $k$.

\begin{definition}[Regular curve coordinates]\label{def:regular-coordinates}
Write $R=\mathcal R=R_{\widehat J}$ for the augmented orbit matrix of Definition~\ref{def:augmented-state}, with the generator columns in the order stated there. Let $L_i$ be its five-column complement obtained by deleting column $i$. A five-column chart retains the specified positive lower bound for $\sigma_{\min}(L_i)$. For a curve $g(t)$, use homogeneous affine matrices and write its body velocity as $g(t)^{-1}g'(t)$ in this generator basis. Where the omitted body component is nonzero, the intrinsic parameter makes it one. The fitted remaining coefficients are the least-squares coefficients $v_i=-L_i^+R_i$, where $R_i$ is the omitted column and $L_i^+=(L_i^TL_i)^{-1}L_i^T$ is the left pseudoinverse of $L_i$; they need not equal the exact body rates.
\end{definition}
\begin{definition}[Sampling data]\label{def:sampling}
For fixed states $h_0,\ldots,h_{m-1}$, the complete evaluation map is $\mathcal A_hJ=([J\circ h_i])_i$, retaining all nonconstant coefficients. A complete sampling certificate consists of the sample nodes, a selected nonzero $N_k$-row minor using all coefficient columns, and its stated output weights. Its evaluation rows are sampled coefficients. A sampled coefficient row means such an evaluation row; a fixed linear-combination row means a fixed linear combination of evaluation rows. The additional samples are the $m-1$ samples beyond the fixed first sample. Their nets are fixed independently of the phase in the original upper-affine coordinates.
\end{definition}
Fix $k\ge14$ in the following theorem. Its proof is given in Subsection~\ref{reg:uniform}.

\begin{theorem}[Regular coefficient estimate]\label{reg:regular-tail}
Suppose that the fifth singular value of the six-column augmented differential
\[
\mathcal R=(A\widehat J,C\widehat J,T_x\widehat J,D\widehat J,T_y\widehat J,T_z\widehat J)
\]
is at least $P^{-C}$.

\noindent\textup{(i)} There is a coefficient-measurable assignment of the original admissible indices to bounded pieces and to one-dimensional pieces of dyadic intrinsic length $H\ge1$. A length-$H$ piece contains at most $P^C\Lambda H$ original labels.

\noindent\textup{(ii)} For each original fixed label $j$, the coefficient set $E_j(H)$ that it is assigned to one of these length-$H$ pieces satisfies, in its original normalized coefficient coordinates,
\begin{equation}\label{reg:tail}
\operatorname{vol}_{N_k} E_j(H)\le P^C H^{-\beta_k}\bigl(1+\log P+\log d_1^{-1}\bigr)^C,\qquad \beta_k=\frac{1105}{1026}\chi_k-2k-153>\chi_k+2.
\end{equation}
Bounded pieces have count $P^C\Lambda$ and coefficient volume at most $P^C$. These coefficient bounds contain no positive power of $d_1^{-1}$.
\end{theorem}


\begin{lemma}[Row selection and block-triangular determinants]\label{reg:lem:transfer}

\leavevmode\noindent\textup{(i)} Suppose $N_k$ fixed linear combinations of sample-output rows form a nonsingular square matrix. If every evaluation row used in the $j$th combination has weight at most $w_j$, there is a nonzero $N_k$-row minor of sampled coefficient rows with total output weight at most $\sum_jw_j$.

\noindent\textup{(ii)} Suppose that a full sampling matrix is triangular in an input/output grade and that the selected output multiplicities in each grade equal the input multiplicities. Its determinant is the product of its square diagonal-block determinants. Off-diagonal terms require no smallness assumption.
\end{lemma}
\begin{proof}
Expanding the determinant of the linear-combination rows multilinearly yields a nonzero evaluation minor. Its $N_k$ rows are distinct, and their total weight is at most $\sum_jw_j$.
\end{proof}

\begin{lemma}[Regular decomposition and sampling nets]\label{reg:regular-preparation}\label{low:4:cf:polynomial}
Fix $k\ge2$ and restrict the augmented admissibility condition to $\sigma_5(\mathcal R)\ge P^{-C}$. Fix the number $m$ of prescribed samples before choosing the comparison constants, thresholds, and grids.

\textup{(i)} The original indices admit an assignment to bounded pieces and prepared one-dimensional pieces of dyadic intrinsic length $H\ge1$. Their counts satisfy
\[
\begin{gathered}
\#\{\text{indices in a bounded piece}\}\le P^C\Lambda,\qquad \#\{\text{indices in a length-}H\text{ piece}\}\le P^C\Lambda H,\\
\#\{\text{pieces at each scale}\}\le P^C(1+\log P+\log d_1^{-1})^C.
\end{gathered}
\]
The actual body-velocity charts satisfy \eqref{reg:prepared-derivatives}, and the sampling nets are fixed independently of the phase in the original upper-affine coordinates, as in Definition~\ref{def:sampling}. All comparison constants, reciprocal thresholds, and bounded-grid cardinalities are bounded by powers of $P$.

\textup{(ii)} Suppose each selected evaluation map has a complete $N_k$-row minor with absolute determinant at least $\Delta>0$, and each of the $m-1$ additional sample nets has at most $N$ elements. The coefficient set $E_j$ for each fixed original label $j$ satisfies
\begin{equation}\label{reg:fixed-net-bound}
\operatorname{vol}_{N_k}E_j\lesssim P^C\Delta^{-1}N^{m-1}.
\end{equation}
\equationalias{low:78:h3r:anchor}
The same bound holds when the complete evaluation map $\mathcal A_h$ is compared to one, $\mathcal A_{h_*}$, with such a minor by
\[
P^{-C}\|\mathcal A_{h_*}G\|^2\le\|\mathcal A_hG\|^2\le P^C\|\mathcal A_{h_*}G\|^2\qquad(G\in V_k).
\]
\end{lemma}
\begin{proof}
\par\medskip\noindent\textit{Step 1. Choose and correct a regular chart.}\par\nobreak\smallskip

Use the fixed-format semialgebraic facts in Lemmas~\ref{geo:conditioning}--\ref{geo:paths}. For minima, the coefficient and physical windows are closed and polynomially bounded on bounded $P$ intervals, and their stopped inequalities are closed. Positivity holds throughout these closed windows, or the required lower bound is explicitly retained. The action on $V_k$ is polynomial of bounded degree; the minors have size at most six; and selected rates and their intrinsic derivatives are rational in the group coordinates and coefficient outputs, with retained nonzero minors in their denominators. Thus the format and closed-family hypotheses hold. The derivative bounds with respect to the normalized curve parameter follow from the divided-difference argument below.

The normal columns $N=(A\widehat J,C\widehat J,T_x\widehat J)$ have least singular value at least $P^{-C}$ by Lemma~\ref{geo:normal-pivot}, after enlarging $P$ by a fixed power to accommodate \eqref{reg:stop} and the retained windows.

\paragraph{Chart choice.}

Use Definition~\ref{def:regular-coordinates} and give first priority to the full-rank region, then to a chart with the prescribed lower bound for $\sigma_{\min}(L_d)$, where $d$ denotes the coordinate corresponding to $D$. In the remaining regular region, take $\sigma_5(R)\ge\kappa$, $\sigma_6(R)<\varepsilon$ and $\sigma_{\min}(L_d)<\mu$. If $w$ is a unit right singular vector for $\sigma_6$, then
\[
\sigma_{\min}(L_i)\ge\kappa|w_i|,\qquad |w_D|<\mu/\kappa.
\]
Indeed, for a unit vector $v$ with $v_i=0$, one has $|v-\langle v,w\rangle w|\ge |w_i|$, and $|Rv|\ge\kappa|v-\langle v,w\rangle w|$. Taking the infimum proves the first inequality, and the condition $\sigma_{\min}(L_d)<\mu$ proves the second. The normal-column lower bound forces the $(D,T_y,T_z)$ part of $w$ to have uniformly positive norm.

Allocate first to $|w_{T_z}|\ge\nu$, where $\sigma_{\min}(L_f)\ge\kappa\nu$, with small fitted $D$ component. On the complement $|w_{T_y}|\gtrsim_P1$, so $\sigma_{\min}(L_e)\ge\kappa|w_{T_y}|$; use this chart with small fitted $D,T_z$ components.

Write $v_{T_y,A}$ and $v_{T_y,C}$ for the fitted $A$ and $C$ components in this chart. Translation stopping gives $|v_{T_y,A}|+|v_{T_y,C}|\ge P^{-C}$. Choose the first-shear cutoff $a_0=P^{-C}$ below half this lower bound. Assign $|v_{T_y,A}|\ge a_0$ to the chart omitting $A$; on its complement $|v_{T_y,C}|\ge P^{-C}$, and use the central $C,T_y$ chart. This change of chart is justified by the same deletion inequality: on the selected $L_e$ chart the fitted $A$ rate differs from $w_A/w_{T_y}$ by an error bounded by a power of the reciprocal pivots times $\varepsilon$. Choosing $\varepsilon$ after $a_0$ gives $|w_A|\ge P^{-C}$ on the first-shear branch, so $L_a$ is quantitatively invertible.

The selected near-kernel vector and all chart tests can be chosen semialgebraically, using the gap between $\sigma_5$ and $\sigma_6$ and a fixed sign convention. Choose $a_0$ first, the terminal cutoff next, the residual $D$ cutoff next, and $\varepsilon$ last; this keeps all least-squares errors inside the prescribed outer margins. For example, the fitted $D$ error in the $f$ chart is at most $\varepsilon/(\kappa\nu^2)$, while $|w_D/w_{T_z}|\le\mu/(\kappa\nu)$. This supplies the required margins in these tests.

\paragraph{Correction to algebraic level sets.}

In a quantitative five-column chart, select five output rows with an invertible $5\times5$ minor and use the local upper-affine coordinate slice omitting the chart direction. The five-output map has first and second derivatives and derivative inverse bounded by $P^C$. For the current augmented state $\widehat J=([J],\ell)$, let $\Pi$ select these five augmented output coordinates and let $r$ vary in that coordinate slice about the identity. The quantitative inverse-function theorem gives a coordinate radius $P^{-C_1}$ and target radius $P^{-C_2}$. For a target $q_{\rm grid}$ in the latter ball, the correction is the unique solution in the former ball of
\[
\Pi(\widehat J\circ r)=q_{\rm grid},\qquad \widehat J\circ r=([J\circ r],\ell\circ r),\qquad \|r-\mathrm{id}\|\le P^{-C_1}.
\]
Choose a fixed output grid of mesh smaller than a fixed fraction of the target radius. Its bounded output range gives $P^C$ targets, independently of $d_1$. Each corrected state lies on a one-dimensional algebraic level set, with boundedly many analytic cells by fixed-format decomposition. Fixed ties assign each original label once. The small correction preserves the enlarged windows, stopping and pivot inequalities. Zero-dimensional cells receive endpoint pieces with bounded counts; the analogous six-output argument gives bounded count in the full-rank region.

\par\medskip\noindent\textit{Step 2. Prepare and count the pieces.}\par\nobreak\smallskip

On a one-dimensional cell, use the parameter of Definition~\ref{def:regular-coordinates}, for which the omitted component of $g^{-1}g'$ is one. The remaining body components are rational functions of the bounded augmented jet and the retained minor. The omitted component of a nonzero tangent cannot vanish, since the transverse minor is invertible. The resulting intrinsic parameter may be an integral rather than an algebraic function; intrinsic derivatives of algebraic functions on the cell are nevertheless algebraic. Include all exact formulas for the frame change used below, their denominators with the stated lower bounds, and their finitely many required intrinsic derivatives in a common monotonicity decomposition. Divide nonzero rates into signed dyadic amplitude bands, and take Whitney intervals in the intrinsic parameter. On an interval of length $H$ this gives
\begin{equation}\label{reg:prepared-derivatives}
\|\widehat J^{(r)}\|\lesssim_P H^{-r},\qquad |u^{(r)}|\lesssim_P \sup_I|u|\,H^{-r}
\end{equation}
for each prepared rate $u$, through the fixed orders needed for frame changes and samples. Apply the one-sided divided-difference estimate from the proof of Lemma~\ref{geo:preparation} on the fixed enlargement of each Whitney interval; the sign and monotonicity of the derivative compare its value at the chosen point with the value supplied by the stencil. Signed bands compare $\sup|u|$ with the local amplitude. Merge intervals shorter than one into the two endpoint neighborhoods of bounded parameter length; a bounded parameter interval uses one piece. Bounded body speed gives count $O_P(\Lambda)$ there. The remaining dyadic intervals have length at least one, so their number is logarithmic in the maximum parameter length.

The number of amplitude bands is logarithmic in the physical range. Indeed, the bounded physical output in a fixed triangular presentation bounds the upper-affine coordinates by $M\lesssim_P d_1^{-1}$. Their total variations are $O_P(M)$ by fixed semialgebraic complexity. Thus the $T_y$ parameter has total length at most
\[
\operatorname{TV}(e)+\sup|d|\operatorname{TV}(f)\lesssim_P M^2=:T_{\mathrm{cap}}.
\]
Other parameters satisfy the same or better bound. Stop the signed dyadic bands at $cT_{\mathrm{cap}}^{-4}$. The terminal band $|u|\le cT_{\mathrm{cap}}^{-4}$ includes zero; its actual rates remain in the exact equations. Discarding it in a comparison gives an error $O_P(T_{\mathrm{cap}}^3T_{\mathrm{cap}}^{-4})$ in the right relative affine change, by the step-three propagator. Relative changes use the homogeneous affine matrices of Definition~\ref{def:regular-coordinates}.

Bounded body speed gives count $O_P(\Lambda)$ per unit interval and $O_P(\Lambda H)$ over length $H$. Fixed ties, parameter-compatible cells, and Borel integration and inversion of the positive parameter make the assignment measurable before any original index is fixed for coefficient integration.

There are finitely many row choices, $P^C$ grid targets, boundedly many cells per target, and at most a power of $1+\log P+\log d_1^{-1}$ bands and parameter intervals. This bounds the total multiplicity and that at each dyadic $H$. Hence $\sum_I H_I^{-\eta}\le P^C(1+\log P+\log d_1^{-1})^C\sum_{n\ge0}2^{-n\eta}$ for every fixed $\eta>0$, as required for common incidence.

\par\medskip\noindent\textit{Step 3. Compare fitted and actual body rates.}\par\nobreak\smallskip

For the fitted coefficients in Definition~\ref{def:regular-coordinates}, the full augmented identity on the curve with normalized parameter is $\widehat J'=R_i+L_iV_{\rm other}$, and therefore
\[
V_{\rm other}=v_i+L_i^+\widehat J',\qquad |V_{\rm other}-v_i|\lesssim_PH^{-1}.
\]
The last bound follows from the quantitative pivot and \eqref{reg:prepared-derivatives}. The fixed inner and outer cutoff margins consequently transfer every small-coefficient and significant-rate test to the actual long-piece field when $H\ge P^C$. Pieces with smaller $H$ have count $O_P(\Lambda)$ and are placed in the term for bounded counts. The derivative terms from the frame change below are also $O_P(H^{-1})$ by the same preparation, while their exact accumulated effects remain in the displayed equations.

\par\medskip\noindent\textit{Step 4. Choose samples and bound their coefficient set.}\par\nobreak\smallskip

For the complete evaluation map of Definition~\ref{def:sampling}, the linear volume formula and Cauchy--Binet give coefficient volume at most $P^C\det(\mathcal A_h^*\mathcal A_h)^{-1/2}\le P^C\Delta^{-1}$. Separate bounded right changes compare the full Gram forms. For a common left affine map $g$, write $gh=(g\circ h_i)_i$. The coefficient pullback of Lemma~\ref{an:character} gives
\[
\mathcal A_{gh}=\mathcal A_h\mathcal L_g,\qquad \det(\mathcal A_{gh}^*\mathcal A_{gh})=|\det\mathcal L_g|^2\det(\mathcal A_h^*\mathcal A_h).
\]
Thus determinant-one changes preserve it, and bounded changes with bounded inverses cost a power of $P$. These are pointwise comparisons, including for phase-dependent tuples.

The fixed frame supplies the first sample. Its bounded sample block at the assigned curve point can be replaced by the identity in the Gram form. To see this replacement precisely, if the assigned state differs from the anchor by a bounded right factor $u$, its coefficient operator is $\mathcal L_uJ=[J\circ u]$, with both $\|\mathcal L_u\|$ and $\|\mathcal L_u^{-1}\|$ bounded by powers of $P$. Equivalently, $P^{-C}\|G\|^2\le\|\mathcal L_uG\|^2\le P^C\|G\|^2$ for every $G\in V_k$. Thus an $m$-sample determinant needs only $m-1$ nets for the additional samples. If each fixed net of Definition~\ref{def:sampling} has size at most $N$, the resulting union gives \eqref{reg:fixed-net-bound}. The choice of selected sample tuple may depend on the phase: round each tuple into the fixed nets by bounded right factors, discard fixed net tuples failing the required Gram lower bound, and take the union. No independence of the phase and its selected tuple is assumed.

\paragraph{Nets in the original coordinate slices.}

For later use, the upper-affine group law is
\[
(a,b,c,d,e,f)(A,B,C,D,E,F)=(a+A,b+B+aD,c+C+aE+bF,d+D,e+E+dF,f+F).
\]
Haar measure is the six-coordinate Lebesgue measure. Let $U$ be the bounded right neighborhood used for the sampling comparisons and $Q^+$ the fixed polynomial enlargement of a comparison box $Q$; right stability means $QU\subseteq Q^+$. Write $w_a,\ldots,w_f$ for the comparison widths after taking their stated maxima with one. Let $N_\xi$ be the covering number of a slice of $Q^+$ with coordinate $\xi$ fixed, by the bounded right slice patches used below. Then
\[
\begin{aligned}
N_\xi&\lesssim\prod_{\eta\in\{a,b,c,d,e,f\}\setminus\{\xi\}}w_\eta &&(\xi=a,d,f),\\
N_e&\lesssim(1+\sup|d|)\prod_{\eta\in\{a,b,c,d,e,f\}\setminus\{e\}}w_\eta,\\
N_b&\lesssim(1+\sup|a|)\prod_{\eta\in\{a,b,c,d,e,f\}\setminus\{b\}}w_\eta.
\end{aligned}
\]
At fixed $e$, a bounded right slice patch must impose $E=-dF$, so take $|F|\lesssim(1+|d|)^{-1}$. At fixed $b$, use $B=-aD$. Disjoint right slice patches prove these bounds. The comparison boxes below contain bounded right neighborhoods and the bounded comparison to the fixed anchor.

\paragraph{Choice of sample points.}

On any one-sided fixed fraction of a prepared interval, every significant normalized rate varies by a prescribed small amount. A homomorphic parameter coordinate, or a nonhomomorphic coordinate whose displayed exact derivative is bounded away from zero, can therefore be sampled at prescribed levels. The normalized sample nodes are close to a fixed progression. The strictly upper-triangular affine matrix flow has a finite ordered-exponential formula of degree three, with polynomial perturbation bounds on bounded normalized sets. More precisely, for an $n\times n$ model minor $B$ with row norms at most $M$ and $|\det B|\ge b$, row errors at most $\varepsilon\le1$ give
\[
 |\det B_{\rm act}-\det B|\le n\varepsilon(M+1)^{n-1}.
\]
Here $n\le N_k$, $M\le P^C$ and $b\ge P^{-C}$, so an inverse-polynomial tolerance preserves at least half the determinant. If some entries are unbounded, apply the exact graded transfer first and perturb only its square diagonal blocks. In the swapped chart use the complete propagator estimate below instead. Its sampling fraction is chosen as $P^{-C}$ before the remaining thresholds; shortening changes the fixed-degree determinant bounds only by powers of $P$. The number $m$ of samples is fixed before these choices.

Fix an assigned label and let its corrected curve point have intrinsic parameter $t_j$. One of the two directions from $t_j$ contains a segment of length at least $c_PH$ in the fixed enlargement of the prepared interval. Choose the sample spacing smaller than $c_PH/(2m)$ on that side. The first curve sample is the corrected anchor; use the bounded-block comparison above. The other samples can be prescribed at fixed original-coordinate levels relative to the unchanged anchor whenever the coordinate derivative has the lower bound stated in the corresponding chart. The bounded initial correction changes these levels by $O_P(1)$; for long pieces this changes their normalized locations by an arbitrarily small amount. Shorter pieces are already in the bounded term. The interpolation determinants below are nonzero at every distinct progression; shortening its spacing changes their positive lower bounds only by powers of the chosen conditioning constants. Thus the construction permits the eighteen samples used below.

\end{proof}

\subsection{Frame changes and scale ranges}\label{reg:gauges}\label{reg:geometric-regimes}
\par\medskip\noindent\textit{Step 1. Normalization for all charts.}\par\nobreak\smallskip

The variables $u,w,r$ are accumulated rates before taking maxima with one; capitals such as $U=\max(1,u)$ denote those maxima. The original rates still occur in their exact differential equations. A significant rate satisfies the quantitative threshold specified in its comparison; it has constant sign and is relatively constant on the prepared interval. Fixed diagonal reflections absorb signs without changing absolute determinants or covering exponents. Tables~\ref{reg:all-regimes} and~\ref{reg:all-costs} collect the case numbers, models, dilations, and one-sample covering costs justified below. Their prescribed sample counts concern the $k\ge14$ tail estimate; the geometric comparisons allow any fixed sample count as above.

Here is the normalization used in every comparison. Let $h=\vartheta_{\rm cmp} H$, where $0<\vartheta_{\rm cmp}<1$ is a conditioning parameter chosen in the order specified for each comparison below, and let $D_s=\operatorname{diag}(s_x,s_y,s_z)$. The normalized relative path is $D_s^{-1}g(t_0)^{-1}g(t_0+hs)D_s$. If its unnormalized body field is $\alpha A+\beta C+\xi T_x+\delta D+\eta T_y+\gamma T_z$, the six normalized coefficients are exactly
\begin{equation}\label{reg:normalized-rates}
\left(h\alpha\frac{s_y}{s_x},\ h\beta\frac{s_z}{s_x},\ \frac{h\xi}{s_x},\ h\delta\frac{s_z}{s_y},\ \frac{h\eta}{s_y},\ \frac{h\gamma}{s_z}\right).
\end{equation}
For a lower shear $\zeta x\partial_y$, its normalized coefficient is $h\zeta s_x/s_y$; diagonal rates acquire the factor $h$. In applying a displayed dilation, replace accumulated displacements by their local versions, such as $h\sup|\alpha|$ or $h^2\sup|\delta|$, and then take the stated maxima with one. Prepared derivative bounds make each significant normalized rate differ from its value at the reference point by $O_P(\vartheta_{\rm cmp})$. Nonzero reference rates can be normalized to the displayed models by bounded diagonal changes and signs. The model sample nodes lie in a fixed normalized interval. Powers of $\vartheta_{\rm cmp}$ in the physical dilation only alter conditioned constants, since $\vartheta_{\rm cmp}$ is fixed independently of $H$. For $L$, translations lie below the total-degree diagonal blocks; for $M_f$, the $C,T_x,D,T_y$ terms lower $\deg_x+\deg_y$. Lemma~\ref{reg:lem:transfer} removes these potentially unbounded entries exactly before any perturbation estimate. Moving shears are undone in the original coordinates, where their coefficient operators and inverses are bounded by $P^C$; no bound on their conjugates by $D_s$ is required. Formula~\eqref{reg:normalized-rates} gives the short-interval comparisons below.

\par\medskip\noindent\phantomsection\textit{Step 2. The chart omitting $D$ (cases 1--9).}\label{reg:d-channels}\par\nobreak\smallskip
Write a prepared curve as $g(a(t),b(t),c(t),t,e(t),f(t))$. Its body velocity is
\[
V=\alpha A+\beta C+\xi T_x+D+\beta_yT_y+\gamma T_z,
\]
where
\[
\alpha=a',\quad \beta=b'-a,\quad \xi=c'-ae'+(at-b)f',\quad \beta_y=e'-tf',\quad \gamma=f'.
\]
Put $\tau=\xi-\beta\beta_y$. The bounded instantaneous affine change $X=x-\beta y$, $Y=y$, $Z=z+\beta_y$ gives the field $(\alpha Y+\tau)\partial_X+Z\partial_Y+\gamma\partial_Z$. Multiplication by $Z$ is injective on the polynomial translation output, and its physical component contains $Be_Y$. The stopped inequality therefore gives
\[
P^{-C}\le P^C\bigl(\|V\widehat J\|+|\alpha|+|\tau|+|\gamma|\bigr).
\]
Apply the chosen positive inverse-power cutoffs to $|\alpha|$, then $|\gamma|$, then $|\tau|$, in that order; a rate is called significant when it exceeds its cutoff. If none is significant, $\|\widehat J'\|\ge P^{-C}$. Bounded semialgebraic variation of the augmented state bounds the length of these remaining intervals by $P^C$, and their count is bounded as required.

\paragraph{Original-coordinate equations and covering net.}
The equations are
\begin{equation}\label{reg:d-equations}
a'=\alpha,\quad d'=1,\quad b'=a+\beta,\quad f'=\gamma,\quad e'=d\gamma+\beta_y,\quad c'=b\gamma+a\beta_y+\xi.
\end{equation}
In a vector expression, $O_P(U_1,\ldots,U_n)$ means an error in coordinate $i$ bounded in absolute value by $P^CU_i$. Starting at the unchanged anchor, integration gives the original-coordinate box
\begin{equation}\label{reg:d-box}
(a,b,c,d,e,f)=O_P(U,HU,HUW,H,HW,W),\qquad N_d\lesssim_P H^3U^3W^3,
\end{equation}
where the last expression is the fixed-$d$ sampling covering number. The box contains all bounded right neighborhoods and the bounded comparison to the actual anchor: this follows term by term from the group multiplication formula in the preceding proof.

\paragraph{First-shear and terminal-rate comparisons.}
When $\alpha$ is significant, $U\sim_PH$ and $W=\max(1,H\sup|\gamma|)$. For bounded $W$, the homogeneous linear chain $A+D$ gives $L$; sample translations lie below its homogeneous diagonal blocks. For large $W$, relative freezing of the two simple rates gives the principal affine chain $A+D+T_z$ and $Q$. The remaining normalized rates tend to zero with the chosen short fraction or the reciprocal large-displacement threshold.

When $\gamma$ is significant, $W\sim_PH$ and $U=\max(1,H\sup|\alpha|)$. For bounded $U$, the integrated $A$ error is small on a short fraction, and a bounded frozen $x$--$y$ shear removes the leading $C$ term, giving $P$. For large $U$, relative preparation gives the principal chain $Q$. The other normalized rates are small, and the original sampling box remains~\eqref{reg:d-box}.

\paragraph{The significant mixed rate.}
For significant $\tau$, a time-dependent change is necessary. Write $G(t)$ for the homogeneous-coordinate matrix of the affine path $g(t)$, and represent affine changes by the same matrix convention. Let
\[
S(t)(x,y,z)=(x-\beta(t)y,y,z+\beta_y(t)),\qquad \widetilde G=GS^{-1}.
\]
Direct matrix differentiation gives
\begin{equation}\label{reg:d-gauge}
\widetilde G^{-1}\widetilde G'=SVS^{-1}-S'S^{-1}=(\alpha+\beta')A+D+\tau T_x+(\gamma-\beta_y')T_z.
\end{equation}
Set $u=H\sup|\alpha+\beta'|$, $w=H\sup|\gamma-\beta_y'|$, $U=\max(1,u)$, and $W=\max(1,w)$. The widths after the time-dependent change of frame $S$ are $(u,Hu,H(1+uw),H,Hw,w)$.

At an anchor $t_0$, write $\beta_0=\beta(t_0)$. Multiplication back to the original coordinates gives exactly
\[
\begin{aligned}
a_o&=a+\beta_0-\beta(t),&b_o&=b+\beta_0d,&d_o&=d,\\
c_o&=c+\beta_0e+(b+\beta_0d)\beta_y(t),&e_o&=e+d\beta_y(t),&f_o&=f+\beta_y(t)-\beta_y(t_0).
\end{aligned}
\]
The additional term $\beta_0e$ has size $O_P(HW)$ and is absorbed in the $c_o$ width $HUW$, since $U\ge1$. Consequently \eqref{reg:d-box} remains the original box, also when $u<1$. Four comparisons exhaust this part:
\[
\begin{array}{c|c|c}
\text{condition}&(s_x,s_y,s_z)&\text{model}\\\hline
u,w\lesssim_P1&(H,H,1)&M\\
u\gg_P1,\ w\lesssim_P1&(HU,H,1)&L\\
w\gg_P1,\ uw\lesssim_P1&(H,Hw,w)&R\\
w,uw\gg_P1&(Huw,Hw,w)&Q.
\end{array}
\]
In the second line translations lie below the homogeneous diagonal blocks. In the third line use local scales $(\vartheta_{\rm cmp} H,\vartheta_{\rm cmp}^2Hw,\vartheta_{\rm cmp} w)$. Formula~\eqref{reg:normalized-rates} makes the $A$ perturbation $O_P(\vartheta_{\rm cmp}^2uw)$, so a fixed short fraction covers every fixed bounded-product range. In the fourth line the local scales have factors $(\vartheta_{\rm cmp}^3,\vartheta_{\rm cmp}^2,\vartheta_{\rm cmp})$, and the residual $T_x$ coefficient is $O_P((\vartheta_{\rm cmp}^2uw)^{-1})$; after fixing $\vartheta_{\rm cmp}$, the large-product threshold controls it. Up to that conditioned factor this is $O_P((uw)^{-1})$. The two nonsimple rates are removed exactly in \eqref{reg:d-gauge}. The comparison therefore remains valid when $u<1$ and $uw$ is large. Choose the cubic fraction and threshold first, then a shorter parabolic fraction for its bounded complement, then the linear and bounded fractions. Their fixed-priority overlapping ranges contain every pair $(u,w)$.

\par\medskip\noindent\textit{Step 3. The residual chart omitting $T_z$ (cases 10--14).}\par\nobreak\smallskip
The initial field is $\alpha A+\beta C+\tau T_x+\delta D+\eta T_y+T_z$. The bounded change of frame $S(x,y,z)=(x-\tau z,y-\eta z,z)$ gives
\[
\widetilde V=\alpha A+\zeta C+\widetilde\delta D+T_z,\qquad \zeta=\beta+\alpha\eta+\tau',\quad \widetilde\delta=\delta+\eta'.
\]
We relabel $\widetilde\delta$ as $\delta$. The least-squares cutoff and preparation make $\delta$ small. Translation nondegeneracy and the small augmented derivative force $|\alpha|+|\zeta|\ge P^{-C}$. The equations are
\[
a'=\alpha,\quad d'=\delta,\quad b'=a\delta+\zeta,\quad e'=d,\quad c'=b,\quad f'=1.
\]
\paragraph{Original coordinates and covering net.}
Let $u=H^2\sup|\delta|$, $v=H\sup|\zeta|$, and $B_0=\max(H,u)$. In particular $u\lesssim_PH^2$, even though the rate $\delta$ is small. The original-coordinate bounds, including both bounded comparisons, are
\begin{equation}\label{reg:f-box}
(a,b,c,d,e,f)=O_P(H,B_0,HB_0,\max(1,u/H),B_0,H),\qquad N_f\lesssim_PH^2B_0^3\max(1,u/H).
\end{equation}
For example, a bounded left comparison adds multiples of $d$ to $b$ and of $e,f$ to $c$, while a bounded right comparison can add $O(H)$ to $b$. These contributions require $B_0$, rather than the width after the change of frame $S$ alone.

\paragraph{Model comparisons.}
For significant $\alpha$, the $M_f$ comparison removes the other four coordinates by its preserved grade, without a smallness assumption. The two nets follow from~\eqref{reg:f-box}.

A spectator variable means the coordinate unchanged by the two-variable model flow. Order rows and columns by its exponent. Each transfer below specifies an ordering with matching input and output multiplicities; Lemma~\ref{reg:lem:transfer} then identifies the determinant with the product of the square diagonal-block determinants.
For significant $\zeta$, one has $v\sim_PH$. When $u\lesssim_PH$, the normalization $(Hv,\max(1,u),H)$ gives $P_f$; the normalized $A$ perturbation is small, and at $A=0$ the $d,e$ entries occur in strictly off-diagonal blocks when rows and columns are ordered by the exponent of the unchanged coordinate. The net follows from~\eqref{reg:f-box}. If $u\gtrsim_PH$ and $|\alpha|u\lesssim_Pv$, use $(Hv,u,H)$ and the same model. The five fixed-$f$ widths are $(Hv/u,v,Hv,u/H,u)$, their product is $Hv^3u\lesssim_PHv^2u^2$, since $v\lesssim_Pu$, so this last expression bounds the net; the first width is at least a fixed inverse power of $P$ and absorbs bounded comparisons. In the remaining range $|\alpha|u\gg_Pv$, let $r=H\sup|\alpha|$. Then $ru\gtrsim_PH^2$, $r\lesssim_PH$, and $(ru,u,H)$ gives $Q$, with a small normalized $C$ perturbation. The original fixed-$f$ widths are $(r,ru/H,ru,u/H,u)$, giving the net $r^3u^4/H^2$. A shorter parabolic fraction covers the bounded complement of the cubic dominance threshold, so there is no omitted transition range.

\par\medskip\noindent\textit{Step 4. The significant first-shear part of the $T_y$ chart (cases 15--16).}\par\nobreak\smallskip
Use the chart omitting $A$, with field $A+\beta C+\tau T_x+\delta D+\eta T_y+\gamma T_z$. The bounded change of frame $S(x,y,z)=(x,y+\beta z+\tau,z)$ gives
\[
\widetilde V=A+(\delta-\beta')D+(\eta+\beta\gamma-\tau')T_y+\gamma T_z.
\]
On long pieces, $|\eta+\beta\gamma-\tau'|\ge P^{-C}$. Otherwise the residual $D,T_z$ cutoffs and the prepared derivatives would leave $A$ as the only significant field, contradicting the stopped lower bound for $y\partial_x$ and the small augmented derivative. Set $u=H\sup|\delta-\beta'|$, $w=H\sup|\gamma|$, $U=\max(1,u)$, $W=\max(1,w)$. Integration gives widths after the change of frame $S$
\[
(a,d,f,b,e,c)=O_P(H,u,w,Hu,H+uw,H^2+Huw).
\]
For $UW\lesssim_PH$, use $P_a$; translations in the unchanged coordinate occur in strictly off-diagonal blocks. For $uw\gg_PH$, use $Q$. The dilations and original fixed-$a$ nets follow from the displayed widths and are listed in the case tables. In the second range both rates before taking maxima with one are bounded below by a fixed inverse power of $P$, since each is at most $P^CH$. All boxes absorb the original bounded comparisons, and the short bounded-product comparison covers each fixed transition ratio.

\par\medskip\noindent\textit{Step 5. The central chart and coordinate interchange (cases 17--20).}\par\nobreak\smallskip
Use the exact $T_y$ parameter, whose differential is $de-d\,df$, not $de$. For $V=\alpha A+\beta C+\tau T_x+\delta D+T_y+\gamma T_z$, the bounded change of frame $S(x,y,z)=(x-\tau y,y,z)$ changes the first two rates to $\alpha+\tau'$ and $\beta-\tau\delta$. The prepared derivative is $O_P(H^{-1})$, and the residual $D$ cutoff retains the central lower bound. Relabeling gives
\[
V=\alpha A+\beta C+\delta D+T_y+\gamma T_z,\qquad |\beta|\sim_P1,\qquad |\alpha|+|\delta|+|\gamma|\le\epsilon_0.
\]
Put
\[
\begin{gathered}
U=\max(1,H\sup|\alpha|),\qquad V_0=\max(1,H\sup|\delta|),\\
W=\max(1,H\sup|\gamma|),\qquad M_0=\max(U,W).
\end{gathered}
\]
The equations give
\begin{equation}\label{reg:central-equations}
a'=\alpha,\quad d'=\delta,\quad f'=\gamma,\quad b'=a\delta+\beta,\quad e'=d\gamma+1,\quad c'=b\gamma+a,
\end{equation}
\begin{equation}\label{reg:central-box}
(a,b,c,d,e,f)=O_P(U,H+UV_0,H(U+W)+UV_0W,V_0,H+V_0W,W).
\end{equation}
The central change of frame also preserves the displayed original box: with $\tau_0=\tau(t_0)$, the relative original coordinates satisfy $b_o=b+\tau_0d$ and $c_o=c+\tau_0e$. The extra terms have sizes $O_P(V_0)$ and $O_P(H+V_0W)$, respectively, and are absorbed by \eqref{reg:central-box}.

\paragraph{Product regimes and model comparisons.}
The two relevant ratios are $UV_0/H$ and $V_0W/H$. Choose the simultaneous cubic comparison and its dominance threshold first. Choose the two single-dominance comparisons and their thresholds next, using their fixed complementary product bounds. Finally choose the bounded-product fraction after all resulting finite transition bounds. The estimates below justify this order; in the comparison after coordinate interchange use the separate $Q/R$ order following \eqref{reg:swapped-error}. Fixed-priority overlaps cover: both ratios bounded; only the first large; both large; and only the second large.

If both products are bounded, $(HM_0,H,M_0)$ gives $M_e$. At $D=0$, the additional entries $a,c,f$ are removed by grading. On a short fraction, the normalized $D$ coefficient is at most $P^C(\epsilon_0+\vartheta_{\rm cmp}\max(UV_0/H,V_0W/H))$, using the local value of $M_0$ after taking the maximum with one. It is small after the residual cutoff and fraction are chosen in that order. Also $e'=1+d\gamma$ differs from $1$ by at most $P^C(\epsilon_0+\vartheta_{\rm cmp} V_0W/H)$, so it is bounded away from zero. The five-dimensional box on an $e$ level has volume $H^2UV_0WM_0$. The nonhomomorphic slice contributes $V_0$, so
\[
N_e\lesssim_PH^2UV_0^2WM_0.
\]
Split $U\ge W$ and $W\ge U$. The auxiliary scale $M_0$ cancels in the complete minor.

If $UV_0/H$ is large and $V_0W/H$ is bounded, use $(UH,H,H/V_0)$. The model field is $A+D+T_y$, hence $L_e$; the central coefficient is $O_P(H/(\vartheta_{\rm cmp} UV_0))$ and the terminal coefficient is $O_P(\vartheta_{\rm cmp} V_0W/H)$ under the local dilation $(\vartheta_{\rm cmp}^2UH,\vartheta_{\rm cmp} H,H/V_0)$. Thus they are small by the ordered threshold choice. On actual $d$ levels, the five fixed-$d$ widths from \eqref{reg:central-box} are $(U,UV_0,HU,H,W)$, up to conditioning factors, because $V_0W\lesssim_PH$ and $W\lesssim_PU$. Their product is $H^2U^3V_0W\lesssim_PH^3U^3$. Thus the original net is at most $H^3U^3$. The $T_y$ term remains part of $L_e$.

If both products are large, use $(UV_0W,V_0W,W)$ and $Q$. Under the local dilation with factors $(\vartheta_{\rm cmp}^3,\vartheta_{\rm cmp}^2,\vartheta_{\rm cmp})$, the central and middle-translation coefficients are $O_P(H/(\vartheta_{\rm cmp} UV_0))$ and $O_P(H/(\vartheta_{\rm cmp} V_0W))$. Fix the fraction before choosing their dominance thresholds. The original fixed-$d$ net is $U^3V_0^3W^3$. Here $U,W\gtrsim_PH/V_0$ and $V_0\lesssim_PH$, which are retained in the logarithmic region.

\par\medskip\noindent\textit{Step 6. Coordinate interchange and the remaining two cases (cases 21--22).}\par\nobreak\smallskip

Suppose finally that $V_0W/H$ is large while $UV_0/H$ is bounded. Keep the original box $(U,H,HW,V_0,V_0W,W)$ and prescribe original $b$ levels. In the coordinates after the change of frame $S$, $b'=a\delta+\beta$; the original relative coordinate instead has $b_o'=a\delta+\beta+\tau_0\delta$. Choose the residual $D$ cutoff after the central lower bound so that $|\tau_0\delta|$ is smaller than its retained margin. The short-fraction choice also controls $a\delta$. Consequently $b_o'$ is quantitatively nonzero, which justifies sampling at original coordinate levels. Exchange $x,y$, let $r=\delta/\beta$, and apply $S(x,y,z)=(x-r y,y,z)$. Direct differentiation gives the full field
\begin{equation}\label{reg:swapped-field}
\widetilde V=\beta D+(r'-r^2\alpha)A+T_x+\gamma T_z+\alpha x\partial_y+r\alpha(-x\partial_x+y\partial_y).
\end{equation}
Reparameterize by $s$ with $ds/dt=\beta(t)$. By the retained bound $|\beta|\sim_P1$, this is invertible on the selected interval. The body field with respect to $s$ is $\widetilde V/\beta$, with the signed value of $\beta$. Set
\[
u_*=H\sup|(r'-r^2\alpha)/\beta|,\qquad X=\max(1,u_*W).
\]
Preparation of the expression gives
\[
u_*\lesssim_P\frac{V_0}{H}\left(1+\frac{UV_0}{H}\right).
\]
Indeed $r=O_P(V_0/H)$, $r'=O_P(V_0/H^2)$, and $\alpha=O_P(U/H)$, which give the two displayed terms after multiplying by $H$. For the local comparison put $\rho=UV_0/H$, $\sigma=V_0W/H$, $x=u_*W$, and $X_{\vartheta_{\rm cmp}}=\max(1,\vartheta_{\rm cmp}^2x)$. Use the local dilation
\[
(s_x,s_y,s_z)=(\vartheta_{\rm cmp} H X_{\vartheta_{\rm cmp}},\vartheta_{\rm cmp}^2HW,\vartheta_{\rm cmp} W).
\]
By \eqref{reg:normalized-rates}, the integrated lower-shear coefficient is at most
\[
P^C\frac{UX_{\vartheta_{\rm cmp}}}{W}\le P^C\left(\frac{U}{W}+\vartheta_{\rm cmp}^2Uu_*\right)\le P^C\left(\frac\rho\sigma+\vartheta_{\rm cmp}^2\rho(1+\rho)\right).
\]
The integrated diagonal coefficient is at most $P^C\vartheta_{\rm cmp}\rho$. Thus the sum of these errors is bounded by
\begin{equation}\label{reg:swapped-error}
P^C\left(\frac\rho\sigma+\vartheta_{\rm cmp}\rho+\vartheta_{\rm cmp}^2\rho(1+\rho)\right).
\end{equation}
The bound for $\rho$ is fixed on this branch. Choose the $\sigma$ dominance threshold so that the first term is small, and then the fraction so that the other two terms are small. The finite transition range in $\sigma$ belongs to the shorter bounded-product comparison. Significant-rate freezing contributes the additional error $O_P(\vartheta_{\rm cmp})$.

\paragraph{Comparison fractions and the complete propagator.}
Choose the comparison parameters in this order. First fix a cubic fraction $\vartheta_Q$ small enough for~\eqref{reg:swapped-error}, then $L$ large enough that $(\vartheta_Q^2L)^{-1}$ is small. For $x\ge L$, the $A,D,T_z$ model is $Q$ and its extra $T_x$ coefficient is at most $P^C(\vartheta_Q^2x)^{-1}$. For $x\le L$, choose $\vartheta_R$ smaller so $\vartheta_R^2L$ is small. Then $X_{\vartheta_R}=1$, the $D,T_x,T_z$ model is $R$, and the extra $A$ coefficient is at most $P^C\vartheta_R^2x$. The comparisons overlap at fixed thresholds. Their local factors depend only on conditioning and degree, preserving the powers of $H,U,V_0,W$.

Let $G_0$ and $G$ be the normalized model and complete propagators. Write $G'=G(V_{\mathrm{core}}+E)$ and $G_0'=G_0V_{\mathrm{core}}$, including in $E$ the lower shear, diagonal terms, rate-freezing errors and remaining core perturbations. With $G=ZG_0$, one has $Z'=Z(G_0EG_0^{-1})$, $Z(0)=I$. Since $\|G_0\|+\|G_0^{-1}\|\le P^C$,
\[
 \|Z-I\|\le\exp\left(P^C\int\|E\|\right)-1
 \le 2P^C\int\|E\|\qquad\left(P^C\int\|E\|\le1\right).
\]
Choose the preceding errors below the sampling-minor tolerance divided by the further polynomial factors in this estimate and in coefficient evaluation. This preserves the complete minor without an $H$ or physical-aspect loss.

\paragraph{Models and the original covering net.}
For bounded $u_*W$ the model is $R$, with $(H,HW,W)$; for large $u_*W$ it is $Q$, with $(HX,HW,W)$. The fixed-$b$ five-dimensional volume is $HUV_0^2W^3$, and the right-relative slice contributes the additional factor $U$. Thus each original sampling net has cardinality at most $HU^2V_0^2W^3\lesssim_PH^3W^3$. The permutation and bounded change of frame compare Gram determinants pointwise.

The preceding arguments use the same finite set of geometric comparisons for every fixed degree. Increasing the sample count changes only the predetermined comparison fractions and conditioning constants. Every perturbed complete minor is compared after all normalized entries have been bounded. Additional entries which may remain large are handled only through the exact square diagonal-block identities proved below.

\subsection{Sampling matrices, weights, and graded transfers}\label{reg:sampling}
The comparisons above use the following normalized affine flows:
\begin{align*}
L_t&=(x+ty+t^2z/2,y+tz,z),& Q_t&=(x+ty+t^2z/2+t^3/6,y+tz+t^2/2,z+t),\\
P_t&=(x,y+tz+t^2/2,z+t),& M_t&=(x+t,y+tz,z),\\
R_t&=(x+t,y+tz+t^2/2,z+t),& M_{f,t}&=(x+ty,y,z+t),\\
P_{f,t}&=(x+tz+t^2/2,y,z+t),& P_{a,t}&=(x+2ty+t^2,y+t,z),\\
M_{e,t}&=(x+tz,y+t,z),& L_{e,t}&=(x+ty+t^2(z+1)/2,y+t(z+1),z).
\end{align*}
Samples may be taken at any fixed distinct nodes. For a fixed integer $M\ge2$, we use $t=0,\ldots,M-1$; shrinking the progression into a fixed comparison interval only changes conditioned constants. Input columns and scalar output rows are indexed by all nonconstant monomials of degree at most $k$. In input column $\mu$, the row at sample $t$ and output exponent $\nu$ is
\begin{equation}\label{reg:sample-entry}
\mathcal M_{(t,\nu),\mu}=[x^{\nu_x}y^{\nu_y}z^{\nu_z}](x+a_ty+b_tz+c_t)^{\mu_x}(y+d_tz+e_t)^{\mu_y}(z+f_t)^{\mu_z}.
\end{equation}
\equationalias{low:4:rd:matrix_entry}\equationalias{low:5:rd5:matrix_entry}\equationalias{low:78:h3r:entry}\equationalias{low:finite-entry}

\paragraph{Sampling weights and coefficient-volume exponents.}
For a selected complete minor, let $R_x,R_y,R_z$ be its output-exponent sums. Diagonal conjugation by $\operatorname{diag}(s_x,s_y,s_z)$ multiplies its unit-scale determinant by
\begin{equation}\label{reg:minor-scaling}
s_x^{\chi_k-R_x}s_y^{\chi_k-R_y}s_z^{\chi_k-R_z}.
\end{equation}
\equationalias{low:4:rd:scaling}\equationalias{low:56:weight}\equationalias{low:78:h3r:character}
Indeed, the $(\nu,\mu)$ entry is multiplied by $s^{\mu-\nu}$, and each input-coordinate exponent sums to $\chi_k$. Factoring columns and rows proves the identity. For $(s_x,s_y,s_z)=(H^a,H^b,H^c)$, the exponent of $H$ in the determinant is
\begin{equation}\label{reg:eq:character}
D=(\chi_k-R_x)a+(\chi_k-R_y)b+(\chi_k-R_z)c.
\end{equation}
The inverse determinant bound $P^CH^{-D}$ and the $M-1$ original-coordinate sample nets, each of size at most $P^CH^n$, give coefficient-volume decay exponent
\begin{equation}\label{reg:eq:witness}
D-(M-1)n.
\end{equation}
This follows from~\eqref{reg:fixed-net-bound}, including its union over all permitted fixed-net tuples.

We will use four exact transfers. First, total degree makes the full $Q$ matrix triangular, with homogeneous diagonal blocks equal to those of $L$. Second, the $x$ exponent makes $Q$ triangular, with diagonal blocks equal to $P$ on $(y,z)$. Third, the $z$ exponent makes $L_e$ triangular in the opposite direction, since its substitutions only add nonnegative powers of $z$; its diagonal blocks are the parabolic flow on $(x,y)$ obtained by setting $z=0$. Fourth, the grade $\deg_y+\deg_z$ makes $R$ triangular, with diagonal blocks equal to $M$. Indeed the translations in $y,z$ lower that grade, while $y\mapsto y+tz$ and $x\mapsto x+t$ preserve it. 

\subsection{Conic, cubic and mixed interpolation}\label{reg:interpolation}
For $M\ge2$, define
\begin{equation}\label{reg:eq:ABT}
 \ell_{n,j}=2(n-2j)+1,\qquad
 \mathsf A_k=\sum_{n=0}^k\sum_{j=0}^{\lfloor n/2\rfloor}j\ell_{n,j},\qquad
 \mathsf B_{k,M}=\sum_{n=0}^k\sum_{j=0}^{\lfloor n/2\rfloor}\sum_{i=0}^{\ell_{n,j}-1}\lfloor i/M\rfloor,
\end{equation}
\[
 \mathsf T_k=\sum_{n=0}^k\sum_{j=0}^{\lfloor n/2\rfloor}\frac{\ell_{n,j}(\ell_{n,j}-1)}2.
\]
The identities and bounds needed below are
\begin{equation}\label{reg:eq:ABTbounds}
 2\mathsf A_k+\mathsf T_k=3\chi_k,\qquad 0\le\mathsf A_k\le\frac12\chi_k,\qquad \mathsf B_{k,M}\le\frac{\mathsf T_k}{M}.
\end{equation}

\begin{lemma}[The conic sums]\label{reg:lem:sums}
The formulas \eqref{reg:eq:ABTbounds} hold for every $k\ge0$.
\end{lemma}
\begin{proof}
For weights $(2,1,0)$, the degree-$n$ character identity is
\[
\sum_{\alpha+\beta+\gamma=n}t^{2\alpha+\beta}=\sum_{j=0}^{\lfloor n/2\rfloor}\sum_{i=0}^{2(n-2j)}t^{2j+i},\qquad \alpha,\beta,\gamma\in\mathbb Z_{\ge0}.
\]
To verify this directly, the multiplicity of weight $r\le n$ in the monomial basis is $\lfloor r/2\rfloor+1$, and the upper half is its reflection. Summing the characters for $n\le k$ gives $2\mathsf A_k+\mathsf T_k=3\chi_k$. The last bound in \eqref{reg:eq:ABTbounds} follows termwise from $\lfloor i/M\rfloor\le i/M$.

For the remaining inequality, elementary geometric-series summation gives
\[
 \sum_{k\ge0}\mathsf A_k t^k=\frac{t^2}{(1-t)^5(1+t)},\qquad
 \sum_{k\ge0}\left(\frac12\chi_k-\mathsf A_k\right)t^k=\frac{t}{2(1-t)^3(1-t^2)}.
\]
For the first identity, write $n=2j+r$ in \eqref{reg:eq:ABT}. The second identity uses $\sum_k\chi_kt^k=t/(1-t)^5$. The final rational function has nonnegative coefficients. This proves the desired inequality in every degree.
\end{proof}

\paragraph{Balanced Hermite selection.}\label{reg:balanced-hermite}
For polynomials of degree at most $\ell-1$ at $M$ fixed distinct nodes, write $\ell=Mq+r$, $0\le r<M$. Select derivative orders $0,\ldots,q-1$ at every node and order $q$ at the first $r$ nodes. These $\ell$ functionals are independent by Hermite interpolation \cite[Proposition 7]{deBoor}, and their derivative-order sum is $\sum_{i=0}^{\ell-1}\lfloor i/M\rfloor$. The applications below identify the quotient or chain with this polynomial space and check the availability and weights of the selected rows separately.

\paragraph{Homogeneous conics.}
Let $\mathcal H_n$ be the homogeneous degree-$n$ space in $(x,y,z)$, with $\mathcal H_n=\{0\}$ for $n<0$, and let $q_h=xz-y^2/2$. Its filtration by powers of $q_h$ has successive quotient dimensions
\[
 \dim \mathcal H_{n-2j}-\dim\mathcal H_{n-2j-2}=2(n-2j)+1=\ell_{n,j}.
\]
Restriction to $(x,y,z)=(t^2/2,t,1)$ identifies the degree-$d$ quotient with all univariate polynomials of degree at most $2d$. Surjectivity follows because every power $t^r$, $0\le r\le2d$, is the restriction of a multiple of $x^{\lfloor r/2\rfloor}y^{r-2\lfloor r/2\rfloor}z^{d-\lceil r/2\rceil}$. The kernel is $(q_h)\cap\mathcal H_d=q_h\mathcal H_{d-2}$: at $z=1$, polynomial division by $x-y^2/2$ gives zero remainder. Homogeneity then gives divisibility by $q_h$ after localizing at $z$. The polynomial $xz-y^2/2$ is irreducible, since its coefficients as a linear polynomial in $x$ are relatively prime in $\mathbb R[y,z]$; it is therefore prime in the polynomial ring and does not divide $z$. Divisibility in the localization consequently implies divisibility in the polynomial ring. Degree additivity gives the stated quotient degree.

Let $F_n\in\mathcal H_n$ be the homogeneous degree-$n$ part of $F$. At a sample node $t$, form the row combination
\begin{equation}\label{reg:eq:virtual-h}
 [u^jv^i]F_n\bigl(u+(v+t)^2/2,v+t,1\bigr).
\end{equation}
This is a linear combination of sampled homogeneous coefficient rows of $F_n\circ L_t$. More precisely, its support consists of
\begin{equation}\label{reg:eq:rows-h}
 (x,y,z)\text{ exponents }(j+p,i-2p,n-j-i+p),\qquad 0\le p\le\lfloor i/2\rfloor,
\end{equation}
with coefficient $\binom{j+p}{j}2^{-p}$ when the exponents are nonnegative. On the $j$th filtration quotient, these rows are exactly the order-$i$ derivatives at $t$ of a polynomial of degree at most $\ell_{n,j}-1$, up to fixed nonzero factors. Rows of filtration order $j$ annihilate every higher filtration level.

Apply the \hyperref[reg:balanced-hermite]{balanced Hermite selection} to each quotient. All quotient blocks, and hence the complete matrix of these linear combinations, are nonsingular. Because $M\ge2$, the largest selected derivative order is at most $n-2j$, so $j+i\le n$. Thus all evaluation rows in \eqref{reg:eq:rows-h} are available in the original degree range.

\begin{proposition}[Homogeneous conic sampling]\label{reg:prop:L}
If $a\ge c$ and $b\ge c$, put $A=a-c$, $B=b-c$, and $d=\max(B,A/2)$.

\noindent\textup{(i)} There is a complete minor for $L$, preserving every homogeneous-degree multiplicity, with determinant exponent $D$ in \eqref{reg:eq:character} at least
\begin{equation}\label{reg:eq:Lgain}
 \mathcal L_k(a,b,c)=(A+B)\chi_k-A\mathsf A_k-d\mathsf B_{k,M}.
\end{equation}

\noindent\textup{(ii)} The same bound holds for the full affine model $Q$.
\end{proposition}
\begin{proof}
A row in \eqref{reg:eq:rows-h} has weight
\[
nc+Aj+Bi+p(A-2B)\le nc+Aj+di.
\]
Its fixed linear-combination row bound is therefore $nc+Aj+di$. The complete input weight is $(a+b+c)\chi_k$, and $\sum_{n=1}^k n\dim\mathcal H_n=3\chi_k$. Summing the bounds for all these linear-combination rows and applying Lemma~\ref{reg:lem:transfer} gives \eqref{reg:eq:Lgain}. The construction preserves output degree in each block, so the total-degree block-triangular argument from $L$ to $Q$ is exact.
\end{proof}

\paragraph{Affine parabolas and the constant quotient.}
For a two-variable polynomial of degree at most $n$, use $q_p=Y-Z^2/2$. Its $q_p$-power filtration has the same quotient dimensions $\ell_{n,j}$. Restriction to $(Y,Z)=(t^2/2,t)$ is onto the polynomials of degree at most $2n$, and its kernel is $q_p\R[Y,Z]_{\le n-2}$ by division by the monic polynomial $Y-Z^2/2$. The same argument on higher powers gives the whole filtration.

The row combination $[u^jv^i]G(u+(v+t)^2/2,v+t)$ uses physical output exponents $(j+p,i-2p)$ with coefficients $\binom{j+p}{j}2^{-p}$. If the two coordinate weights are $(a,b)$ and $d=\max(b,a/2)$, its evaluation rows have weights at most $aj+di$.

Apply this construction in every fixed block with fixed exponent of the remaining variable: for the exponent $s$, the remaining degree is $n=k-s$. In the block $s=0,j=0$, the constant coefficient is absent at every sample because the original space is modulo constants. Replace the $2k+1$ value/derivative functionals by $2k$ functionals of orders
\begin{equation}\label{reg:eq:constant-correction}
 1+\lfloor i/M\rfloor,\qquad i=0,\ldots,2k-1,
\end{equation}
using the \hyperref[reg:balanced-hermite]{balanced Hermite selection} on the derivative of the restriction polynomial, of degree at most $2k-1$. They therefore determine the restriction polynomial modulo constants. Their kernel leaves exactly the higher $q_p$ filtration, modulo the discarded constant. Every other block is unchanged. The derivative-order sum increases, relative to the full construction for the chain containing the constant, by
\[
 2k-\lfloor 2k/M\rfloor\le2k.
\]

\begin{proposition}[Affine conic sampling]\label{reg:prop:P}

\leavevmode\noindent\textup{(i)} For weights $(a,b)\ge0$ on the normal and tangential parabolic coordinates, and an arbitrary spectator weight, there is a complete minor preserving spectator exponents and having determinant exponent $D$ in \eqref{reg:eq:character} at least
\begin{equation}\label{reg:eq:Pgain}
 \mathcal P_k(a,b)=(a+b)\chi_k-a\mathsf A_k-\max(b,a/2)(\mathsf B_{k,M}+2k).
\end{equation}
It applies to $P$, $P_f$, and $P_a$ with the appropriate coordinate permutation.

\noindent\textup{(ii)} It also applies to $Q$ with $(a,b)$ equal to its $(y,z)$ weights, and to $L_e$ with $(a,b)$ equal to its $(x,y)$ weights.
\end{proposition}
\begin{proof}
The filtration and Hermite argument prove that the complete matrix of linear-combination rows is nonsingular, including the modified constant-free block. The spectator exponent has exactly the input multiplicities, so its character cancels. Summing the normal and tangential row-weight bounds and applying Lemma~\ref{reg:lem:transfer} gives \eqref{reg:eq:Pgain}. The factor $2$ in $P_a$ changes the parabola from $Y=Z^2/2$ to $Y=Z^2$ and changes fixed coefficients, not row supports or exponent bounds. The transfers to $Q$ and $L_e$ use that lemma with the $x$- and $z$-grade triangular identities established above.
\end{proof}

\paragraph{Graded Jordan chains}
Homogenize the degree-at-most-$k$ space to $\Sym^k\R^4$ in $(x,y,z,w)$. The generator of $Q_t$ is
\[
N=y\partial_x+z\partial_y+w\partial_z.
\]
Give the variables natural weights $(3,2,1,0)$ and put
\[
U=3x\partial_y+4y\partial_z+3z\partial_w,\qquad H=[U,N]=3x\partial_x+y\partial_y-z\partial_z-3w\partial_w.
\]
Direct calculation gives $[H,U]=2U$ and $[H,N]=-2N$. On $\operatorname{Sym}^k\mathbb R^4$, the natural-weight operator is $(H+3kI)/2$. After complexification, the classification and complete reducibility of finite-dimensional $\mathfrak{sl}_2$ modules \cite[Theorem 11.16]{EtingofSL2} decompose this space into graded chains. A lowest vector $Nv=0$, $Hv=-\lambda v$ has $\lambda\in\mathbb Z_{\ge0}$ and chain $v,Uv,\ldots,U^\lambda v$, with
\begin{equation}\label{reg:eq:sl2}
NU^rv=r(\lambda-r+1)U^{r-1}v.
\end{equation}
The matrices are rational. Choose a rational basis of each $\ker N\cap\ker(H+\lambda I)$ and apply $U$ repeatedly. The classification shows that the resulting chains form a basis after complexification, hence already over $\mathbb R$. This also proves that the chain changes preserve the natural grades and may be taken rational.

The degree-$k$, highest-weight-$\lambda$ chain multiplicities are the coefficients of $t^kv^\lambda$ in the binary-cubic covariant series \cite[Lemma~2.1 and p.~8]{BedratyukCovariants}. In the present variables this series is
\begin{equation}\label{reg:eq:chain-series}
 \mathcal C(t,v)=\frac{1+t^3v^3}{(1-tv^3)(1-t^2v^2)(1-t^4)}.
\end{equation}
Thus, for every
\begin{equation}\label{reg:eq:quadruples}
a,b,c\ge0,\qquad e\in\{0,1\},\qquad a+2b+4c+3e=k,
\end{equation}
there is one chain of length $3a+2b+3e+1$.

Rescale the basis of a chain to $e_0,\ldots,e_{\ell-1}$ so that $Ne_r=e_{r-1}$ for $r\ge1$ and $Ne_0=0$. Its natural grades are consecutive, say $b,b+1,\ldots,b+\ell-1$. For $F=\sum_rc_re_r$, the coefficient of $e_i$ in $e^{tN}F$ is $\sum_{r\ge i}c_rt^{r-i}/(r-i)!$. This is the $i$th derivative of $p(t)=\sum_rc_rt^r/r!$. The \hyperref[reg:balanced-hermite]{balanced Hermite selection} therefore gives actual chain-coordinate outputs of grades $b+\lfloor i/M\rfloor$, $0\le i<\ell$, at the $M$ nodes. The input-grade sum is $b\ell+\ell(\ell-1)/2$ and the output-grade sum is $b\ell+\sum_{i<\ell}\lfloor i/M\rfloor$.

For a chain of length $\ell$, define
\[
e_M(\ell)=\sum_{i=0}^{\ell-1}\bigl(i-\lfloor i/M\rfloor\bigr).
\]
The rational chain changes preserve each natural-weight space, so every dual chain-coordinate row is a linear combination of sampled monomial coefficient rows of that weight. Undoing the input change of basis leaves the independent Hermite rows invertible on all original columns. Lemma~\ref{reg:lem:transfer} therefore selects a complete minor of sampled monomial coefficients with the same total natural gain.

The removed constant $w^k$ is the bottom vector of the unique chain of length $3k+1$. Quotienting by that $N$-invariant line shortens that chain to length $3k$, with consecutive natural grades still intact. Consequently the complete natural gain is
\begin{equation}\label{reg:eq:EQ}
E^Q_M(k)=\sum_{\eqref{reg:eq:quadruples}}e_M(3a+2b+3e+1)-e_M(3k+1)+e_M(3k).
\end{equation}

\paragraph{A lower bound from the chain decomposition}
Let
\begin{equation}\label{reg:eq:AQ}
A_Q(k)=\sum_{\eqref{reg:eq:quadruples}}\frac{(3a+2b+3e+1)(3a+2b+3e)}2-3k.
\end{equation}
Applying $\lfloor i/M\rfloor\le i/M$ on the actual, shortened chains gives
\begin{equation}\label{reg:eq:EQbound}
E^Q_M(k)\ge(1-M^{-1})A_Q(k).
\end{equation}

\begin{lemma}\label{reg:lem:AQ}
For every $k\ge2$, $A_Q(k)\ge(13/4)\chi_k$.
\end{lemma}
\begin{proof}
With $D_v=v\partial_v$, the generating series of \eqref{reg:eq:AQ} is
\[
\begin{aligned}
 \sum_{k\ge0}A_Q(k)t^k&=\left.\frac{D_v^2+D_v}{2}\mathcal C(t,v)\right|_{v=1}-\frac{3t}{(1-t)^2}\\
 &=\frac{t(3t^6-6t^5+3t^4+t^3+3t^2+6t+3)}{(1-t)^5(1+t)(1+t^2)}.
\end{aligned}
\]
Subtract $(13/4)t/(1-t)^5$ and expand by residue classes modulo four. The resulting coefficient formulas for $G_k=A_Q(k)-(13/4)\chi_k$ are
\begin{equation}\label{reg:eq:Qgap}
\begin{aligned}
 4G_{4r}&=r(48r^2+58r-29),\\
 4G_{4r+1}&=48r^3+94r^2+9r-1,\\
 4G_{4r+2}&=48r^3+130r^2+65r+7,\\
 4G_{4r+3}&=48r^3+166r^2+139r+33.
\end{aligned}
\end{equation}
For $k\ge2$, the first two residue classes have $r\ge1$, and their displayed polynomials are positive there. The last two are positive for $r\ge0$. Hence $G_k\ge0$ for every $k\ge2$.
\end{proof}

\begin{proposition}[Sampling bound for \texorpdfstring{$Q_t$}{Qt}]\label{reg:prop:Q}
In this proposition $u,v,w$ denote logarithmic diagonal gaps. For $u,v,w\ge0$ and diagonal weights
\[
 (a,b,c)=(u+v+w,v+w,w),
\]
the complete $Q$ sampling matrix has a nonzero minor with determinant exponent $D$ in \eqref{reg:eq:character} at least
\begin{equation}\label{reg:eq:Qgain}
 \min(u,v,w)E^Q_M(k)\ge\min(u,v,w)\frac{13}{4}(1-M^{-1})\chi_k\qquad(k\ge2).
\end{equation}
\end{proposition}
\begin{proof}
Use the evaluation minor obtained from \eqref{reg:eq:EQ}. For a nonzero sample entry with input exponent $\mu$ and output exponent $\nu$, $\nu_x\le\mu_x$, $\nu_x+\nu_y\le\mu_x+\mu_y$, $|\nu|\le|\mu|$. Select a nonzero permutation term in the minor and sum these inequalities. They give
\[
 \Delta_1=\chi_k-R_x\ge0,\qquad \Delta_2=2\chi_k-R_x-R_y\ge0,\qquad \Delta_3=3\chi_k-R_x-R_y-R_z\ge0.
\]
Their sum is the natural gain $E^Q_M(k)$. The anisotropic gain is $u\Delta_1+v\Delta_2+w\Delta_3$, proving the first inequality. The second follows from \eqref{reg:eq:EQbound} and Lemma~\ref{reg:lem:AQ}.
\end{proof}


\paragraph{Mixed polynomial chains.}\label{reg:mixed-chains}
Fix $k\ge14$ for the mixed calculation used in Subsection~\ref{reg:uniform}.
For $M_f=(x+ty,y,z+t)$, homogenization gives $N_0=y\partial_x+w\partial_z$. With $U_0=x\partial_y+z\partial_w$ and $H_0=[U_0,N_0]$, the bidegree $(a,b)$ space, $a+b=k$, is $\operatorname{Sym}^a(\mathbb R^2)\otimes\operatorname{Sym}^b(\mathbb R^2)$. The Clebsch--Gordan rule \cite[Theorem 11.18]{EtingofSL2} gives one chain of each length $k-2j+1$, $0\le j\le\min(a,b)$. These chains have rational lowest vectors $(xw-yz)^j y^{a-j}w^{b-j}$ and are obtained by applying $U_0$, so the decomposition descends from the complex to the real space. The natural mixed grade is $(H_0+kI)/2$; thus both this grade and bidegree are preserved. Summing over $j\le a\le k-j$ gives multiplicity $k-2j+1$ for a chain of that length.

The natural mixed weight is $\deg_x+\deg_z$. Both the bidegree and this weight are respected by the chain changes of basis. Applying the \hyperref[reg:balanced-hermite]{balanced Hermite selection} and Lemma~\ref{reg:lem:transfer} therefore gives a physical complete minor with gain
\begin{equation}\label{reg:eq:Emixed}
E_M^{\rm mix}(k)=\sum_{j=0}^{\lfloor k/2\rfloor}(k-2j+1)e_M(k-2j+1)-e_M(k+1)+e_M(k).
\end{equation}
The last two terms again shorten the chain containing the constant before estimation. The selected rows preserve the bidegree multiplicities, including the grade $\deg_x+\deg_y$. Coordinate permutation gives the corresponding grades for $M$ and $M_e$.

Put
\[
A_{\rm mix}(k)=\sum_{j=0}^{\lfloor k/2\rfloor}\frac{(k-2j+1)^2(k-2j)}2-k.
\]
Then $E_4^{\rm mix}(k)\ge(3/4)A_{\rm mix}(k)$. If $k=2r$, replace $k-2j+1$ by $2h+1$; the summand becomes $4h^3+4h^2+h$, summed over $0\le h\le r$, with $2r$ subtracted afterwards. If $k=2r+1$, the summand becomes $4h^3+10h^2+8h+2$, and $2r+1$ is subtracted. Summing these polynomials and expanding $\chi_k$ gives
\begin{align*}
A_{\rm mix}(2r)-\tfrac32\chi_{2r}&=\frac{r(4r^2+9r-19)}{12},\\
A_{\rm mix}(2r+1)-\tfrac32\chi_{2r+1}&=\frac{4r^3+15r^2-7r-6}{12}.
\end{align*}
The first is nonnegative for $r\ge2$; the second is positive for $r\ge1$, as is seen by replacing $r$ with $1+s$. Thus
\begin{equation}\label{reg:eq:mixed-bound}
E_4^{\rm mix}(k)\ge\frac98\chi_k\qquad(k\ge3).
\end{equation}



\subsection{Proof of the uniform coefficient estimate}\label{reg:uniform}\label{reg:statement}\label{reg:scale-tables}

Throughout this subsection $k\ge14$. The notation and augmented stop are those of Section~\ref{reg:common-pieces}.
\paragraph{Scale domains and covering costs.}

The parameters in Table~\ref{reg:all-regimes} are logarithms to base $H$ of the positive quantities named there. The case of bounded $H$ has already been assigned to bounded pieces; for $H>1$, the displayed polytopes suppress powers of $P$. For example, $u\le P^CH$ gives $\log_Hu\le1+C\log P/\log H$. These shifts change each fixed linear objective by $O(\log P/\log H)$ and hence change the final power bound by at most a factor $P^C$. If $(a_x,a_y,a_z)$ is the row of Table~\ref{reg:all-costs}, then $(s_x,s_y,s_z)=(H^{a_x},H^{a_y},H^{a_z})$, and the same row's $n$ means that one additional sample costs at most $P^CH^n$. Case 9 deliberately uses the logarithm of the original rate $u<1$. In case 22, $T=\max(1,u_*W)$; the original net $HU^2V_0^2W^3$ is bounded by $P^CH^3W^3$, and enlarging to $1\le T,W\le P^CH$ can only weaken the eventual estimate. Each row follows by taking logarithms of the original boxes established in Subsection~\ref{reg:geometric-regimes}.

{\small
\begin{longtable}{r p{.31\textwidth} c p{.17\textwidth} p{.27\textwidth}}
\caption{The twenty-two regular cases. The parameter column lists the quantities whose logarithms to base $H$ are $\theta$. The denominator in the model column is the number of sampling points.}\label{reg:all-regimes}\\
ID & Case & Model/$m$ & Parameters & Logarithmic region\\\hline
\endfirsthead
ID & Case & Model/$m$ & Parameters & Logarithmic region\\\hline
\endhead
1 & $D$: first shear, bounded terminal & $L/18$ & --- & $\{0\}$\\
2 & $D$: first shear, large terminal & $Q/18$ & $W$ & $0\le\theta_1\le1$\\
3 & $D$: terminal, bounded first shear & $P/18$ & --- & $\{0\}$\\
4 & $D$: terminal, large first shear & $Q/18$ & $U$ & $0\le\theta_1\le1$\\
5 & $D$: mixed bounded & $M/5$ & --- & $\{0\}$\\
6 & $D$: mixed linear & $L/18$ & $U$ & $0\le\theta_1\le1$\\
7 & $D$: mixed parabolic & $R/4$ & $w$ & $0\le\theta_1\le1$\\
8 & $D$: cubic, $u\ge1$ & $Q/18$ & $(u,w)$ & $0\le\theta_i\le1$\\
9 & $D$: cubic, $u<1$ & $Q/18$ & $(u,w)$ & $-1\le\theta_1\le 0,\ 0\le\theta_2\le 1,\ -\theta_1-\theta_2\le 0$\\
10 & $T_z$: first shear, $u\le H$ & $M_f/4$ & --- & $\{0\}$\\
11 & $T_z$: first shear, $u\ge H$ & $M_f/4$ & $u$ & $1\le\theta_1\le 2$\\
12 & $T_z$: central, $u\le H$ & $P_f/18$ & $\max(1,u)$ & $0\le\theta_1\le1$\\
13 & $T_z$: ordered parabolic & $P_f/18$ & $u$ & $1\le\theta_1\le 2$\\
14 & $T_z$: ordered cubic & $Q/18$ & $(r,u)$ & $0\le\theta_1\le 1,\ 1\le\theta_2\le 2,\ -\theta_1-\theta_2\le -2$\\
15 & $A$: bounded product & $P_a/18$ & $(U,W)$ & $0\le\theta_i\le1,\ \theta_1+\theta_2\le 1$\\
16 & $A$: large product & $Q/18$ & $(u,w)$ & $0\le\theta_i\le1,\ -\theta_1-\theta_2\le -1$\\
17 & Central bounded, $U\ge W$ & $M_e/4$ & $(U,V_0,W)$ & $0\le\theta_i\le1,\ \theta_1+\theta_2\le 1,\ \theta_2+\theta_3\le 1,\ -\theta_1+\theta_3\le 0$\\
18 & Central bounded, $W\ge U$ & $M_e/4$ & $(U,V_0,W)$ & $0\le\theta_i\le1,\ \theta_1+\theta_2\le 1,\ \theta_2+\theta_3\le 1,\ \theta_1-\theta_3\le 0$\\
19 & Central, first product large & $L_e/18$ & $(U,V_0)$ & $0\le\theta_i\le1,\ -\theta_1-\theta_2\le -1$\\
20 & Central, both products large & $Q/18$ & $(U,V_0,W)$ & $0\le\theta_i\le1,\ -\theta_1-\theta_2\le -1,\ -\theta_2-\theta_3\le -1$\\
21 & Swapped, bounded $u_*W$ & $R/4$ & $(U,V_0,W)$ & $0\le\theta_i\le1,\ \theta_1+\theta_2\le 1,\ -\theta_2-\theta_3\le -1$\\
22 & Swapped, large $u_*W$ & $Q/18$ & $(T,W)$ & $0\le\theta_i\le1$\\
\end{longtable}
}

{\small
\begin{longtable}{r p{.49\textwidth} p{.30\textwidth}}
\caption{Diagonal exponents and original covering exponents for one additional sample.}\label{reg:all-costs}\\
ID & $(a_x,a_y,a_z)$ & $n$\\\hline
\endfirsthead
ID & $(a_x,a_y,a_z)$ & $n$\\\hline
\endhead
1 & $(2,\ 1,\ 0)$ & $6$\\
2 & $(2+\theta_1,\ 1+\theta_1,\ \theta_1)$ & $6+3\theta_1$\\
3 & $(2,\ 2,\ 1)$ & $6$\\
4 & $(2+\theta_1,\ 2,\ 1)$ & $6+3\theta_1$\\
5 & $(1,\ 1,\ 0)$ & $3$\\
6 & $(1+\theta_1,\ 1,\ 0)$ & $3+3\theta_1$\\
7 & $(1,\ 1+\theta_1,\ \theta_1)$ & $3+3\theta_1$\\
8 & $(1+\theta_1+\theta_2,\ 1+\theta_2,\ \theta_2)$ & $3+3\theta_1+3\theta_2$\\
9 & $(1+\theta_1+\theta_2,\ 1+\theta_2,\ \theta_2)$ & $3+3\theta_2$\\
10 & $(1,\ 0,\ 1)$ & $5$\\
11 & $(1,\ 0,\ 1)$ & $1+4\theta_1$\\
12 & $(2,\ \theta_1,\ 1)$ & $5$\\
13 & $(2,\ \theta_1,\ 1)$ & $3+2\theta_1$\\
14 & $(\theta_1+\theta_2,\ \theta_2,\ 1)$ & $-2+3\theta_1+4\theta_2$\\
15 & $(2,\ 1,\ \theta_2)$ & $4+2\theta_1+\theta_2$\\
16 & $(1+\theta_1+\theta_2,\ \theta_1+\theta_2,\ \theta_2)$ & $2+4\theta_1+3\theta_2$\\
17 & $(1+\theta_1,\ 1,\ \theta_1)$ & $2+2\theta_1+2\theta_2+\theta_3$\\
18 & $(1+\theta_3,\ 1,\ \theta_3)$ & $2+\theta_1+2\theta_2+2\theta_3$\\
19 & $(1+\theta_1,\ 1,\ 1-\theta_2)$ & $3+3\theta_1$\\
20 & $(\theta_1+\theta_2+\theta_3,\ \theta_2+\theta_3,\ \theta_3)$ & $3\theta_1+3\theta_2+3\theta_3$\\
21 & $(1,\ 1+\theta_3,\ \theta_3)$ & $1+2\theta_1+2\theta_2+3\theta_3$\\
22 & $(1+\theta_1,\ 1+\theta_2,\ \theta_2)$ & $3+3\theta_2$\\
\end{longtable}
}

\begin{proof}[Proof of Theorem~\ref{reg:regular-tail}]
\par\medskip\noindent\textit{Step 1. Common cubic lower bound.}\par\nobreak\smallskip

Here $u,v,w$ denote logarithmic diagonal gaps. Consider the region
\begin{equation}\label{reg:eq:central-domain}
 0\le u,v,w\le1,\qquad u+v\ge1,\qquad v+w\ge1,
\end{equation}
with weights $(u+v+w,v+w,w)$. Use the same $M$ nodes for all three certificates above. By \eqref{reg:eq:ABTbounds}, Proposition~\ref{reg:prop:L} gives the lower bound
\begin{equation}\label{reg:eq:central-L}
 \left(\frac12u+\frac32v-\frac2M\max\left(v,\frac{u+v}{2}\right)\right)\chi_k.
\end{equation}
To justify substitution of $\mathsf A_k\le\chi_k/2$, first write its coefficient as $-(u+v)+2d/M\le0$, where $d=\max(v,(u+v)/2)\le u+v$. The same calculation with Proposition~\ref{reg:prop:P} gives
\begin{equation}\label{reg:eq:central-P}
 \left(\frac12v+\frac32w-\frac2M\max\left(w,\frac{v+w}{2}\right)\right)\chi_k-2k.
\end{equation}
Here $\max(w,(v+w)/2)\le1$, so the full constant-quotient correction is at most $2k$. The functions in the parentheses are coordinatewise nondecreasing for $M\ge2$.

By \eqref{reg:eq:central-domain}, reducing $u,w$ to $1-v$ can only decrease these two lower bounds. Put $s=\min(v,1-v)\in[0,1/2]$. The larger of the two conic bounds is at least
\begin{equation}\label{reg:eq:conic-max}
 \left(\frac32-\frac2M-\left(1-\frac2M\right)s\right)\chi_k-2k.
\end{equation}
Indeed the larger argument is $1-s\ge1/2$, on which the relevant maximum is $1-s$. Meanwhile $\min(u,v,w)\ge s$, and the principal certificate gives at least $(13/4)(1-M^{-1})s\chi_k$.

Set $M=18$. The largest of the three complete minor gains is therefore at least
\[
 \max\left\{\frac{25}{18}-\frac89s,\frac{221}{72}s\right\}\chi_k-2k.
\]
The decreasing and increasing affine functions intersect at $s=20/57$, where both equal $1105/1026$. Hence
\begin{equation}\label{reg:eq:central-complete}
 \max D\ge\frac{1105}{1026}\chi_k-2k.
\end{equation}
Every case to which we apply this bound has additional-sample net exponent $n\le9$. The $M=18$ samples give seventeen additional nets, with total cost exponent $17n\le17\cdot9=153$. Thus \eqref{reg:eq:witness} gives exactly the tail exponent \eqref{reg:tail}. The three choices of sampling data use the same sample tuple; their finite choice changes only a constant, not the number of samples.

\par\medskip\noindent\textit{Step 2. Cubic cases.}\par\nobreak\smallskip
The case numbers are those of Table~\ref{reg:all-regimes}. Each coefficient exponent is $D-(M-1)n$ from \eqref{reg:eq:witness}, using the determinant exponent and the net cardinality in that case.

\paragraph{The six ordered cubic cases.}
For IDs 2, 4, 8, 14, 16 and 20, put $(u,v,w)=(a_x-a_y,a_y-a_z,a_z)$, using the diagonal exponents and covering exponents $n$ of Table~\ref{reg:all-costs}. The resulting gap coordinates are recorded in Table~\ref{app:table-ordered-cubic}. The parameter restrictions in Table~\ref{reg:all-regimes} give $0\le u,v,w\le1$, $u+v\ge1$, $v+w\ge1$ and $n\le9$ in each of these six cases. Thus \eqref{reg:tail} covers all six.

\paragraph{The two remaining cubic cases.}
In ID 9 the weights are $(1-r+v,1+v,v)$ with $0\le r\le v\le1$, and $n=3+3v\le6$. In ID 22 the weights are $(1+u,1+v,v)$ with $0\le u,v\le1$, and again $n=3+3v\le6$.

For a homogeneous conic bound valid also when $a-b<0$, use the separate estimates $\mathsf A_k\le\chi_k/2$ and $\mathsf B_{k,M}\le3\chi_k/M$. In ID 9 this gives
\[
D_L\ge\left(\frac43-\frac r2\right)\chi_k\qquad(M=18),
\]
since $A=1-r$, $B=1$, and $d=1$. The affine parabolic certificate uses weights $(1+v,v)$; here $d=(1+v)/2$ and its coefficient of $\mathsf A_k$ is nonpositive. It gives
\[
D_P\ge\left(\frac49+\frac{13}{9}v\right)\chi_k-2k
 \ge\left(\frac49+\frac{13}{9}r\right)\chi_k-2k.
\]
The two affine bounds intersect at $r=16/35$, with common value $116/105$. Thus their maximum is at least $(116/105)\chi_k-2k$.

For ID 22, the homogeneous bound is at least $(4/3-v/2)\chi_k$, since $A=1+u-v$, $B=1$, and $d=1$. The same parabolic bound applies. Its combination with the homogeneous bound gives the identical result, with $r=v$. The product of the seventeen additional-sample net cardinalities is at most $P^CH^{102}$, so these two cases have exponent at least
\begin{equation}\label{reg:eq:otherQ}
 \frac{116}{105}\chi_k-2k-102>\beta_k.
\end{equation}
The separate use of $\mathsf B_{k,M}\le3\chi_k/M$ here is deliberate: replacing $\mathsf A_k$ by its upper bound after combining its coefficient would have the wrong sign near $r=1$.

\par\medskip\noindent\textit{Step 3. Conic and mixed cases.}\par\nobreak\smallskip
ID 1 has $L$ weights $(2,1,0)$; its gain is at least $(1-M^{-1})\mathsf T_k\ge(17/9)\chi_k$. ID 6 has $L$ weights $(1+\theta,1,0)$, $0\le\theta\le1$; its gain is at least $(25/18+\theta/2)\chi_k$. Both have $n\le6$.

IDs 3, 12, 13, and 15 have the two active parabolic weights $(2,1)$, with an arbitrary spectator scale. Their gains are at least $(17/9)\chi_k-2k$ by Proposition~\ref{reg:prop:P}. The largest exponent in the net cardinality among them is $7$. ID 19, the central-linear model $L_e$, has active parabolic weights $(1+\theta_1,1)$ and spectator weight $1-\theta_2$. The exact $z$-grade transfer gives gain at least $(25/18+\theta_1/2)\chi_k-2k$, independently of that spectator weight; here $n=3+3\theta_1\le6$. Therefore every case in this subsection has tail exponent at least
\begin{equation}\label{reg:eq:conic-tail}
 \frac{25}{18}\chi_k-2k-119>\beta_k.
\end{equation}

\paragraph{Mixed families and the full parabolic model $R$.}
Use the mixed interpolation bound~\eqref{reg:eq:mixed-bound}.

The remaining cases have the following maximum total covering exponents:
\begin{center}\small
\begin{tabular}{cclrr}
\toprule
ID & Model & Preserved grade & Samples & Total covering exponent\\\midrule
5 & $M$ & $\deg_y+\deg_z$ & 5 & $12$\\
7 & $R$ via $M$ & $\deg_y+\deg_z$ & 4 & $18$\\
10 & $M_f$ & $\deg_x+\deg_y$ & 4 & $15$\\
11 & $M_f$ & $\deg_x+\deg_y$ & 4 & $27$\\
17, 18 & $M_e$ & $\deg_x+\deg_z$ & 4 & $15$\\
21 & $R$ via $M$ & $\deg_y+\deg_z$ & 4 & $18$\\\bottomrule
\end{tabular}
\end{center}
In the $R$ cases the weights are $(1,1+v,v)$. Their $v$ component is exactly the preserved grade $\deg_y+\deg_z$, so it cancels. Lemma~\ref{reg:lem:transfer} makes the determinant equal to the mixed determinant, with every translation retained in off-diagonal blocks. The analogous cancellation in $M_f$ and $M_e$ is the original grading. The five-state case has $E_5^{\rm mix}\ge E_4^{\rm mix}$. Consequently these cases all have exponent at least
\begin{equation}\label{reg:eq:mixed-tail}
 \frac98\chi_k-27>\beta_k.
\end{equation}
This finishes the coverage of all twenty-two cases.

\par\medskip\noindent\textit{Step 4. Actual samples and completion of the regular estimate.}\par\nobreak\smallskip
Subtracting the required incidence threshold from \eqref{reg:tail} gives
\[
 \beta_k-(\chi_k+2)=\frac{79}{1026}\chi_k-2k-155.
\]
For $k=14+r$, direct polynomial expansion yields
\begin{equation}\label{reg:eq:positive}
 \beta_{14+r}-(\chi_{14+r}+2)
 =\frac{79r^4}{24624}+\frac{2449r^3}{12312}+\frac{113681r^2}{24624}+\frac{560687r}{12312}+\frac{131}{513}.
\end{equation}
Every coefficient is positive. Thus the assertion holds for all $k\ge14$.

Fix $k\ge14$ before choosing the geometric comparison fractions. The evaluation-row constructions above give finitely many rational complete minors. At their normalized model tuples their absolute determinants have positive lower bounds depending on $k$ and the chosen progression. Lemma~\ref{reg:regular-preparation} allows the prescribed sample count, and Subsection~\ref{reg:geometric-regimes} gives the normalized model comparisons, including the complete variation-of-constants estimate \eqref{reg:swapped-field}. Choose the comparison errors smaller than the continuity radii of these finitely many determinants. The corresponding actual Gram determinants then satisfy the same powers of $H$, up to $P^C$. All large additional entries have already been removed by triangular identities with equal row and column grade multiplicities, so determinant continuity is applied only to bounded normalized entries.

Apply~\eqref{reg:fixed-net-bound} with the original nets of Tables~\ref{reg:all-regimes} and~\ref{reg:all-costs}. The exponent~\eqref{reg:eq:witness} and the bounds above give $P^CH^{-\beta_k}$ for each fixed choice. In the six ordered cubic cases the three possible minors use the same eighteen sample states; their finite choice contributes a constant.

The finitely many chart, sign, change of frame, and row choices cost a power of $P$. Lemma~\ref{reg:regular-preparation} bounds the number of bands and prepared intervals by $P^C(1+\log P+\log d_1^{-1})^C$. Summing the nonnegative coefficient estimates with the original label fixed over these common choices proves~\eqref{reg:tail}. That lemma also supplies the count and bounded-piece assertions, completing the proof of Theorem~\ref{reg:regular-tail}.
\end{proof}

For $q=(1+\eta)(p-1)<\beta_k$, Theorem~\ref{reg:regular-tail} and the count in Lemma~\ref{reg:regular-preparation} supply the regular input to Proposition~\ref{an:common-incidence}.

\section{Coefficient coordinates and finite slope estimates}\label{atlas:section}

Use the fixed-degree conventions and conditioning parameter $P\ge2$ of Section~\ref{geo:section}. Coefficient measure is taken on $V_k$ or an explicitly specified slice. This section constructs the local coefficient maps and fixed-frame integration used in Section~\ref{surf:section}; Section~\ref{atlas:parent} restores the coefficients lost under normal differentiation. The finite slope estimates apply to the derivative degrees six through nine.

\subsection{Restriction and the normal filtration}\label{atlas:restriction}

The examples in Subsection~\ref{geo:normalization} show that a lower-order pivot does not permit removal of higher normal coefficient directions. The atlases retain these directions and, when differentiation describes an atlas, every coefficient in its kernel.

\begin{definition}[Coefficient maps and jets]\label{def:coefficient-atlas}
\leavevmode\noindent\textup{(i)} A complete coefficient atlas is a specified square map from coefficients to the stated outputs with nonzero Jacobian on the full coefficient space or the stated normalization slice. This gives local coordinates; the bound on the number of preimages used in integration is supplied separately by Lemma~\ref{geo:coefficient-map}.

\noindent\textup{(ii)} A fixed-graph map holds the graph fixed; the moving-graph map uses the selected graph of the varying phase.

\noindent\textup{(iii)} External data are the Taylor coefficients of the coordinate or seed map through the orders used by the stated output list. After choosing a representative $\zeta$ from their fixed net, hold these coefficients fixed during coefficient differentiation and integration. Thus the map is $F\mapsto\mathcal T_\zeta(F)$ on the stated branch or slice, and its Jacobian differentiates only those coefficient coordinates.

Subscripts listing variables denote raw partial derivatives, so $J_{x^ay^bz^c}=\partial_x^a\partial_y^b\partial_z^cJ$. Ordinary jets are Taylor coefficients in the specified external coordinates. Gauss jets are Taylor coefficients in the specified Gauss coordinates. Ordered null jets are $(Y^iZ^jf)/(i!j!)$. The selected graph and the intrinsic fields remain phase-dependent when coefficients vary; their coefficient derivatives are retained when the external seed parameters are fixed.
\end{definition}
\begin{definition}[Jet index sets]\label{def:jet-indices}
All pairs $(i,j)$ below belong to $\mathbb Z_{\ge0}^2$. For an integer $d\ge0$ define
\[
I_d=\{(i,j):0\le i,j\le d\},\qquad T_d=\{(i,j):i+j\le d\},\qquad I_d^{\mathrm w}=\{(i,j):i+2j\le2d\}.
\]
In an index-set difference, $0$ denotes $(0,0)$. The sets $I_d,T_d,I_d^{\mathrm w}$ specify rectangular, triangular and weighted jets, respectively. Sets with a negative subscript are empty. A subscript set on a scalar function denotes its Taylor coefficients at the evaluation point, in the specified fixed external tangential coordinates. Derivatives can be used instead, changing only constants depending on $k$. For $d\ge3$, set
\[
E_d=I_{d-3}\cup\{(d-2,0),(0,d-2)\},\qquad E_d^{\mathrm w}=I_{d-3}^{\mathrm w}\cup\{(0,d-2),(1,d-3)\}.
\]
An index set is downward closed if decreasing either nonnegative coordinate keeps the index in the set. Every $E_d$ or $E_d^{\mathrm w}$ used here is downward closed and contains the complete jet $T_{d-2}$. At a normalized evaluation point the graph has $c(0)=c_y(0)=c_z(0)=0$. In the real isotropic chart it additionally has $c_{zz}(0)=0$.
\end{definition}

\begin{lemma}[The normal filtration]\label{surf:normal-filtration}

\leavevmode\noindent\textup{(i)} Let $Q=Ay^2+Byz$ with $B\ne0$. On polynomials of degree at most $k$, including the constant coefficient, the fixed-graph map
\[
R\longmapsto\bigl((\partial_x^rR)|_{x=Q}\bigr)_{I_{k-2r}},\qquad 0\le r\le\lfloor k/2\rfloor,
\]
is square and invertible. Its absolute determinant, with Taylor coefficients in the tangential variables and normal derivatives as displayed, is $\prod_{r=0}^{\lfloor k/2\rfloor}(r!)^{(k-2r+1)^2}|B|^{\tau_k}$, where
\[
 \tau_k=\sum_{r=0}^{\lfloor k/2\rfloor}\binom{k-2r+2}{3}.
\]

\noindent\textup{(ii)} For an arbitrary quadratic graph $Q=Ay^2+Byz+Cz^2$, the weighted index sets $I_{k-2r}^{\mathrm w}$ give the same factorial factor times $|A|^{\tau_k}$ in the absolute determinant. The weighted map is invertible when $A\ne0$.

\end{lemma}
\begin{proof}
Let $Q(Y,z)=AY^2+BYz$ with $B\ne0$, and let $m\ge0$ be fixed. Restriction of the degree-at-most-$m$ polynomial space to $X=Q$ has kernel $(X-Q)\mathbb R[X,Y,z]_{\le m-2}$, where a negative degree means the zero space. Its image has dimension $(m+1)^2$. The square grid outputs $[Y^iz^j]H(Q,Y,z)$, $0\le i,j\le m$, form coordinates on that image, with restriction determinant $B^{m(m+1)(m+2)/6}$ up to sign in the monomial quotient basis described below.
For each $d$, let $\mathcal C_d$ be the span of all binary monomials of degree at most $d$, together with all monomials $x^ay^bz^c$ with $a\ge1$ and $a+b+c=d$. Its dimension is $(d+1)^2$. At $Q=Byz$, every index $(i,j)\in I_d\setminus T_d$ is the restriction of exactly one of these additional monomials: take $a=i+j-d$, $b=d-j$, and $c=d-i$. Its coefficient is $B^a$. Together with the binary monomials, these give a diagonal restriction matrix after reordering. The sum of these powers $a$ is $\binom{d+2}{3}$. Adding $Ay^2$ increases the first tangential index and decreases the second in each nonleading term. Ordering by the first index makes the projected restriction matrix triangular with the same diagonal. This proves its determinant formula in the grid chart.

For the weighted chart first take $Q=Ay^2$. For $(i,j)\in I_d^{\mathrm w}\setminus T_d$, the unique additional preimage has $a=i+j-d$, $b=2d-i-2j$, and $c=j$. These are nonnegative and $a>0$, and its coefficient is $A^a$. The sum of the powers is again $\binom{d+2}{3}$. Give $x,y,z$ weights $2,1,2$. The terms $Byz$ and $Cz^2$ have strictly larger weight than $Ay^2$, so they enter strictly triangular blocks and do not change this determinant.

Let $q=x-Q$. The kernel of restriction to the whole graph is $q\mathbb R[x,y,z]_{\le k-2}$: polynomial division in $x$ proves divisibility, and the nonzero highest degree part $-Q$ proves the degree bound on the quotient. This kernel has codimension $(k+1)^2$. Since the projected restriction on $\mathcal C_k$ has that same dimension and is invertible, the kernel of projected restriction has exactly the same dimension as the whole-graph kernel and therefore equals it. Consequently
\[
 \mathbb R[x,y,z]_{\le k}=\mathcal C_k\oplus q\mathbb R[x,y,z]_{\le k-2}.
\]
Repeating yields a basis consisting of $q^r\mathcal C_{k-2r}$. The coefficient change has a nonzero constant determinant. At one division stage, the monomials missing from $\mathcal C_k$ are exactly the positive-$x$ monomials of degree at most $k-1$. The term $xR$ in $qR=xR-QR$ maps the monomial basis of degree at most $k-2$ bijectively to these missing rows. Every positive-$x$ term of $QR$ has total degree one larger than the corresponding row of $xR$, and its binary terms lie in $\mathcal C_k$. Descending total-degree order therefore makes the remaining block triangular with diagonal one. Repeating this argument proves the constant basis determinant. On the filtration, the order-$r$ normal restriction has diagonal factor $r!$ and the restriction determinant just computed; lower normal orders vanish on subsequent filtration terms. Multiplying the diagonal determinants proves the assertion. It also proves
\[
 \sum_{r=0}^{\lfloor k/2\rfloor}(k-2r+1)^2=\binom{k+3}{3}.
\]
\end{proof}

Removing the constant input and its identity output row from Lemma~\ref{surf:normal-filtration} gives the complete fixed-graph map on $V_k$:
\begin{equation}\label{low:78:h3n:filtration-map}
 (F|_{x=Q})_{I_k\setminus\{(0,0)\}},\qquad
 \bigl((\partial_x^rF)|_{x=Q}\bigr)_{I_{k-2r}},\quad 1\le r\le\lfloor k/2\rfloor.
\end{equation}
Its determinant is unchanged. Only the original global constant of $F$ is removed; derivative phases retain their constant, as in Definition~\ref{def:normal-restrictions}.

\subsection{Coefficient maps on varying graphs}\label{atlas:moving}\label{low:78:h3n:section}
\begin{lemma}[Coefficient coordinates including \texorpdfstring{$F_{xx}$}{Fxx}]\label{surf:strong-atlas}
Let $D\ge6$ and $F\in V_D$. Let $x=c(y,z)$ satisfy $F_x(c(y,z),y,z)=0$, and put $\lambda=F_{xx}|_c\ne0$ and $\phi_r=\partial_x^rF|_c$. Let the first derivative of the tangential coordinate map whose required Taylor coefficients are held fixed be $M\in GL_2(\mathbb R)$, and write its transformed quadratic graph as $Ay^2+byz+Cz^2$.

\noindent\textup{(i)} The output list
\begin{equation}\label{surf:ordinary-list}
 (\phi_0)_{I_D\setminus\{(0,0)\}},\quad (\phi_2)_{I_{D-3}},\quad c_{E_D},\quad (\phi_r)_{I_{D-2r}}\quad(3\le r\le\lfloor D/2\rfloor)
\end{equation}
has $N_D$ entries. The absolute coefficient Jacobian is
\begin{equation}\label{surf:strong-jacobian}
C_D|\det M|^{\chi_D}|\lambda(0)|^{-|E_D|}|G_D(A,b,C)|,
\end{equation}
where $C_D>0$ depends only on coordinate conventions,
\begin{equation}\label{surf:quadratic-factor}
G_D(A,b,C)=b^{t_D}g_D(AC/b^2),\qquad g_D(0)\ne0,
\end{equation}
\equationalias{low:78:h3n:det-character}
and
\begin{equation}\label{surf:tD}
t_D=\chi_D-3(D-2)^2-2-\sum_{r=3}^{\lfloor D/2\rfloor}r(D-2r+1)^2.
\end{equation}
The displayed expression is interpreted polynomially when $b=0$.

\noindent\textup{(ii)} Replacing $I_d,E_D$ by $I_d^{\mathrm w},E_D^{\mathrm w}$ gives a nonzero constant times $|\det M|^{\chi_D}|\lambda(0)|^{-|E_D|}|A|^{t_D}$.

\noindent\textup{(iii)} These identities are independent of higher graph jets and of higher external tangential jets. The output maps have uniformly finite inverse fibers on every regular algebraic branch where their determinants are nonzero.
\end{lemma}
\begin{proof}
\par\medskip\noindent\textit{Step 1. Ordinary and weighted injectivity.}\par\nobreak\smallskip
First fix $Q=yz$ and include the constant critical-value row. If the critical rows of a variation $R$ vanish, then $R|_Q=0$ and $R=qH$, where $q=x-Q$ and $\deg H\le D-2$. The second-normal restriction is $R_{xx}|_Q=2H_x|_Q$. It is supported on $I_{D-3}$, so its retained rows force it to vanish identically. For $i,j\in\mathbb{Z}_{\ge0}$, set $c_a=[x^ay^{i-a}z^{j-a}]H$ for $0\le a\le\min(i,j)$. Then
\[
[y^iz^j]H(yz,y,z)=\sum_{a=0}^{\min(i,j)}c_a,\qquad
[y^{i-1}z^{j-1}]H_x(yz,y,z)=\sum_{a=0}^{\min(i,j)}a c_a\quad(i,j\ge1).
\]
If $i=D-2$ and $j>0$, the degree inequality $i+j-a\le D-2$ and the inequality $a\le j$ force $a=j$. Thus that boundary coefficient vanishes. The same argument applies with $i,j$ exchanged. The restriction $H|_Q$ is consequently supported on $E_D$. Its first-normal rows, represented by the graph rows, force $H|_Q=0$. Write $H=qH_1$. Since $H_x|_Q=H_1|_Q=0$, divide once more to obtain $R=q^3H_2$. The remaining higher-normal rows determine $H_2$ by Lemma~\ref{surf:normal-filtration}. Thus $R=0$.

At $Q=y^2$, put $d=D-2$. For a restricted index $(i,j)$ the possible powers $a$ in $x^ay^{i-2a}z^j$ satisfy $\max(0,i+j-d)\le a\le\lfloor i/2\rfloor$. An index outside $I_{d-1}^{\mathrm w}\cup T_d$ has only one possible positive power, and the equation $H_x|_Q=0$ eliminates its coefficient. The union $I_{d-1}^{\mathrm w}\cup T_d$ equals $E_D^{\mathrm w}$. Thus the weighted graph rows give the same two divisions, and the weighted filtration completes the injectivity proof.

\par\medskip\noindent\textit{Step 2. Output count and moving graph.}\par\nobreak\smallskip
The number of outputs including the constant is
\[
 (D+1)^2+2(D-2)^2+2+\sum_{r=3}^{\lfloor D/2\rfloor}(D-2r+1)^2=\binom{D+3}{3}.
\]
Hence the fixed-graph maps just proved injective are square. The constant input coefficient affects only the constant critical-value row, with derivative one. Deleting that row and column therefore gives the asserted map on $V_D$ with the same determinant. For the moving graph their variations are
\[
 \delta c=-\frac{R_x|_c}{\lambda},\qquad \delta\phi_0=R|_c,\qquad \delta\phi_r=R_{x^r}|_c+\phi_{r+1}\delta c.
\]
Every retained positive-order scalar index set is contained in the downward-closed graph set. The Leibniz rule therefore eliminates $\phi_{r+1}\delta c$ by graph rows. Multiplication by $\lambda$ on the graph jet is triangular with diagonal $\lambda(0)$, giving the factor $\lambda(0)^{-|E_D|}$.

\par\medskip\noindent\textit{Step 3. Weight filtration.}\par\nobreak\smallskip
Give the input monomial $x^ay^bz^c$ weight $2a+b+c$, and an order-$r$ restriction row at $(i,j)$ weight $2r+i+j$. Since the graph starts in degree two, an entry vanishes if its row weight is below its input weight. The homogeneous quadratic map is invertible and preserves weight, so the row and column counts agree separately at every weight. The complete matrix is therefore block triangular with square diagonal blocks. Higher graph terms and higher external seed terms change only strictly triangular blocks. Its determinant depends only on the quadratic graph and the linear part of the external coordinate map.

\par\medskip\noindent\textit{Step 4. Quadratic determinant and coefficient character.}\par\nobreak\smallskip
Under $y\mapsto sy$, $z\mapsto s^{-1}z$, the grid output has equal total first and second index sums. Its determinant is invariant. Every monomial in the quadratic determinant consequently contains the same powers of $A$ and $C$. Homogeneity in the quadratic coefficients gives degree $t_D$: the input $x$ powers sum to $\chi_D$, and subtracting the total row normal orders gives exactly \eqref{surf:tD}. This proves the form $b^{t_D}g_D(AC/b^2)$. Injectivity at $Q=yz$ proves $g_D(0)\ne0$. The second weight filtration used in Lemma~\ref{surf:normal-filtration} gives the weighted factor $A^{t_D}$.

Lemma~\ref{an:character} applied to $(x,y,z)\mapsto(x,M(y,z))$ gives the coefficient factor $|\det M|^{\chi_D}$. After this pullback the quadratic coefficients are precisely $A,b,C$ appearing in \eqref{surf:strong-jacobian}.

\par\medskip\noindent\textit{Step 5. Multiplicity and coefficient volume.}\par\nobreak\smallskip
The root equation, retained pivot, and finite derivative outputs give a fixed-format algebraic map. On the selected nonzero-Jacobian branches, Lemma~\ref{geo:coefficient-map} supplies the bounded multiplicity and coefficient-volume estimate.
\end{proof}

\paragraph{Shared finite-degree conventions and specializations}

Use the full $N_k$-dimensional monomial coefficient measure and the augmented stop~\eqref{geo:translation-stop}, whose physical matrix may be singular. The normal pivot and rank-four models are those of Lemma~\ref{geo:normal-pivot} and Theorem~\ref{geo:classification}, including the degree restrictions, converse, and polynomial parameter bounds. Degree seven has the same model types as degree six; degree eight adds $q_0^4$, and degree nine adds $q_\gamma^3$.

Use the constant-free filtration map~\eqref{low:78:h3n:filtration-map}. At a fixed conditioned model, replacing a selected normal block by graph jets gives local coefficient coordinates. The geometric assignment and scalar widths enter separately in the coefficient count.

For the coefficient map in finite degree, use exactly \eqref{surf:ordinary-list}. We absorb its fixed Taylor-coordinate constant into the quadratic polynomial and retain the scalar exponent $t_k$ of \eqref{surf:tD}. Thus the determinant convention in the finite calculations is
\begin{equation}\label{low:78:h3n:exact-jac}
 |\det M|^{\chi_k}|\lambda(0)|^{-|E_k|}|G_k(A,b,C)|,
\end{equation}
The factorization~\eqref{surf:quadratic-factor} is unchanged by absorbing the same fixed constant into $G_k$ and $g_k$. Here $t_k$ is given by \eqref{surf:tD} with $D=k$. The weighted formula, independence of higher graph and external seed jets, and bounded inverse multiplicity are those of Lemma~\ref{surf:strong-atlas}.

For $k=7$ the output dimensions in \eqref{surf:ordinary-list} are $63+25+27+4=119$, and $t_7=121$. The third-normal block has four entries. For $k=8$ the dimensions are $80+36+38+9+1=164$. Here $t_8=189$.

On determinant-one horizontal ribbons this character is one; on a conditioned collar it is a single conditioned factor. Changing the complete lower graph-jet outputs afterwards is a separate row-coordinate operation, with the triangular lower-jet factor stated in the corresponding count; it does not apply the full $\chi_k$ input character a second time.

\subsection{Recovery of all parent coefficients}\label{atlas:parent}\label{low:78:smult:section}

\begin{definition}[Normal restrictions and derivative phases]\label{def:normal-restrictions}
On an order-$m$ sheet of $F$, write $\phi_j=(\partial_x^jF)|_c$ and $\phi=\phi_0$. For the reduction with $m\ge3$, the derivative phase is $G=\partial_x^{m-2}F$, and the term parent phase refers to the original polynomial $F$. The graph is critical for $G$, whose critical value is $\phi_{m-2}$. Retain the constant coefficient of $G$ and all lower normal blocks; only the original global constant of $F$ is discarded.
\end{definition}
The lower normal blocks recover the coefficients of $F$ as follows.

\begin{proposition}[Lifting a derivative phase atlas]\label{low:78:smult:atlas}
Let $k\ge3$, $3\le m\le k$, $r=m-2$, and let $c$ be an order-$m$ sheet of $F\in V_k$. Suppose a complete ordinary coefficient atlas for $G=\partial_x^rF$ at this graph is given, including its constant coefficient direction.

\noindent\textup{(i)} Append the normal-scalar outputs of Definition~\ref{def:normal-restrictions}
\[
(\phi_j)_{\mathcal T_{k-j}},\qquad \mathcal T_{k-j}=T_{k-j},\qquad 0\le j<r,
\]
omitting the constant row only when $j=0$. The resulting map is a complete local coefficient atlas for $F$.

\noindent\textup{(ii)} If the fixed external tangential parametrization has linear part $M$, its Jacobian determinant is, up to a fixed nonzero constant, the derivative-phase atlas determinant times
\[
(\det M)^{\sum_{j=0}^{r-1}\binom{k-j+2}{3}}.
\]
\end{proposition}
\begin{proof}
Retain the constant of $G$ and decompose the polynomial uniquely as
\[
F(x,y,z)=\sum_{j=0}^{r-1}\frac{x^j}{j!}H_j(y,z)
+\frac{1}{(r-1)!}\int_0^x(x-s)^{r-1}G(s,y,z)\,ds,
\qquad \deg H_j\le k-j,
\]
where only $H_0$ is taken modulo constants. This is an invertible linear coefficient change with a fixed nonzero determinant. Its dimensions agree, since
\[
\binom{k-r+3}{3}
+\sum_{j=0}^{r-1}\binom{k-j+2}{2}-1
=\binom{k+3}{3}-1.
\]

The derivative phase atlas determines $G$ locally and therefore determines the selected graph $c$ locally. For $0\le j<r$, set $A_{j,G}(y,z)=\frac1{(r-j-1)!}\int_0^{c(y,z)}(c(y,z)-s)^{r-j-1}G(s,y,z)\,ds$. With these fixed, the highest remaining scalar satisfies $\phi_{r-1}=H_{r-1}+A_{r-1,G}$, where $A_{r-1,G}$ is a fixed function of $G$ and $c$. Its binary jet through degree $k-r+1$ determines $H_{r-1}$. Descending in $j$, one has
\[
\phi_j=H_j+\sum_{\ell=j+1}^{r-1}\frac{c^{\ell-j}}{(\ell-j)!}H_\ell+A_{j,G}.
\]
Every term except $H_j$ is already determined. The binary jet through degree $k-j$ therefore determines $H_j$. For $j=0$ the constant is irrelevant, and the remaining rows determine $H_0$ modulo constants.

This argument is also a differential factorization. The derivative phase outputs form the first invertible block. With their variations zero, $\delta G=0$ and hence $\delta c=0$. The lower blocks are triangular in descending $j$, with diagonal variation $\delta H_j$ in each scalar. The terms $(\partial_x^{j+1}F)|_c\,\delta c$ occur in the off-diagonal blocks. If the external tangential parametrization has linear part $M$, the map from the coefficients of $H_j$ to its Taylor coefficients through degree $k-j$ is triangular by total degree with determinant $(\det M)^{\binom{k-j+2}{3}}$. Higher terms of the parametrization affect only off-diagonal blocks in this total-degree ordering. Multiplying the lower-block determinants gives the asserted determinant factor.
\end{proof}

In a complete normalized parent admissible set, every retained positive normal scalar satisfies $|\phi_j|\le P^C$ for $j\ge1$. The parent value $\phi_0$ satisfies $|\partial_y\phi_0|+|\partial_z\phi_0|\le P^C$ in the tangent-normalized frame; in the null chart this gives $|Y\phi_0|+|Z\phi_0|\le P^C$. The coefficient calculation uses these respective bounds for the lower scalars. For $m\ge3$, the critical value of $G=\partial_x^{m-2}F$ is $\phi_{m-2}$. The bounded-gradient widths used for the parent value $\phi_0$ do not require the graph to be critical for $F$.

For the degree-seven and degree-eight applications, the complete derivative phase block includes its constant. The dimensions and external characters of this block and the lower parent blocks are recorded in Table~\ref{app:table-parent-blocks}. The derivative-phase character and the binary parent characters proved above have total exponent $\chi_k$, giving the external character in \eqref{low:78:h3n:exact-jac}.

\begin{lemma}[Recovery of every lower parent coefficient]\label{surf:parent-lift}
Suppose $F\in V_k$, $2\le m\le\lfloor k/2\rfloor$, the selected sheet satisfies $F_{x^{m-1}}|_c=0$, and $\kappa_m=F_{x^m}|_c$ satisfies $P^{-C}\le|\kappa_m|\le P^C$. Put $\phi_j=(\partial_x^jF)|_c$ and $D=k-m+2$, and assume $D\ge6$.

\noindent\textup{(i)} If $m=2$, use \eqref{surf:ordinary-list} in degree $k$. If $m\ge3$, use the degree-$D$ coefficient map for $G=\partial_x^{m-2}F$, including its constant critical-value row, and append
\[
 (\phi_j)_{T_{k-j}},\qquad 0\le j\le m-3,
\]
omitting only the constant of $\phi_0$. The resulting output map has $N_k$ components.

\noindent\textup{(ii)} Its determinant is a fixed nonzero constant times
\begin{equation}\label{surf:parent-jacobian}
 (\det M)^{\chi_k}\kappa_m^{-|E_D|}G_D(A,b,C).
\end{equation}
The weighted alternative has $A^{t_D}$ in place of $G_D$.

\noindent\textup{(iii)} The ordinary map is a local coefficient atlas wherever $G_D(A,b,C)\ne0$; the weighted map is a local coefficient atlas wherever $A\ne0$.
\end{lemma}
\begin{proof}
For $m=2$, the coefficient map and determinant formula are those of Lemma~\ref{surf:strong-atlas}; the condition $D=k\ge6$ is exactly its degree hypothesis. For $m\ge3$, use the differentiation-kernel coordinates $(G,H_{m-3},\ldots,H_0)$ in the proof of Proposition~\ref{low:78:smult:atlas}. The degree-$D$ outputs depend only on $G$, whose critical equation and nonzero second-normal pivot are precisely the order-$m$ hypotheses. The lower parent outputs are triangular in descending $j$, with diagonal blocks given by the complete binary jets of $H_j$. Their moving-graph terms lie in the off-diagonal blocks. This block determinant factorization holds even when the degree-$D$ block is singular. Lemma~\ref{surf:strong-atlas}, with the constant critical-value row restored, gives its determinant. The dimension and external character identities are
\[
 \binom{D+3}{3}+\sum_{j=0}^{m-3}\binom{k-j+2}{2}-1=N_k,
 \qquad \chi_D+\sum_{j=0}^{m-3}\binom{k-j+2}{3}=\chi_k.
\]
Indeed the first identity is the differentiation-kernel decomposition in that proposition, and the second follows by summing Pascal's identity. Each complete binary parent block contributes $(\det M)^{\binom{k-j+2}{3}}$. Multiplying these factors by the degree-$D$ coefficient Jacobian proves~\eqref{surf:parent-jacobian}; the fixed differentiation determinant is absorbed into the stated nonzero constant. The same argument gives the weighted alternative. The stated nonzero conditions, together with the nonzero normal pivot and invertible $M$, then give local coefficient coordinates by the inverse-function theorem.
\end{proof}

\subsection{Coefficient integration at fixed frames}\label{atlas:fixed-frame-integration}

\begin{lemma}[Conditional Haar volume at the fixed frame]\label{surf:haar}\label{low:4:sf:lem:haar}
Use the group $\mathcal U$ of~\eqref{geo:group}, writing $c_0$ for its translation coordinate $c$, and let $H_0=\{e=f=0\}$, $H_3=\{d=e=f=0\}$, or $H=\{a=f=0\}$. Let $\mathcal W\subset V_k$ be the original normalized semialgebraic coefficient window of fixed format. Suppose that every $J\in\mathcal W$ satisfies $\|J\|\le P^C$ and $s_{\min}(R_J^{\rm ph}|_{\operatorname{Lie}K})\ge P^{-C}$, where $R_J^{\rm ph}X=[XJ]$ and $K$ is the specified subgroup.

\noindent\textup{(i)} Then, for any fixed polynomial $J_*$,
\[
 \operatorname{vol}_{K}\{h\in K:[J_*\circ h]\in\mathcal W\}\lesssim P^{C_k},\qquad K\in\{H_0,H_3,H\}.
\]

\noindent\textup{(ii)} The same estimate holds on a fixed left coset. The coefficients of $J_*$ and the location of that left coset need not be bounded. For $H$, which is normal, the distinction between left and right cosets is immaterial.
\end{lemma}
\begin{proof}
The orbit-window set is semialgebraic of fixed format, treating the coefficients of $J_*$ as parameters. Lemmas~\ref{geo:conditioning}--\ref{geo:paths} give uniformly finitely many components and fixed-format paths from a chosen point of each component to every other point.

Along such a path $h(t)$, every coefficient of $[J_*\circ h(t)]$ is bounded by $P^C$. Each has a bounded number of monotonicity intervals, hence total variation at most $P^{C_k}$. If $v=h^{-1}h'$ is its body velocity, differentiation of the complete coefficient vector and the assumed subgroup pivot give
\[
 \int\|v(t)\|\,dt\le P^{C_k}\sum_{\alpha}\operatorname{Var}\bigl([J_*\circ h(t)]_\alpha\bigr)\le P^{C_k}.
\]
The affine group is represented by upper-unipotent $4\times4$ matrices. Its integral equation contains only iterated integrals of orders one, two, and three. Therefore all group coordinates relative to the chosen starting point are bounded by a fixed polynomial in $\int\|v\|$, hence by $P^{C_k}$. Each component is contained in a left translate of a subgroup coordinate box of that size. Haar measure is Lebesgue measure in these triangular coordinates, and translations preserve it. Sum the finitely many boxes. On a left coset $g_0K$, write $h=g_0k$ and absorb $g_0$ into the fixed polynomial $J_*\circ g_0$. 
\end{proof}

\begin{definition}[Slices and point factors]\label{def:point-factor}
A normalization slice retains every coefficient not fixed by its stated normalization equations. A conditional Haar fiber is the subgroup or coset integral left after fixing the indicated slice and endpoint variables. In those endpoint variables, the point factor is the candidate-area to good-area ratio $A_{\rm cand}/a_0$, where the candidate endpoint area is at most $A_{\rm cand}$ and the common good endpoint set $G_F$ satisfies $|G_F|\ge a_0>0$.
\end{definition}
\begin{lemma}[Complete coefficient--point disintegration]\label{surf:null-density}

\leavevmode\noindent\textup{(i)} On an order-$m$ null sheet, impose the four coefficient constraints
\[
J_{x^{m-1}}(0)=J_{x^{m-1}y}(0)=J_{x^{m-1}z}(0)=J_{x^{m-1}zz}(0)=0.
\]
Write $J=[F\circ h]$, $\kappa=J_{x^m}(0)$, and $B=c_{yz}(0)\ne0$ for its normalized graph. With all remaining $N_k-4$ coefficients as slice coordinates, the joint incidence density is
\begin{equation}\label{surf:null-density-eq}
dF\,de\,df=C_{k,m}\,2|\kappa|^4|B|\,dJ_{N_k-4}\,dh_{H_0}\,d\eta\,df,\qquad \eta=e-df.
\end{equation}

\noindent\textup{(ii)} In the Gauss chart, whose constraints replace the last row by $J_{x^{m-1}yz}(0)=0$, the corresponding density at a fixed Gauss point is
\begin{equation}\label{surf:gauss-density}
dF=C'_{k,m}|\kappa|^4|c_{yy}(0)|^2\,dJ_{N_k-4}\,dh_H.
\end{equation}
Here $dJ_{N_k-4}$ is Lebesgue measure in the remaining monomial coefficient coordinates, and $dh_{H_0},dh_H$ are Haar measures in the group coordinates of Lemma~\ref{surf:haar}. The constants are positive fixed Taylor-coordinate factors.

\noindent\textup{(iii)} Under the bounds on every retained coefficient in Definition~\ref{def:admissibility} and fixed-frame subgroup-pivot bounds of Lemma~\ref{geo:original-subgroup-pivots}, a fixed full slice output map $\mathcal T$ with $|\det D\mathcal T|\ge\delta>0$ and a uniformly bounded number of preimages gives
\[
\operatorname{vol}\mathcal E\le P^{C_k}\delta^{-1}V_{\rm out}\,A_{\rm cand}/a_0,
\]
where $V_{\rm out}$ is its output-box volume and $A_{\rm cand}/a_0$ is the point factor of Definition~\ref{def:point-factor}. For the $H_3$ alternative, include also the length of the separately integrated endpoint-shear interval.
\end{lemma}
\begin{proof}
\par\medskip\noindent\textit{Joint null density.}\par\nobreak\smallskip

Use the infinitesimal columns $T_x=\partial_x$, $A=y\partial_x$, $C=z\partial_x$, $D=z\partial_y$, and $T_y=\partial_y$. In the null constraints the columns $(T_x,A,C,D)$ have successive pivots $\kappa,\kappa,\kappa,2J_{x^{m-1}yz}$. Differentiating the graph equation gives $J_{x^{m-1}yz}=-\kappa B$. Elimination by the first three columns therefore gives absolute determinant $2|\kappa|^4|B|$. All other coefficient coordinates have an identity derivative. Pullback by $h$ on the full filtered coefficient space has determinant one, since its action on every homogeneous quotient is unipotent. Keeping the two physical center variables gives the full $(N_k+2)$-variable density before the last change of variables.

\par\medskip\noindent\textit{Fixing $(\eta,f)$.}\par\nobreak\smallskip

The substitution $e=\eta+df$ has determinant one in the complete set of integration variables. At fixed $(\eta,f)$ factor $h=k g_0$, where $k\in H_0$ and $g_0(x,y,z)=(x,y+\eta,z+f)$. The coordinates of $k$ are $(a,b,c_0-a\eta-bf,d)$, a triangular change with determinant one. Thus the fixed frame restriction is imposed on $(J\circ g_0^{-1})\circ k^{-1}$ on a single $H_0$ fiber. Fixing the full slice outputs fixes one of boundedly many possible $J$ by Lemma~\ref{geo:coefficient-map}. Lemma~\ref{surf:haar}, with the fixed frame pivot, bounds the remaining conditional Haar integral by $P^{C_k}$. Here $\eta$ is a coordinate in the joint change of variables. The endpoint map $(e,f)\mapsto(\eta,f)$ at a fixed polynomial is not required to have a nonzero Jacobian.

\par\medskip\noindent\textit{Gauss fibers at fixed $(a,f)$.}\par\nobreak\smallskip

For the Gauss constraints use columns $(T_x,T_y,C,D)$. Their triangular diagonal is $(\kappa,-\kappa c_{yy},\kappa,-\kappa c_{yy})$, giving \eqref{surf:gauss-density}. The map $h\mapsto(a,f)$ is a homomorphism, and $H$ is its normal kernel. Fixing the Gauss point fixes its coset. The same coefficient area formula and fixed-frame Haar bound apply.

\par\medskip\noindent\textit{Keeping the endpoint shear: fix $(d,\eta,f)$.}\par\nobreak\smallskip

If an output map retains the actual endpoint shear as an external variable, keep the four null constraints and the factor $2|\kappa|^4|B|$ in~\eqref{surf:null-density-eq}. The shear coordinate $d$ is additive under multiplication in $H_0$, and its zero set is $H_3$. Hence
\[
dh_{H_0}=dh_{H_3}\,dd.
\]
Fix the slice outputs and $(d,\eta,f)$, equivalently $(d,e,f)$. The remaining group integral is on an $H_3$ coset. Its three normalization columns have determinant $\kappa^3$, and Lemma~\ref{surf:haar}, with the fixed-frame pivot, bounds the conditional Haar integral by $P^{C_k}$. Integrating this bound over the stated interval of $d$ completes the alternative disintegration. The $H_0$ disintegration already includes this variable in its fiber.

\par\medskip\noindent\textit{Integration over good endpoints.}\par\nobreak\smallskip

Finally, let $\mathcal E$ be the original coefficient set and $G_F$ its common good endpoint set, with $|G_F|\ge a_0$. Write $V_{\mathrm{out}}$ for the full slice-output box volume and $A_{\mathrm{cand}}$ for the candidate endpoint area in the joint variables of the chosen disintegration. For the $H_0$ and Gauss disintegrations, Tonelli's theorem and the preceding bounds give
\[
a_0\operatorname{vol}(\mathcal E)\le\int_{\mathcal E}|G_F|\,dF\le P^{C_k}\delta^{-1}V_{\mathrm{out}}A_{\mathrm{cand}}.
\]
Division by $a_0$ gives the asserted point factor. For the $H_3$ alternative, the right side also includes the length of the separately integrated endpoint-shear interval, as above. The semialgebraic branch counts are bounded by the fixed format, so they only change $P^{C_k}$.
\end{proof}

\subsection{Finite slope estimates in degrees six through nine}\label{atlas:slopes}\label{low:78:s38det:section}\label{low:9:h9d:slope-section}\label{low:78:h3n:det7-section}

For the main theorem of Paper I, the finite case needed is $D=9$ at $(k,m)=(14,7)$; the other cases are retained for the shared estimates used by Paper II.

\begin{lemma}[Finite sampling near polynomial roots]\label{atlas:finite-slope-grid}
Fix a nonzero real polynomial $p$ with $p(0)\ne0$, and an integer $\mu\ge1$ bounding the multiplicity of every real root of $p$.

\noindent\textup{(i)} For fixed $a>0$ and $C_0>0$, there is a finite equally spaced set $\mathcal T\subset[0,a]$, containing zero, and $c>0$ such that
\[
H\ge1,\quad |x|\le C_0H|\delta|\quad\Longrightarrow\quad \max_{t\in\mathcal T}|p(x+t\delta)|\ge cH^{-\mu}.
\]

\noindent\textup{(ii)} For fixed $p$, if $a$, $a^{-1}$ and $C_0$ are bounded by powers of $P$, the constant $c^{-1}$ can also be bounded by a power of $P$; the grid cardinality depends only on $p$.
\end{lemma}
\begin{proof}
Choose disjoint fixed small neighborhoods of the finitely many real roots, separated from zero. Off these neighborhoods, $|p|$ has a positive lower bound, also for large $|x|$, so the node $t=0$ suffices. Otherwise $|x|\ge c_p>0$, and $|\delta|\ge c_p/(C_0H)$. Choose more than $\deg p+1$ equally spaced nodes. If $a|\delta|$ is small relative to the root separation, all samples near zeros can be near only the one root close to $x$; one of the two endpoint nodes has distance at least $a|\delta|/2$ from it. The local root factorization gives a lower bound by a constant times $\min(1,(a|\delta|)^\mu)$. If $a|\delta|$ is larger, the sample spacing has a fixed positive lower bound. Choose smaller root neighborhoods with diameters less than that bound. Each contains at most one sample, so the number of nodes forces one sample outside all of them, again giving a fixed lower bound. These estimates prove the assertion; their dependence on $a$ and $C_0$ involves only fixed powers.
\end{proof}

\begin{lemma}[Anisotropic selection for a fixed polynomial]\label{surf:slope-selection}
Let $p$ be a fixed nonzero real polynomial with $p(0)\ne0$ and $\deg p\le d$. On a wide patch with transverse length $w\ge H\ge1$, suppose the normalized quadratic graph is $(r/2)y^2+Byz$, with $B$ and $B^{-1}$ conditioned, and the actual relative slope satisfies $|e|\le P^Cw$. There is a fixed finite set of perturbations $e+t_jw/H$, with fixed $|t_j|\le\varepsilon_P$, so that each change $\Delta e=t_jw/H$ satisfies $|\Delta e|H/w\le\varepsilon_P$, such that at one of them the slope determinant factor $|p(r(e+t_jw/H)/B)|$ is bounded below by $P^{-C}H^{-d}$.
\end{lemma}

\begin{proof}
If $p$ is constant the assertion is immediate. Otherwise apply Lemma~\ref{atlas:finite-slope-grid} with $x=re/B$, $\delta=rw/(BH)$, and multiplicity bound $d\ge1$. The slope hypothesis gives $|x|\le P^CH|\delta|$. The case $r=0$ follows from $p(0)\ne0$, and the permitted $t_j$ lie in an arbitrarily small fixed interval. Thus the shifts have bounded anisotropic size and require only a fixed number of choices.
\end{proof}

Use $G_D$, $g_D$ and $t_D$ from Lemma~\ref{surf:strong-atlas}, in the Taylor-coordinate convention of~\eqref{surf:ordinary-list}. For every $D\ge6$, define
\begin{equation}\label{atlas:slope-polynomial}
p_D(s)=(1+s)^{t_D}g_D\!\left(\frac{s(s+2)}{4(1+s)^2}\right).
\end{equation}
This is a polynomial of degree at most $t_D$, since it is the homogeneous polynomial $G_D$ evaluated on the sheared quadratic coefficients. Also $p_D(0)=g_D(0)\ne0$. For $6\le D\le9$, put $\mu_D=t_D-2\deg g_D$. The exact determinant data needed for finite slope selection are
\begin{equation}\label{atlas:slope-data}
\begin{array}{c|rrr}
D&t_D&\deg g_D&\mu_D\\\hline
6&73&35&3\\
7&121&58&5\\
8&189&92&5\\
9&282&137&8.
\end{array}
\end{equation}
The polynomials are nonzero at zero. Every complex root has multiplicity at most two, except that $g_8$ has a triple root at $t=1$. Supplementary Section~S7 specifies the common matrix entries, row and column orders, exact factors, and squarefreeness certificates.

\begin{lemma}[Finite-degree slope selection]\label{atlas:finite-slopes}\label{low:6:wc36:slope-repair}\label{low:78:s37i:slope-repair}\label{low:78:h3n:slope-repair}\label{low:slope7}\label{low:78:s38det:multiplicity}\label{low:slope8}\label{low:9:h9d:slope}\label{low:slope9}

\leavevmode\noindent\textup{(i)} For $6\le D\le9$, the polynomial $p_D$ of~\eqref{atlas:slope-polynomial} has a zero of exact order $\mu_D$ at $s=-1$ and all its other real zeros have multiplicity at most two.

\noindent\textup{(ii)} Suppose $w\ge H\ge1$, $|r|+|B|\le P^C$, $|B|\ge P^{-C}$, and $|e|\le P^Cw$. For every sufficiently small fixed $a_P>0$ bounded below by an inverse power of $P$, there is a fixed finite set $\mathcal T\subset[0,a_P]$, independent of $w,H,r,B,e$, such that one $t\in\mathcal T$, with $E=e+tw/H$, satisfies
\[
\left|G_D\left(\frac r2,B+rE,BE+\frac r2E^2\right)\right|\ge P^{-C'}H^{-\mu_D}.
\]
The number of choices is bounded in terms of $D$. No lower bound for $r$ is required.
\end{lemma}
\begin{proof}
The exact factors and Euclidean greatest-common-divisor checks in supplementary Section~S7 give~\eqref{atlas:slope-data} and the stated root multiplicities of $g_D$. For real $s\ne-1$,
\[
\frac{s(s+2)}{4(1+s)^2}=\frac14-\frac1{4(1+s)^2}<\frac14,\qquad
\frac{d}{ds}\frac{s(s+2)}{4(1+s)^2}=\frac1{2(1+s)^3}\ne0.
\]
Thus the substitution preserves finite root multiplicity and excludes the triple root of $g_8$. Since $t_D\ge2\deg g_D$, $p_D$ is a polynomial. Its leading $g_D$ term gives exact order $t_D-2\deg g_D=\mu_D$ at $s=-1$; all lower terms vanish to higher order. Also $p_D(0)=g_D(0)\ne0$.

Put $x=re/B$ and $\delta=rw/(BH)$. The assumed slope range gives $|x|\le P^CH|\delta|$, and the determinant in the statement equals $B^{t_D}p_D(x+t\delta)$. Lemma~\ref{atlas:finite-slope-grid}, with the multiplicity bound $\mu_D$, and the conditioned factor $|B|^{t_D}$ prove the estimate. When $r=0$ the value is the nonzero constant $B^{t_D}g_D(0)$. The tolerated slope changes satisfy $|\Delta e|H/w\le\varepsilon_P$ for the chosen small conditioning constant and are made after rounding, so they add no seed grid.
\end{proof}

For a parent phase of degree $k=9$, the cubic-normal ($m=3$) and quartic-normal ($m=4$) cases use $D=8$ and $D=7$, respectively, so both have cost $P^CH^5$.


\section{Surface coefficient estimates in degrees at least ten}\label{surf:section}

This section estimates the coefficient volumes for the contained null patches and Gauss patches supplied by the common geometric assignment, then sums these estimates. The resulting surface estimate holds for $k\ge10$, including the degrees ten through thirteen used in Paper II.

\subsection{Null coefficient bounds by normal order}\label{surf:null-counts}

Fix $k\ge7$, a normal order $2\le m\le\lfloor k/2\rfloor$ supplied by the assignment, and $D=k-m+2$. Use the normal restrictions of Definition~\ref{def:normal-restrictions} and geometric certificates of Definition~\ref{def:geometric-certificate}, with a fixed positive proportion of common good points and conditioned $\kappa_m,B$. Prepare the scalar, curvature and seed lists simultaneously. Use the bounds stated after Proposition~\ref{low:78:smult:atlas} for the values of the positive normal scalars and the tangential derivatives of the parent value $\phi_0$. Also retain $|r|\le P^C$. No lower normal scalar is a denominator. In the width calculations below, the parent scalar is $\phi_0$, whose tangential gradient is bounded. For $m\ge3$, the critical value $\phi_{m-2}$ of $G$ belongs to the bounded positive normal scalars.

Use the index sets and finite-degree conventions of Section~\ref{atlas:restriction}. The coefficient measure is the full $N_k$-dimensional parent measure. Restore the complete lower graph jet before deleting the four null graph-normalization rows with indices $(0,0),(1,0),(0,1),(0,2)$. Since $T_{D-2}\subset E_D,E_D^{\rm w}$, these rows are present in every atlas. Lemma~\ref{surf:null-density} gives transverse determinant $2\kappa_m^4B$, the remaining $N_k-4$ slice directions, and the fixed-frame conditional Haar integral. Its three-column alternative has determinant $\kappa_m^3$ and requires a separate integration over the actual endpoint shear interval. These are alternative disintegrations. Fix the rounded external seed before coefficient integration as in Definition~\ref{def:coefficient-atlas}; Lemma~\ref{geo:coefficient-map} bounds the output map's inverse multiplicity. Use the point factor of Definition~\ref{def:point-factor}.

\begin{definition}[Derivative widths]\label{def:derivative-widths}
For vertically long ribbons and Gauss strips, write the long derivative index first. An output width is the length of the allowed interval for the indicated derivative. On coordinate scales $L,W$, the scale factors for bounded scalars are $L^{-i}W^{-j}$. For a scalar with bounded gradient, including the parent value $\phi_0$, these factors are $L^{1-i}$ for $j=0,i>0$ and $L^{-i}W^{1-j}$ for $j>0$. A purely longitudinal row has $j=0$. These conventions specify the scale factors; the prepared bounds include the stated amplitude and any additional coframe factor.
\end{definition}

For the ordinary determinant of Lemma~\ref{surf:strong-atlas}, use $p_D$ from~\eqref{atlas:slope-polynomial} and the slope estimates of Subsection~\ref{atlas:slopes}. We use the following inverse-determinant exponents:
\begin{equation}\label{surf:mu}
 \mu_D=\begin{cases}3,&D=6,\\5,&D=7,8,\\8,&D=9,\\t_D,&D\ge10.\end{cases}
\end{equation}
The finite-degree values follow from Lemma~\ref{atlas:finite-slopes}; the last follows from Lemma~\ref{surf:slope-selection}. For $k\ge10$, every required derivative degree obeys $D=k-m+2\ge\lceil k/2\rceil+2\ge7$, so these alternatives cover all normal orders. For $k\ge14$, $D\ge9$, and $D=9$ occurs only when $(k,m)=(14,7)$; this case uses the degree-nine estimate in Lemma~\ref{atlas:finite-slopes}.

Here are the index sums used in the volume estimates:
\[
 \mathsf{s}(d)=\frac{d(d+1)^2}{2},\qquad \mathsf{t}(d)=\binom{d+2}{3},
\]
\[
 \mathsf{h}(D)=\mathsf{s}(D)+\mathsf{s}(D-3)+\sum_{r=3}^{\lfloor D/2\rfloor}\mathsf{s}(D-2r),\qquad \mathsf{v}(D)=\mathsf{t}(D)+\mathsf{t}(D-3)+\sum_{r=3}^{\lfloor D/2\rfloor}\mathsf{t}(D-2r).
\]
Put $U=\mathsf{h}(D)+\chi_k-\chi_D$. The index sums are
\[
\begin{aligned}
\sum_{(i,j)\in I_d}(i,j)&=(\mathsf{s}(d),\mathsf{s}(d)),\\
\sum_{(i,j)\in I_d^{\mathrm w}}(i,j)&=(\mathsf{s}(d)+\mathsf{t}(d),\mathsf{s}(d)-\mathsf{t}(d)),\\
\sum_{(i,j)\in T_d}(i,j)&=(\mathsf{t}(d),\mathsf{t}(d)).
\end{aligned}
\]
Thus all retained scalar blocks have index sums $(U,U)$ in the ordinary chart and $(U+\mathsf{v}(D),U-\mathsf{v}(D))$ in the weighted chart. The lower parent triangles account for exactly $\chi_k-\chi_D$ in each index. For later use define
\[
c_{k,m}=\begin{cases}k(k+1),&m=2,\\k(k+1)/2,&m\ge3,\end{cases}\quad p_{k,m}=\begin{cases}2k,&m=2,\\k,&m\ge3,\end{cases}\quad c_{k,m}^{\mathrm w}=\begin{cases}k^2,&m=2,\\k(k+1)/2,&m\ge3.\end{cases}
\]
Here $c_{k,m}$ counts ordinary rows of the parent scalar of positive short index, $p_{k,m}$ counts weighted rows of the parent scalar of zero short index and positive long index, and $c_{k,m}^{\mathrm w}$ counts weighted rows of the parent scalar of positive short index. The graph index sums are
\begin{align*}
a_D&=\sum_{(i,j)\in E_D,\ i\ge2}(i-2)=\frac{(D-3)(D-4)^2}{2},\\
b_D&=\sum_{(i,j)\in E_D,\ i\ge2}j=\frac{(D-4)(D-3)(D-2)}{2},\\
e_D&=|E_D\setminus T_{D-2}|=\frac{(D-3)(D-4)}2,\\
x_D&=\sum_{E_D\setminus T_{D-2}}i=\frac{(D-2)(D-3)(D-4)}3,\\
 \sigma_D&=\sum_{E_D^{\mathrm w}\setminus T_{D-2}}i=\mathsf{s}(D-3)+\mathsf{t}(D-3)+1-\mathsf{t}(D-2),\\
\sum_{(i,j)\in E_D\setminus T_{D-2}}(j-2)&=x_D-2e_D.
\end{align*}
All extra ordinary graph rows have both indices at least two.

\paragraph{Wide ribbons}

\begin{proposition}[Wide ribbon volume]\label{surf:wide}
Let $k\ge10$. Suppose a contained ribbon has parametrization $\Psi(\xi,\tau)=(\xi+\Gamma(\tau),\tau)$ in $(y,z)$ coordinates, with lengths $w\ge H\ge1$. Assume $|\Gamma^\prime(\tau)-d(\xi+\Gamma(\tau),\tau)|\le P^Cw$ and $|\Gamma^{(j)}|\le P^CD_*H^{1-j}$ for the prescribed finite orders $j\ge2$, with $D_*\le P^CwH$, and the fixed-frame endpoint area factor is at most $P^CH^2$.

\noindent\textup{(i)} Under the complete parent hypotheses of Lemma~\ref{geo:moderate}, the coefficient set satisfies
\begin{equation}\label{surf:wide-eq}
 \operatorname{vol}E_j(w,H)\le P^{C_k}w^{-A_{k,m}}H^{-B_{k,m}},
\end{equation}
where
\[
A_{k,m}=U-c_{k,m}-k+a_D,\qquad B_{k,m}=U+b_D-3k-1-\mu_D.
\]

\noindent\textup{(ii)} At fixed coordinate area $S=wH$, the dyadic sum over $1\le H\le S^{1/2}$, with $w=S/H$, is at most
\begin{equation}\label{surf:wide-sum}
P^{C_k}(1+\log S)S^{-\omega_{k,m}},\qquad \omega_{k,m}=\min\left(A_{k,m},\frac{A_{k,m}+B_{k,m}}2\right).
\end{equation}
\end{proposition}
\begin{proof}
A bounded scalar row $(i,j)$ has width $P^Cw^{-i}H^{-j}$. A row of the parent scalar with $i>0$ has width $P^Cw^{1-i}H^{-j}$; a row with $i=0$, $j>0$ has width $P^CwH^{1-j}$. Therefore the complete scalar product is
\[
P^{C_k}w^{-U+c_{k,m}+k}H^{-U+k}.
\]
Keep original graph coordinates for rows with horizontal index at most one, and use curved coordinates for rows with horizontal index at least two. The identity $\partial_\xi^2(c\circ\Psi)=c_{yy}\circ\Psi$ converts their prepared bounded-curvature derivatives into widths $w^{2-i}H^{-j}$. Their product is $w^{-a_D}H^{-b_D}$. This partial output change has determinant one. Indeed all original rows being restored lie in the complete lower graph jet $T_{D-2}$; inverse shear is triangular, with diagonal one, when ordered by horizontal index. Nonlinear seed terms only add lower-total-order rows. Extra graph rows stay in curved coordinates. Each graph row uses exactly one of these two bounds: its original-coordinate interval or its prepared curvature interval.

Round the seed derivative of order $j$ on a mesh $c_Pw/H^j$, for $1\le j\le k$. The first derivative costs $P^CH$, and each later derivative costs at most $P^CH^2$. The total is $P^{C_k}H^{2k-1}$. To justify retention of all derivative widths, prepare the full scalar and curvature derivatives through every total order appearing in the finite chain rule before rounding. Interpolate the finitely many seed jets. The induced coordinate change is a transverse displacement $\delta\Gamma(\tau)$ with $|\partial_\tau^j\delta\Gamma|\le c_PwH^{-j}$ and zero value at the centered endpoint. The Leibniz rule shows that the bounded-scalar norm with weights $w^iH^j$ is preserved. For the parent scalar, use weights $w^{i-1}H^j$ when $i>0$, and $w^{-1}H^{j-1}$ when $i=0,j>0$; every transverse differentiation of the scalar supplies the inverse factor needed to absorb $\delta\Gamma$. These are finite-dimensional triangular jet equations, so their norms grow by at most a fixed factor on the unit interpolation interval. The same argument applies to the bounded scalar $c_{yy}$ for the curvature rows. Thus one fixed rounded seed jet retains the stated widths on the common good set.

Apply the finite slope selection after rounding. Its bounded anisotropic shifts preserve these widths and its determinant loss is $P^{C_k}H^{\mu_D}$. The external linear parts have determinant of absolute value one, so the explicit character in \eqref{surf:parent-jacobian} introduces no additional length. All output maps are now fixed before coefficient integration. Lemma~\ref{surf:null-density} contributes the point factor $P^CH^2$. Combining these five factors gives
\[
\begin{aligned}
\operatorname{vol}E_j(w,H)&\le P^{C_k}(w^{-U+c_{k,m}+k}H^{-U+k})(w^{-a_D}H^{-b_D})H^{2k-1}H^{\mu_D}H^2\\
&=P^{C_k}w^{-A_{k,m}}H^{-B_{k,m}}.
\end{aligned}
\]

For the final sum write $w=S/H$ with $1\le H\le S^{1/2}$. Then $w^{-A}H^{-B}=S^{-A}H^{A-B}$, whose dyadic geometric sum gives \eqref{surf:wide-sum}, with a logarithmic factor when $A=B$.
\end{proof}

\paragraph{Vertically long ribbons}

\begin{proposition}[Vertically long ribbon volume]\label{surf:vertical}
Let $k\ge10$. Suppose the actual ribbon has parametrization $\Psi(\xi,\tau)=(\tau+\Gamma(\xi),\xi)$ in $(y,z)$ coordinates, where the longitudinal variable $\xi$ has length $L$, the transverse variable $\tau$ has length $W$, and $L\ge W\ge1$. Put $S=LW$ and $K_0=L/W$. Assume $|\Gamma^\prime(\xi)-d(\tau+\Gamma(\xi),\xi)|\le P^CL$, $|\Gamma^{(j)}|\le P^CL^{3-j}$ for the prescribed finite orders $j\ge2$, and a fixed-frame point factor at most $P^CL^2K_0$. Let $E_j^{\mathrm g}$ and $E_j^{\mathrm w}$ be the coefficient subsets using the ordinary and weighted charts, respectively.

\noindent\textup{(i)} They satisfy
\begin{equation}\label{surf:vertical-eq}
\begin{aligned}
\operatorname{vol}E_j^{\mathrm g}&\le P^{C_k}S^{-(g_L+g_W)/2}K_0^{-(g_L-g_W)/2},\\
\operatorname{vol}E_j^{\mathrm w}&\le P^{C_k}S^{-(w_L+w_W)/2}K_0^{-(w_L-w_W)/2}.
\end{aligned}
\end{equation}
where
\begin{align*}
g_L&=U+x_D-5k-2,&g_W&=U-c_{k,m}+x_D-2e_D+k+1,\\
w_L&=U+\mathsf{v}(D)-2p_{k,m}-6D-2,&w_W&=U-\mathsf{v}(D)-c_{k,m}^{\mathrm w}+2D+1.
\end{align*}

\noindent\textup{(ii)} The union $E_j^{\mathrm g}\cup E_j^{\mathrm w}$ covers the relevant coefficient set, since the pair of charts covers every actual first slope.
\end{proposition}
\begin{proof}
Use the widths of Definition~\ref{def:derivative-widths}, before including the purely longitudinal amplitude. In the ordinary chart there are $k$ purely longitudinal rows, so the scalar product is $L^{-U+k}W^{-U+c_{k,m}}$. The interval width for each purely longitudinal row is then multiplied by its amplitude bound $1+|e|\le P^CL$.

Retain the complete original graph window $T_{D-2}$ and use curvature widths only on $E_D\setminus T_{D-2}$. Since the short direction is the original $y$ direction, the second short derivative of the curved graph is the bounded scalar $c_{yy}$. The extra product is $L^{-x_D}W^{-x_D+2e_D}$. Restoring the complete lower jet is triangular with determinant of absolute value one. As in the wide case, this makes the ordinary lower intervals and the additional curvature intervals distinct output coordinates.

The seed meshes are $c_PW/L^j$. The first range $P^CL$ costs $P^CLK_0$, and each higher range $P^CL^{3-j}$ costs $P^CL^2K_0$. For $k$ jets the cost is $P^{C_k}L^{2k-1}K_0^k$. Use the finite jet interpolation proved in Proposition~\ref{surf:wide}, with transverse displacement derivatives bounded by $c_PWL^{-j}$ and the widths of Definition~\ref{def:derivative-widths}. Prepare the full auxiliary scalar and curvature derivative lists jointly before rounding, as in that proof, and retain the purely longitudinal amplitude of the parent scalar. For the ordinary subset $E_j^{\mathrm g}$, using $K_0=L/W$, we obtain
\[
\begin{aligned}
\operatorname{vol}E_j^{\mathrm g}&\le P^{C_k}(L^{-U+k}W^{-U+c_{k,m}})L^k(L^{-x_D}W^{-x_D+2e_D})L^{2k-1}K_0^kL^2K_0\\
&=P^{C_k}L^{-g_L}W^{-g_W}.
\end{aligned}
\]

To choose the determinant chart, write the quadratic graph before the relative shear as $(r/2)y^2+Byz$, with $B$ conditioned. In the long-first coordinates the coefficients are
\[
A=eB+re^2/2,\qquad b=B+re,\qquad C=r/2,\qquad b^2-4AC=B^2.
\]
Choose a sufficiently small conditioned constant $\varepsilon$. Either $|A|\ge\varepsilon(1+|e|)$, or the last identity and the formula for $A$ imply $|b|\ge P^{-C}$ and $|AC/b^2|$ is as small as required. To see the latter assertion for $|e|\ge1$, put $s=re/B$. Then $|s+2|\lesssim_P\varepsilon$, so $|1+s|$ is bounded below and $|s(s+2)/(1+s)^2|\lesssim_P\varepsilon$. For $|e|\le1$, boundedness of $C$ and $b^2=B^2+4AC$ gives the assertion directly. The second alternative uses the ordinary determinant because $g_D(0)\ne0$; the first uses the weighted determinant.

In the weighted chart the raw scalar exponent pair is $(U+\mathsf{v}(D)-p_{k,m},U-\mathsf{v}(D)-c_{k,m}^{\mathrm w})$. Retain the full purely longitudinal amplitude $L^{p_{k,m}}$. The greatest longitudinal index is $J=2D$, and $2D\ge k$, so every lower parent triangle is included. The seed cost is $L^{2J-1}K_0^J$. Remove nonlinear seed jets from the graph coordinates by triangular operations: a term of order $a\ge2$ replaces $(i,j)$ by $(i-a,j+1)$, which stays in the retained weighted set because $i-a+2(j+1)\le i+2j$. After this removal the extra graph intervals have product at most $P^{C_k}(1+|e|)^{\sigma_D}$. The inverse determinant contributes $|A|^{-t_D}$. Since $|A|\ge\varepsilon(1+|e|)$ on this branch and $t_D\ge\sigma_D$ as proved below, $(1+|e|)^{\sigma_D}|A|^{-t_D}\lesssim_P(1+|e|)^{\sigma_D-t_D}\le1$. Multiplying by the seed net and $L^2K_0$ gives $w_L,w_W$.
\end{proof}

\paragraph{Parabolic collars}

The collar linear parts lie in the prepared small neighborhood of the triangular null-model matrix from Lemmas~\ref{geo:collar-endpoints} and~\ref{geo:collar-jets}. At that model the quadratic graph has zero purely transverse coefficient and conditioned mixed coefficient $B$, so $g_D(0)\ne0$ gives the grid determinant lower bound; the small error in the linear part preserves it. Higher external jets do not enter this determinant.

\begin{lemma}[A net for collar parameters]\label{surf:collar-net}
Let an actual fixed-width parabolic collar of long length $L\ge1$ have the endpoint parametrization
\[
 \Psi(s,t)=(A(s)+D(s)t+V(s)t^2/2,\ B(s)+t),\qquad A(0)=B(0)=0.
\]
Assume the bounds on the original phase and the prepared derivative estimates of Lemma~\ref{geo:collar-jets}. Use its raw seed derivatives through order $2k+1$ and full scalar jets through total order $2k$; the seed bounds are~\eqref{geo:collar-seed-widths} with $(B_*,D_*,V_*)=(B,D,V)$. Keep the actual values $D_0,V_0$, bounded by $P^C$. Suppose the reciprocal of
\[
J=A'-DB'+(D'-VB')t+V't^2/2
\]
has the prepared bounds~\eqref{geo:reciprocal-collar-jacobian} through the required finite orders. The reconstruction map can then use a fixed seed net with the following properties.

\noindent\textup{(i)} Its cardinality is at most $P^{C_k}L^{5k-1}$.

\noindent\textup{(ii)} The same net preserves every retained scalar width.
\end{lemma}
\begin{proof}
\par\medskip\noindent\textit{Step 1. Net size.}\par\nobreak\smallskip

Use meshes multiplied by a sufficiently small $c_P$ as follows:
\[
\begin{array}{c|c|c}
\text{jet}&\text{mesh}&\text{total number of choices}\\\hline
A_1&L^{-1}&P^CL\\
A_j,\ 2\le j\le k&L^{-j}&P^{C_k}L^{2(k-1)}\\
B_j,\ 1\le j\le k&L^{-j-1}&P^{C_k}L^{2k}\\
D_j,\ 1\le j\le k&L^{-j}&P^{C_k}L^k\\
D_0,\ V_j,\ 0\le j\le k&1,\ L^{-j}&P^{C_k}.
\end{array}
\]
Their product is $P^{C_k}L^{5k-1}$. A retained ordinary row has longitudinal index at most $k$, including rows $(k,j)$ with $j>0$. Its value needs only the seed jets through longitudinal order $k$.

\par\medskip\noindent\textit{Step 2. Coordinate interpolation.}\par\nobreak\smallskip

For comparison, use the full preparation lists of Lemma~\ref{geo:collar-jets}: scalar order $2k$ and unrounded seed order $2k+1$. Interpolate the gridded jets while leaving those temporary higher seed jets unchanged. The induced parameter vector field is
\[
X^s=J^{-1}\bigl(\delta A+\delta D\,t+\delta V\,t^2/2-(D+Vt)\delta B\bigr),\qquad X^t=\delta B-B'X^s.
\]
It vanishes at the centered endpoint, and the displayed meshes and the Leibniz rule give $|\partial_s^i\partial_t^jX|\le P^Cc_PL^{-i}$. The potentially largest term combines $D_a=O_P(L^{1-a})$ with $\delta B_{i-a}=O_P(c_PL^{-(i-a)-1})$, whose product is $O_P(c_PL^{-i})$. Every other product obeys the same or a stronger bound. The reciprocal-Jacobian estimates persist during interpolation: the potentially adverse term $\delta D_0B_{i+1}$ is $O_P(c_PL^{-i})$, and the finer $B$ mesh treats $D\delta B'$ in the same manner. Differentiate $J^{-1}J=1$ to obtain the higher reciprocal estimates inductively. The zeroth-order Jacobian stays conditioned when $c_P$ is small.

\par\medskip\noindent\textit{Step 3. Scalar widths.}\par\nobreak\smallskip

At the centered endpoint, use the following finite jet norms, with $(i,j)\in\mathbb Z_{\ge0}^2$:
\[
\begin{aligned}
\|f\|_{\rm b}&=\max_{i+j\le2k}L^i|\partial_s^i\partial_t^jf(0)|,\\
\|f\|_{\rm c}&=\max\left\{\max_{1\le i\le2k}L^{i-1}|\partial_s^if(0)|,\ \max_{\substack{j\ge1\\i+j\le2k}}L^i|\partial_s^i\partial_t^jf(0)|\right\}.
\end{aligned}
\]
Use $\|\cdot\|_{\rm b}$ for bounded scalars and $\|\cdot\|_{\rm c}$ for the nonconstant jet of the parent scalar. Choose $c_P$ sufficiently small that the finite system coefficients below are bounded on the unit interpolation interval. Differentiating the interpolation equation $\dot f=X\cdot\nabla f$ gives a finite system with bounded weighted coefficients. Although the atlas rows by themselves are not closed under this equation, the full jet of total order $2k$ is closed: the only possible derivative of order $2k+1$ is multiplied by $X(0)=0$. For the parent scalar $\phi_0$, use the second norm. If all transverse derivatives strike $X^s$, the remaining scalar derivative is purely longitudinal and contains at least one long derivative; its bound larger by one power of the length is exactly canceled by that derivative. An $X^t\partial_t$ term already has a transverse derivative on the scalar. The same finite system argument therefore preserves the widths of the parent scalar. Integration over the unit interpolation interval costs only a fixed conditioning factor.

\par\medskip\noindent\textit{Step 4. The fixed output map.}\par\nobreak\smallskip

After comparison, discard the temporary higher seed jets: no retained output depends on them. Choose a polynomial coordinate map with these Taylor coefficients. Its linear part is invertible with the bounds already retained. Graph coordinates remain in bounded windows because all its endpoint seed jets are bounded when $L\ge1$, and the implicit graph pivot is conditioned. The linear part remains in its conditioned atlas chart. 

At an endpoint away from the original seed, the actual relation is $A'-DB'=1+ut_0+V't_0^2/2$ before reparametrization. The meshes and reciprocal-Jacobian comparison retain this factor.
\end{proof}

\begin{proposition}[Collar volume in every sufficiently large degree]\label{surf:collar-tail}
For $k\ge9$, every normal order $2\le m\le\lfloor k/2\rfloor$ on an actual collar supplied by Lemmas~\ref{geo:collar}, \ref{geo:collar-endpoints}, and~\ref{geo:collar-jets} has
\[
 \operatorname{vol}E_j(L)\le P^{C_k}L^{-\beta_{k,m}},\qquad \beta_{k,m}=\mathsf{h}(D)+\chi_k-\chi_D-6k-2\ge\chi_k+4.
\]
\end{proposition}
\begin{proof}
All positive normal scalars are bounded, whereas the parent scalar uses bounded first derivatives. Thus the raw scalar product is $L^{-U+k}$. The complete graph window is bounded. The rounded first tangential matrix is a small conditioned perturbation of the triangular null model, so the ordinary determinant is conditioned by $g_D(0)\ne0$; higher graph terms do not affect that determinant. This includes cubic graph invariants. The seed net of Lemma~\ref{surf:collar-net} and the fixed-frame point-and-shear bound give $\operatorname{vol}E_j(L)\le P^{C_k}L^{-U+k}L^{5k-1}L^3=P^{C_k}L^{-(U-6k-2)}$. The factor $L^3$ uses the $H_3$ alternative of Lemma~\ref{surf:null-density}, with area factor $L^2$ and endpoint shear range $O_P(L)$.

To check its size, put $E(d)=\mathsf{h}(d)-\mathsf{h}(d-1)-\mathsf{t}(d)$. Direct summation gives
\[
\begin{aligned}
E(2n)&=\frac{4(n-3)^3+9(n-3)^2+53(n-3)+96}{6},\\
E(2n+1)&=\frac{4(n-3)^3+15(n-3)^2+65(n-3)+126}{6}.
\end{aligned}
\]
Thus $E(d)>12$ for $d\ge7$. Increasing $m$ by one decreases $D$ by one and decreases $U$ by $E(D)>0$. The smallest exponent therefore occurs at $D_* =\lceil k/2\rceil+2$. Since $k\le2D_*-4$,
\[
 \beta_{k,m}-(\chi_k+2)\ge \mathsf{h}(D_*)-\chi_{D_*}-12D_*+20.
\]
The right side equals $2$ at $D_*=7$ and increases by $E(d)-12>0$ thereafter. This proves the claimed margin for all $k\ge9$.
\end{proof}

\subsection{Gauss coefficient counts}\label{surf:gauss-counts}

\paragraph{Gauss coordinate maps and coefficient volumes}
The atlas and volume formulas below hold for $k\ge6$ and $2\le m\le\lfloor k/2\rfloor$. Here we use them for $k\ge7$ on the conditioned normal branches of Section~\ref{geo:section}.

Use the complementary chart of Theorem~\ref{geo:gauss}, with conditioned $r_0=c_{yy}$, $s_0=\det(D^2c)/r_0$ and their inverses, and the selected branch $y=Y_F(p,z)$ for $p=c_y$ in~\eqref{geo:gauss-equations}. Retain every regular branch, at most $D_0(D_0-1)$ for $D_0=k-m+1$, and the frame, graph identities and body differential in~\eqref{geo:gauss-frame}--\eqref{geo:gauss-body}. Evaluate graph derivatives at $(Y_F(p,z),z)$ with the phase and branch fixed. Implicit differentiation and the complete parent window bound the prescribed finite derivatives; the theorem supplies the contained strips, full-parent buffer, count and fixed-frame endpoint ranges. Bounded intrinsic paths give polynomial relative-frame bounds as in the proof of Lemma~\ref{surf:haar}. Both Hessian signs are allowed, without a real isotropic direction. The cubic invariant belongs to the small-$c_{yy}$ chart of Section~\ref{surf:null-counts}.

\begin{lemma}[Complete Gauss atlases]\label{surf:gauss-atlas}
Let $k\ge6$ and $2\le m\le\lfloor k/2\rfloor$. Put $d_g=k-2m+2$.

\noindent\textup{(i)} Use the full normal-filtration output list
\[
 (\phi_0)_{I_k\setminus\{(0,0)\}},\qquad (\phi_j)_{I_{k-2j}},\qquad 1\le j\le\lfloor k/2\rfloor,
\]
with the block $j=m-1$ replaced by graph coefficients on $I_{d_g}$. Fix the required Taylor coefficients of the external map from strip parameters to $(p,z)$. For the scalar outputs, compose this map with $(p,z)\mapsto(Y_F(p,z),z)$ and take the indicated Taylor coefficients. During coefficient differentiation the external map stays fixed, but $Y_F$ varies with $F$. The graph coordinates are taken at the physical origin in directions fixed at the chosen reference chart with its stated inverse bounds. The weighted alternative uses $I_d^{\mathrm w}$. These are full $N_k$-dimensional maps.

\noindent\textup{(ii)} In sufficiently small complete-coefficient neighborhoods of conditioned quadratic-profile models, a finite choice of grid or weighted chart has inverse determinant bounded by $P^{C_k}$, including all scalar, graph, and lower parent directions.
\end{lemma}
\begin{proof}
\par\medskip\noindent\textit{Step 1. The model determinant.}\par\nobreak\smallskip

The row count is the identity in Lemma~\ref{surf:normal-filtration}, with only the constant of $\phi_0$ omitted. At a homogeneous quadratic-profile model the positive normal scalars are constant. The fixed normal determinant is a nonzero multiple of $B^{\tau_k}$ in the grid chart near a graph $Q=Byz+Cz^2$, after exchanging the tangential axes if needed; in the weighted chart it is a nonzero multiple of $A^{\tau_k}$. Replacing normal block $m-1$ by graph rows contributes $(-1/\kappa_m)^{(d_g+1)^2}$, because
\[
 \delta c=-\frac{R_{x^{m-1}}|_Q}{\kappa_m}.
\]
For blocks with normal order $r>m-1$, all required graph indices belong to the retained graph set, so row elimination removes the moving-graph term $\phi_{r+1}\delta c$. For a lower block of normal order $r<m-1$ at tangential index $(i,\ell)$, any unremoved correction has input weight $2(m-1)+i+\ell$, strictly above the principal row weight $2r+i+\ell$. Thus these corrections lie in strictly triangular blocks of the homogeneous matrix.

\par\medskip\noindent\textit{Step 2. Variation at a fixed Gauss point.}\par\nobreak\smallskip

Write the model as $F_0=f(x-Q)+\ell_y y+\ell_z z$, with $f^{(m-1)}(0)=0$, and put $S_r=R_{x^r}|_Q$. At a fixed Gauss point,
\[
 \begin{aligned}
 \delta Y&=-\frac{(\delta c)_y}{r_0},\qquad
 \delta\Phi_r=S_r+f^{(r+1)}(0)\delta c\quad(r\ge1),\\
 \delta\Phi_0&=S_0+f'(0)\delta c-\frac{\ell_y}{r_0}(\delta c)_y,
 \end{aligned}
\]
where $\Phi_r=\phi_r(Y_F(p,z),z)$ and the right sides are evaluated along the model Gauss inverse. In particular, the parent moving-graph term is retained even when $m>2$. Step~1 treats the $f^{(r+1)}(0)\delta c$ terms. The remaining parent correction at index $(i,\ell)$ has input weight $2(m-1)+i+\ell+1>i+\ell$, so it is strictly triangular too. Off the normalized slice, graph jets are at the fixed physical origin and scalar rows at the fixed Gauss point; these coincide on the slice. The reference physical directions stay fixed during coefficient differentiation. Their conditioned linear substitution reduces the principal matrix to Lemma~\ref{surf:normal-filtration}.

\par\medskip\noindent\textit{Step 3. Physical normalization rows.}\par\nobreak\smallskip

Since $d_g\ge2$, both retained graph lists contain the complete quadratic jet. Express this jet in the physical $y,z$ derivatives using the fixed linear change from the reference directions. This change of outputs and its inverse are bounded by $P^{C_k}$, so they change the full determinant only by a factor with the same conditioning bounds. The graph equation relates $c(0),c_y(0),c_z(0),c_{yz}(0)$ to the four coefficient constraints in~\eqref{surf:gauss-density} by a triangular change with conditioned nonzero diagonal. These four graph rows have zero derivatives along the normalization slice. The full Jacobian therefore splits into a conditioned four-row constraint block and an invertible $(N_k-4)$-dimensional slice block. Remove the four physical graph rows, with indices $(0,0),(1,0),(0,1),(1,1)$, and keep the transverse group directions in that density. No other coefficient direction is removed.

\par\medskip\noindent\textit{Step 4. Uniform charts.}\par\nobreak\smallskip

A conditioned nonsingular quadratic form has either a conditioned long quadratic coefficient $A$, or a small $A$ and a conditioned mixed coefficient $B$. The first case uses the weighted chart. In the second, the grid determinant is a small perturbation of its nonzero value at $A=0$. The determinant formulas, followed by a polynomially fine finite net of reference directions, give at most $P^{C_k}$ charts with model Jacobian $J_0$ satisfying $|\det J_0|\ge P^{-C_k}$ and $\|J_0^{-1}\|\le P^{C_k}$. The external map into $(p,z)$ stays fixed during coefficient differentiation; $c$ and $Y_F$ do not. Implicit differentiation through auxiliary order $2k+2$ bounds the first and second derivatives of this coefficient map by $P^{C_k}$, with only powers of the retained normal and curvature pivots in the denominators. Choose the complete coefficient neighborhood, reference-direction mesh and rounding constant $c_P$ after the model lower bound, then the large-$L$ threshold, so the actual full Jacobian $J$ satisfies
\[
 \|J_0^{-1}(J-J_0)\|\le\tfrac12,
 \qquad |\det J|\ge2^{-N_k}|\det J_0|\ge P^{-C_k}.
\]
Indeed, higher centerline jets are $O_P(L^{1-j})$ for $j\ge2$, and the rounded errors $c_PW/L^j$ are at most $c_P$ since $1\le W\le L$. All tolerances are inverse powers of $P$, independent of $L,W$; strips below the conditioning threshold are absorbed into $P^{C_k}$.
\end{proof}

For the following calculation define
\[
 \mathsf{a}(k)=\sum_{r=0}^{\lfloor k/2\rfloor}\mathsf{s}(k-2r)=\chi_k+\tau_k.
\]
The identity follows by summing the complete monomial columns, or from the explicit formulas
\[
\begin{array}{c|cc}
k&\mathsf{a}(k)&\tau_k\\\hline
2n&n(n+1)(6n^2+14n+7)/6&n(n+1)^2(n+2)/3\\
2n+1&(n+1)(n+2)(3n^2+7n+3)/3&(n+1)(n+2)(2n^2+6n+3)/6.
\end{array}
\]

\begin{proposition}[Complete Gauss coefficient volume]\label{surf:gauss-tail}
Let an actual Gauss strip have lengths $L\ge W\ge1$, area $S=LW$, and aspect $K=L/W$, with the hypotheses on the original phase of Theorem~\ref{geo:gauss}. Put
\[
G=\mathsf{a}(k)-\mathsf{s}(d_g),\qquad \Delta=\tau_k-\mathsf{t}(d_g),\qquad \beta^{\mathrm G}_{k,m}=G-\frac{k(k+2)}2.
\]
Let $E_j^{\mathrm g}$ and $E_j^{\mathrm w}$ be the coefficient subsets using the grid and weighted Gauss charts, respectively.

\noindent\textup{(i)} Their union covers the coefficient set, and
\[
\begin{aligned}
\operatorname{vol}E_j^{\mathrm g}&\le P^{C_k}S^{-\beta^{\mathrm G}_{k,m}}K^{-\gamma_k^{\mathrm g}},\\
\operatorname{vol}E_j^{\mathrm w}&\le P^{C_k}S^{-\beta^{\mathrm G}_{k,m}}K^{-\gamma_{k,m}^{\mathrm w}},
\end{aligned}
\]
where
\[
 \gamma_k^{\mathrm g}=\frac{k^2}2-k-1,\qquad \gamma_{k,m}^{\mathrm w}=\Delta+\frac{k^2}2-3k-1.
\]

\noindent\textup{(ii)} Both aspect exponents are positive for $k\ge6$. The least area exponent occurs at $m=2$, and is greater than $\chi_k+2$ for $k\ge6$ and greater than $\chi_k+4$ for $k\ge10$.
\end{proposition}
\begin{proof}
Use the widths for bounded scalars and for the parent scalar of Definition~\ref{def:derivative-widths}. The graph rows have bounded joint intervals. Removing their index sum from the complete scalar sum gives the raw grid pair
\[
 (G-k,\ G-k(k+1))
\]
and the raw weighted pair
\[
 (G+\Delta-2k,\ G-\Delta-k^2).
\]
For example, the grid block for the parent scalar has $k$ positive purely longitudinal rows and $k(k+1)$ rows of positive short index. In the weighted block these counts are $2k$ and $k^2$, which explains every amplitude subtraction.

The centerline jet of order $j$ has range $P^CL^{1-j}$ and mesh $c_PW/L^j$, so it costs $P^CK$. The grid output needs $k$ jets and the weighted output needs $2k$ jets. Their total costs are $P^{C_k}K^k$ and $P^{C_k}K^{2k}$. In the anisotropic interpolation argument of Proposition~\ref{surf:vertical}, $q$ displacement factors of total longitudinal order $a$ contribute $W^qL^{-a}$, canceling the derivative factor $L^aW^{-q}$ supplied by the jointly prepared auxiliary jets. The parent widths are preserved as well, using $W\le L$ for purely longitudinal rows. Lemma~\ref{surf:gauss-atlas} supplies the conditioned complete determinant. The candidate endpoint area is $P^CL^2$ and the good area is at least $P^{-C}LW$, so the fixed-frame point factor from \eqref{surf:gauss-density} is $P^CK$. Thus
\[
\begin{aligned}
\operatorname{vol}E_j^{\mathrm g}&\le P^{C_k}L^{-G+k}W^{-G+k(k+1)}K^{k+1},\\
\operatorname{vol}E_j^{\mathrm w}&\le P^{C_k}L^{-G-\Delta+2k}W^{-G+\Delta+k^2}K^{2k+1}.
\end{aligned}
\]
Using
\[
L^{-a}W^{-b}=S^{-(a+b)/2}K^{-(a-b)/2}
\]
gives exactly the asserted exponents. No separate endpoint shear integration occurs in this Gauss disintegration.

Since $\mathsf{s}(d)$ is increasing, $G$ is smallest at $m=2$. Its area margin above $\chi_k+2$ is
\[
 \mathcal G_k=\tau_k-\mathsf{s}(k-2)-\frac{k(k+2)}2-2.
\]
For $k=2n$ and $2n+1$, put $u=n-3$. Substitution of the displayed sums gives respectively
\[
 \mathcal G_{2n}=\frac{u^4+4u^3+5u^2+11u+12}{3},\qquad \mathcal G_{2n+1}=\frac{2u^4+12u^3+25u^2+39u+39}{6}.
\]
These are positive for $n\ge3$. For $k\ge10$ they are at least $34$, so the stronger claimed margin follows. Finally $d_g\le k-2$, and $\tau_k=\mathsf{t}(k)+\mathsf{t}(k-2)+\cdots$ gives $\Delta\ge \mathsf{t}(k)$. Hence both aspect exponents are positive for $k\ge6$, directly from the displayed formulas.
\end{proof}

\subsection{Exponent bounds and the surface estimate}\label{surf:completion}

\paragraph{The surface estimate and summation}

We finish the arithmetic checks for the null ribbon bounds. Put $f_D=\mathsf{h}(D)-\chi_D$. For $m=2$, $k=D$; for $m\ge3$ the possible range is $D+1\le k\le2D-4$. Each of the three area margins in Propositions~\ref{surf:wide} and~\ref{surf:vertical}, after subtraction of $\chi_k+2$, decreases with $k$ when $D$ is fixed. Substituting the largest $k$ in a lifted case yields the lower bounds
\begin{align*}
B_D^{\mathrm{wide}}&=f_D+\frac{a_D+b_D-t_D}{2}-\frac{(2D-4)(2D-3)}4-2(2D-4)-\frac52,\\
B_D^{\mathrm{grid}}&=f_D+x_D-e_D-D^2-\frac D2+\frac52,\\
B_D^{\mathrm{weighted}}&=f_D-D^2-\frac D2-\frac32.
\end{align*}
The first bounds the mean wide exponent. The ordinary mean-wide and grid margins exceed these bounds by $D^2/2-2D-5$, and the ordinary weighted margin exceeds its bound by $D^2/2-7D/2-1$. These differences are positive for $D\ge10$.

Write $D=2n+\varepsilon$, $\varepsilon\in\{0,1\}$, and $u=n-5\ge0$. The separate wide margin $A_{k,m}-(\chi_k+2)$ is at least $B_D^{\mathrm{wide}}$: the table below gives their difference at the largest lifted degree $k=2D-4$ and at the ordinary degree $k=D$, and smaller lifted $k$ improves that difference. This uses no monotonicity in $m$.

The common ordinary-grid and weighted aspect bounds are
\[
\begin{aligned}
\frac{g_L-g_W}{2}&\ge e_D+\frac{(D+1)(D+2)}4-3(D+1)-\frac32,\\
\frac{w_L-w_W}{2}&=\mathsf{v}(D)-p_{k,m}-4D-\frac32+\frac{c_{k,m}^{\mathrm w}}2\ge\mathsf{v}(D)-4D-\frac52.
\end{aligned}
\]
The weighted inequality uses $c_{k,m}^{\mathrm w}/2-p_{k,m}\ge-1$ in both the ordinary and lifted cases. The following table gives coefficients in ascending powers $1,u,u^2,u^3,u^4$. The wide excess is $A_{k,m}-(\chi_k+2)-B_D^{\mathrm{wide}}$ at the indicated degree; the aspect rows are the right-hand lower bounds just displayed.
\begingroup
\small
\setlength{\arraycolsep}{4pt}
\[
\begin{array}{l|c|ccccc}
\text{quantity}&\varepsilon&1&u&u^2&u^3&u^4\\\hline
B_D^{\mathrm{wide}}&0&15&54&155/6&4&1/6\\
&1&49&995/12&385/12&13/3&1/6\\
B_D^{\mathrm{grid}}&0&323/2&436/3&143/3&20/3&1/3\\
&1&247&1189/6&349/6&22/3&1/3\\
B_D^{\mathrm{weighted}}&0&133/2&61&65/3&4&1/3\\
&1&103&515/6&169/6&14/3&1/3\\\hline
\text{Wide excess, lifted}&0&130&98&179/6&4&1/6\\
&1&187&1571/12&433/12&13/3&1/6\\
\text{Wide excess, ordinary}&0&162&124&203/6&4&1/6\\
&1&233&1931/12&481/12&13/3&1/6\\\hline
\text{Grid aspect bound}&0&39/2&37/2&3&0&0\\
&1&59/2&43/2&3&0&0\\
\text{Weighted aspect bound}&0&571/2&631/3&167/3&20/3&1/3\\
&1&811/2&1627/6&397/6&22/3&1/3\\
t_D-\sigma_D&0&216&121&83/3&4&1/3\\
&1&284&911/6&205/6&14/3&1/3.
\end{array}
\]
\endgroup
These identities follow from the finite sums defining $\mathsf h,\mathsf v,t_D$ and the formulas for sums of powers up to degree three. Each row has nonnegative coefficients and a positive constant term. Thus all the required area margins and aspect bounds are positive, and $t_D>\sigma_D$, for $D\ge10$. This justifies the weighted determinant absorption and both coordinate-aspect sums.

When $D=7,8,9$ and $k\ge10$, only lifted cases arise. Their largest possible degrees are $10,12,14$. Direct substitution of \eqref{surf:mu} gives the following margins above $\chi_k+2$ at these limiting degrees:
\[
\begin{array}{c|c|cccc}
D&\max k&\text{wide}&\text{grid}&\text{weighted}&t_D-\sigma_D\\\hline
7&10&17&30&12&85\\
8&12&41&115/2&47/2&119\\
9&14&82&100&41&162.
\end{array}
\]
The grid aspect exponent increases with $k\ge10$ and is at least $2,6,11$ in these three derivative degrees. The weighted lower bound $\mathsf{v}(D)-4D-5/2$ is already positive at $D=7$. Hence all coordinate-aspect sums converge, except for the logarithmic factor explicitly retained in \eqref{surf:wide-sum}. The uniform area exponents furnished by the three ribbon counts are at least $\chi_k+17$, $\chi_k+32$, and $\chi_k+14$, respectively.

\begin{theorem}[Complete surface coefficient sets in every degree $k\ge10$]\label{surf:common-tail}
For the common surface assignment of Theorems~\ref{geo:common-allocation} and~\ref{geo:gauss}, let $S$ and $K$ be its assigned area and geometric aspect.

\noindent\textup{(i)} At each original fixed frame $j$, the coefficient set that it belongs to an assigned surface piece at these scales satisfies, in its complete normalized parent coefficient space,
\begin{equation}\label{surf:common-tail-eq}
 \operatorname{vol}E_j(S,K)\le P^{C_k}S^{-\chi_k-7/2}K^{-1/2}.
\end{equation}
The estimate includes every normal order, both real isotropic charts, the Gauss charts including elliptic models, all rounded Taylor coefficients of the external coordinate and seed maps, and their coordinate-aspect sums.

\noindent\textup{(ii)} If the original physical frame has determinant $v_j$, the corresponding original coefficient measure has the additional exact factor $v_j^{-\chi_k}$.
\end{theorem}
\begin{proof}
Lemma~\ref{geo:moderate}, Theorem~\ref{geo:gauss} and Lemmas~\ref{geo:collar}, \ref{geo:collar-endpoints}, and~\ref{geo:collar-jets} provide precisely the contained patches, common good area, prepared jets, and endpoint ranges used in the preceding counts. The original subgroup pivots are those of Lemma~\ref{geo:original-subgroup-pivots}, so Lemma~\ref{surf:null-density} integrates every coefficient set at the fixed frame.

The coordinate areas in the preceding estimates and the assigned area $S$ are comparable up to powers of $P$, by the geometric assignment. The wide and vertical estimates therefore give area powers strictly stronger than $\chi_k+4$. Their logarithms can therefore be absorbed into that weaker bound. The collar estimate gives power at least $\chi_k+4$ with no coordinate-aspect sum, since its width is bounded above and below by conditioning constants. Bounded caps have bounded area and are absorbed by $P^{C_k}$; a cap with an unbounded horizontal length is counted by its actual fixed-width collar. The Gauss estimate also has area exponent strictly larger than $\chi_k+4$ and a summable positive coordinate-aspect exponent. There are only boundedly many normal orders, sign branches, and algebraic charts at fixed $k$. Combining them gives $P^{C_k}S^{-\chi_k-4}$.

The assigned width is at least $P^{-C}$, so the assigned aspect satisfies $K\le P^CS$. Consequently
\[
S^{-\chi_k-4}\le P^CS^{-\chi_k-7/2}K^{-1/2},
\]
which proves \eqref{surf:common-tail-eq}. By Lemma~\ref{an:character}, returning to the original coefficient variables at the fixed frame multiplies volume by $v_j^{-\chi_k}$.
\end{proof}

\section{Completion of the proof}\label{an:assembly}\label{an:assembly-section}

\subsection{Coefficient incidence}\label{completion:incidence}\label{an:coefficient-assembly}

\begin{proposition}[Coefficient incidence]\label{an:unified-incidence}\label{an:higher-incidence}
Fix $k\ge14$ and $\chi_k+3<p<\chi_k+3+1/16$. Choose $\eta>0$ so that $q=(1+\eta)(p-1)<\min(\beta_k,\chi_k+7/2)$. Put
\[
\gamma=0.
\]
Let $\{T_j\}$ be a fixed family in one physical dyadic shape family, with packing multiplicity $\Lambda$, and let $\mathcal A_j(F)$ be its complete stopped admissibility condition, including the physical restrictions and the common assignments proved above. Suppose its coefficient windows and reciprocal pivots are bounded by powers of $P$.

\noindent\textup{(i)} For fixed $b_j\ge0$,
\begin{equation}\label{an:unified-window-incidence}
\int_{V_k}\left(\sum_jb_jv_j\mathbf1_{\mathcal A_j(F)}\right)^p\,dF\lesssim_{p,k} P^{A_{p,k}}\Lambda^{p-1}(1+\log P+\log d_1^{-1})^{B_{p,k}}d_2^{-\gamma}\sum_jb_j^pv_j^{p-\chi_k}.
\end{equation}

\noindent\textup{(ii)} For the physical family and low-linear selection, this gives~\eqref{an:physical-incidence}.

\noindent\textup{(iii)} The same bounds hold for any measurable selected subcollection of the admissibility condition.
\end{proposition}
\begin{proof}
\par\medskip\noindent\textit{Step 1. Assignment and label counts.}\par\nobreak\smallskip
Work in a fixed RQ bin. Lemma~\ref{geo:normal-pivot} excludes rank below four, and Lemma~\ref{reg:regular-preparation} assigns the quantitative full-rank region to bounded pieces and the regular region to prepared pieces. The remaining labels have the exhaustive surface assignment of Theorem~\ref{surf:common-tail}. The inner and outer windows and priority choices are fixed before an original label is fixed for coefficient integration.

On unbanded pieces set $s_I=1$ on bounded pieces, $s_I=H_I$ on regular curve pieces, and $s_I=S_I$ on surface pieces. The common assignments give count $P^C\Lambda s_I$ and
\[
\sum_I s_I^{-\eta}\le P^C(1+\log P+\log d_1^{-1})^C.
\]
The surface assignment has the stated area multiplicity and logarithmically many aspects. Lemma~\ref{reg:regular-preparation} supplies regular multiplicity. The resulting dyadic sums converge, with subunit scales costing powers of $P$.

\par\medskip\noindent\textit{Step 2. Fixed-frame coefficient moments.}\par\nobreak\smallskip
At an original fixed anchor, the surface tail~\eqref{surf:common-tail-eq} in Theorem~\ref{surf:common-tail} satisfies Proposition~\ref{an:common-incidence} since $q<\chi_k+7/2$ and its aspect exponent is $1/2>0$. The choice $\eta=1/[16(\chi_k+3)]$ gives
\begin{equation}\label{an:q-margin}
q<\chi_k+2+\frac18<\beta_k,
\end{equation}
because $\beta_k-\chi_k-2\ge131/513$; hence $\sum_HH^{q-\beta_k}<\infty$. Bounded coefficient sets have complete window volume $P^C$.


Proposition~\ref{an:common-incidence}, followed by finite H\"older over the piece types, proves~\eqref{an:unified-window-incidence}.

\par\medskip\noindent\textit{Step 3. Physical reconstruction.}\par\nobreak\smallskip
For the physical selection, $P\lesssim K_{\mathrm{lin}}^CL_R^C$, $\Lambda\lesssim_\vartheta L_R^3$, and $d_1\gtrsim R^{-1}$. Here the physical mesh $\vartheta$ is fixed, and $\rho=c_0\sqrt\lambda\sim L_R^{-1}$, so the factor $\vartheta^{-15}$ in \eqref{an:local-packing} is part of the constant. The conditioning and physical logarithms are therefore powers of $K_{\mathrm{lin}}$ and $L_R$. Finite H\"older over the RQ bins costs another power of $L_R$, proving~\eqref{an:physical-incidence}. Positivity gives the assertion for selected subcollections.
\end{proof}

\subsection{The upper bound}\label{completion:upper-bound}

Apply Proposition~\ref{an:physical-assembly} to Proposition~\ref{an:unified-incidence}. Here $\gamma=0$, so positivity follows directly from $p>\chi_k+3$. The coefficient incidence estimate gives the smooth-amplitude upper bound on the interval $\chi_k+3<p<\chi_k+3+1/16$. Once it is known at any $p_0$, it extends to every $p>p_0$ by the pointwise bound

\begin{equation}\label{an:larger-p}
|E_a(F)|^p\le\|a\|_1^{p-p_0}|E_a(F)|^{p_0}.
\end{equation}
For each $p>\chi_k+3$, choose $p_0<p$ in the established interval.


Apply Lemma~\ref{an:cube-reduction} with $d=3$. Every exponent already proved satisfies $p>\chi_k+3>\chi_k$, so its hypothesis on $p$ holds. This proves the upper bound for $\mathcal I_k$.

\subsection{Necessity and endpoint divergence}\label{completion:sharpness}\label{an:sharpness-subsection}

\begin{proposition}[Sharpness in every degree]\label{an:necessity}
For every $k\ge2$ and $0<p\le\chi_k+3$, one has $\mathcal I_k\notin L^p(V_k)$.
\end{proposition}
\begin{proof}
Use coordinates $(u,y,z)$ for a seed polynomial $J\in V_k$. Set its three coefficients of $u^{k-1}$, $u^{k-1}y$, and $u^{k-1}z$ equal to zero. Allow $c=[u^k]J$ to range in $(1,1+\delta)$ and each other free coefficient to range in $(-\delta,\delta)$. Thus the seed set has dimension $N_k-3$ and positive relative measure. Fix $\epsilon>0$ and introduce parameters
\[
\frac25<t<\frac35,\qquad |a|,|b|<\frac1{20},\qquad F(x,y,z)=J\left(\frac{x-t-ay-bz}{\epsilon},y,z\right)\pmod{\mathbb R}.
\]
The three extra parameters make this a map between spaces of the same dimension $N_k$. Factor its differential through the full affine coefficient substitution $S(x,y,z)=((x-t-ay-bz)/\epsilon,y,z)$. Lemma~\ref{an:character} gives its absolute determinant $\epsilon^{-\chi_k}$. Before this substitution, the free seed columns are the identity columns in all rows except the three prescribed zero rows; differentiation in $t,a,b$ gives the additional columns
\[
-\epsilon^{-1}J_u,\qquad -\epsilon^{-1}yJ_u,\qquad -\epsilon^{-1}zJ_u.
\]
In the three omitted rows these columns form the diagonal matrix with diagonal entries $-kc/\epsilon$. Indeed, the degree bound excludes $u^ky$ and $u^kz$, and the prescribed zero coefficients remove the possible off-diagonal contributions. The full absolute Jacobian is consequently
\begin{equation}\label{an:necessity-jacobian}
(kc)^3\epsilon^{-\chi_k-3}.
\end{equation}

This coefficient parameterization is injective for fixed $\epsilon$. The coefficient ratios satisfy $[x^{k-1}]F/[x^k]F=-kt$, $[x^{k-1}y]F/[x^k]F=-ka$, and $[x^{k-1}z]F/[x^k]F=-kb$. They recover $t,a,b$, and inverse affine substitution then recovers all coefficients of $J$. The change-of-variables theorem applied to this injective smooth map with nonzero Jacobian shows that its image has measure $\gtrsim_\delta\epsilon^{-\chi_k-3}$.

At the central seed, the normal integral is
\begin{equation}\label{an:normal-gamma}
\int_{\mathbb R}e^{iu^k}\,du=\begin{cases}\dfrac2k\Gamma(1/k)e^{i\pi/(2k)},&k\text{ even},\\[1mm]\dfrac2k\Gamma(1/k)\cos(\pi/(2k)),&k\text{ odd}.\end{cases}
\end{equation}
This follows from \cite[equations (5.9.8)--(5.9.9)]{DLMF} and $\Gamma(1+1/k)=\Gamma(1/k)/k$.

For fixed $(y,z)\in[0,1]^2$, put $s=t+ay+bz\in[3/10,7/10]$ and $L(y,z)=\int_{-s}^{1-s}e^{iJ(v/\epsilon,y,z)}\,dv$. For $k\ge3$, apply \cite[Lemma 6.4]{BGZZK} to $q(v)=(J(v/\epsilon,y,z)-J(0,y,z))/(2\pi)$ with $A=(2\pi)^{-1/k}\epsilon^{-1}$. Its leading coefficient is $cA^k$, and its other coefficients satisfy $|[v^j]q|\lesssim_k \delta A^j$ for $1\le j<k$. The cited asymptotic has error $O_k(\epsilon^{k-1})=o(\epsilon)$ and a leading constant close to a fixed nonzero number. Since $J(0,y,z)=O_k(\delta)$, multiplication by $e^{iJ(0,y,z)}$ preserves a common sector of angle less than $\pi$, as in the proof of \cite[Lemma 6.2]{BGZZK}.

For $k=2$, the three prescribed zero coefficients give $J(u,y,z)=cu^2+J(0,y,z)$. Formula~\eqref{an:normal-gamma} and one integration by parts in each tail give
\[
L(y,z)=\epsilon e^{iJ(0,y,z)}c^{-1/2}\sqrt\pi e^{i\pi/4}+O(\epsilon^2).
\]
Thus in both cases, for sufficiently small fixed $\delta$ and then small $\epsilon$, there is $|\omega_k|=1$ such that $\operatorname{Re}(\omega_kL(y,z))\gtrsim_k\epsilon$ throughout $[0,1]^2$. Integrating this inequality gives $|\mathcal I_k(F)|\gtrsim_k\epsilon$; the zero-constant representative changes only its constant phase factor.

Take $\epsilon=2^{-n}$. The image sets are disjoint if $\delta<2^k-1$, because their $x^k$ coefficients lie in the disjoint intervals $(2^{nk},(1+\delta)2^{nk})$. Using their measure lower bound from~\eqref{an:necessity-jacobian} and the integral lower bound gives
\[
\int_{V_k}|\mathcal I_k(F)|^p\,dF\gtrsim_p\sum_{n\ge n_0}2^{-n(p-\chi_k-3)}.
\]
This series diverges for $p\le\chi_k+3$, including equality.
\end{proof}

\begin{proof}[Proof of Theorem~\ref{root:main}]\label{an:main-proof}
Subsection~\ref{completion:upper-bound} proves the smooth-amplitude and cube upper bounds for every $k\ge14$ and $p>\chi_k+3$. Proposition~\ref{an:necessity} proves divergence for $0<p\le\chi_k+3$, including equality. Together these prove the theorem.
\end{proof}
\appendix
\section{Reference tables}\label{app:reference-tables}
\renewcommand{\thetable}{A.\arabic{table}}
\renewcommand{\theHtable}{A.\arabic{table}}
\setcounter{table}{0}


\begin{table}[H]
\centering\small
\caption{Gap coordinates for the six ordered cubic cases in Table~\ref{reg:all-regimes}. The additional-sample exponents $n$ are given in Table~\ref{reg:all-costs}. In IDs 2 and 4, $\theta=\theta_1$.}\label{app:table-ordered-cubic}
\begin{tabular}{cc}
\toprule
ID & $(u,v,w)$\\\midrule
2 & $(1,1,\theta)$\\
4 & $(\theta,1,1)$\\
8 & $(\theta_1,1,\theta_2)$\\
14 & $(\theta_1,\theta_2-1,1)$\\
16 & $(1,\theta_1,\theta_2)$\\
20 & $(\theta_1,\theta_2,\theta_3)$\\\bottomrule
\end{tabular}
\end{table}

\begin{table}[H]
\centering\small
\caption{Derivative-phase and parent-block dimensions and character exponents for the applications of Proposition~\ref{low:78:smult:atlas}. The derivative phase includes its constant and supplies the first character exponent; the remaining exponents come from the complete binary parent jets.}\label{app:table-parent-blocks}
\[
\begin{array}{c|c|c|c|c|c}
k&m&D=k-m+2&\binom{D+3}{3}&\text{parent dimensions}&\text{character exponents}\\\hline
7&3&6&84&35&126+84=210\\
8&3&7&120&44&210+120=330\\
8&4&6&84&36+44&126+84+120=330
\end{array}
\]
\end{table}

\section{The binary cubic case}\label{pl:section}
The binary cubic example displays the complete-jet count and its sharp threshold. We use factorial coefficients and $e(t)=\exp(2\pi i t)$ here and in the planar supplement; fixed changes of these conventions preserve integrability. Put
\begin{equation}\label{pl:moment-surface}
\mathbb X(r,s)=\left(r,s,\frac{r^2}{2},rs,\frac{s^2}{2},\frac{r^3}{6},\frac{r^2s}{2},\frac{rs^2}{2},\frac{s^3}{6}\right),
\end{equation}
\begin{equation}\label{pl:extension}
Ea(\xi)=\int_{\R^2}a(u)e(P_\xi(u))\,du,\qquad P_\xi(u)=\xi\cdot\mathbb X(u),\qquad \xi\in\R^9.
\end{equation}
Write $E1=E\mathbf1_{[0,1]^2}$. The following case of \cite[Theorem 1.2]{BGZZK} will be proved using the same coefficient-counting principle as the ternary argument.
\begin{corollary}[Binary cubic integrability]\label{pl:main}
For every $0<p<\infty$,
\begin{equation}\label{pl:main-equivalence}
E1\in L^p(\R^9)\quad\Longleftrightarrow\quad p>12.
\end{equation}
For $p>12$ and $a\in C_c^\infty((-1,2)^2)$, $\|Ea\|_{L^p}$ is bounded by finitely many smooth seminorms of $a$, uniformly on a fixed compact support class.
\end{corollary}

Use the frame conventions and affine frequency decomposition in Supplement~S1.1, especially Theorem~S1.2, with its notation $A_R$, $K=R^\beta$, $V=\delta^2/\tau$, and $\mathcal J_{\mathfrak t}$.

\subsection{The sharp threshold}\label{pl:upper-bound}
\begin{proposition}[Annular upper bound]\label{pl:annular-proposition}
For every $p>12$, there is $C_p<\infty$ independent of $\beta$ such that
\begin{equation}\label{pl:smooth-annular}
\int_{A_R}|Ea(\xi)|^p\,d\xi\lesssim_{p,\beta} \|a\|_{C^{N_{p,\beta}}}^pR^{(12-p)/3+C_p\beta},
\end{equation}
with finite $N_{p,\beta}$.
\end{proposition}
\begin{proof}
At one type, finite-sum H\"older, \textup{(S1.11)}, and \textup{(S1.13)} give
\begin{equation}\label{pl:scale-count}
\begin{split}
\mathfrak I_{\delta,\tau}:=\int_{A_R}\left|\sum_{J\in\mathcal J_{\mathfrak t}}E_Ja(\xi)\mathbf1_{\Omega_J}(\xi)\right|^p\,d\xi
&\lesssim_{p,\beta,a}K^{C_p}\tau^{-1}V^{-1}V^pV^{-10}\\
&=K^{C_p}\tau^{-1}V^{p-11}=K^{C_p}\delta^{2p-22}\tau^{10-p}.
\end{split}
\end{equation}
Since $\tau\ge\delta$ and $p>12$, \textup{(S1.10)} yields
\begin{equation}\label{pl:annular-bound}
\mathfrak I_{\delta,\tau}\lesssim_{p,\beta,a}K^{C_p}\delta^{p-12}\lesssim_{p,\beta,a}K^{C_p}R^{(12-p)/3}.
\end{equation}
The $O(L_R^2)$ types cost a fixed logarithmic power by H\"older, absorbed in $K^{C_p}$. The remainder contributes $O(R^{9-Mp})$, and $M$ can be chosen large. This proves \eqref{pl:smooth-annular}; bounded $R$ follows from $|Ea|\le\|a\|_1$.
\end{proof}

Choose $C_p\beta<(p-12)/6$ in \eqref{pl:smooth-annular} and sum the dyadic annuli to obtain the smooth-amplitude assertion of Corollary~\ref{pl:main}. A fixed affine rescaling takes any fixed compact amplitude support into $U_{\rm pl}$; the coefficient character $\chi_{2,3}=10$ from Lemma~\ref{an:character} transfers the smooth bound, with a constant depending on that support. Lemma~\ref{an:cube-reduction}, with $d=2$, $k=3$, and $\chi_{2,3}=10<p$, then gives the upper bound for the square.

\begin{proof}[Completion of the proof of Corollary~\ref{pl:main}]
The endpoint construction in the lower-bound paragraph of Supplement~S1.2 gives, for each sufficiently large dyadic $\Lambda$, $\gtrsim\Lambda^2$ disjoint coefficient sets $\Omega_{m,n,\Lambda}$ of volume $\sim\Lambda^{10}$ on which $|E1|\gtrsim\Lambda^{-1}$ (Lemmas~S1.4 and~S1.5). Their disjointness also holds between dyadic levels, and hence
\[
\int_{\bigcup_{m,n}\Omega_{m,n,\Lambda}}|E1(\xi)|^p\,d\xi\gtrsim\Lambda^2\Lambda^{10}\Lambda^{-p}=\Lambda^{12-p}.
\]
The dyadic sum diverges for every $p\le12$, including the endpoint. Together with the upper bound, this proves \eqref{pl:main-equivalence}.
\end{proof}


\section{Estimates used in the companion paper}\label{app:shared-estimates}
Use the fixed-degree and constant-dependence conventions of Section~\ref{an:section}, and the notation of Sections~\ref{geo:section}, \ref{reg:section}, and \ref{atlas:section}.

\subsection{Summation over dyadic rate ranges}\label{an:band-summation}

\begin{proposition}[Weighted incidence with a dyadic rate parameter]\label{low:6:r3:weighted-holder}
Fix $p>1$, $0<a<1$, $\eta>0$, $q=(1+\eta)(p-1)$, $b>q$, $\tau>0$, an integer $m_0\ge0$, and an exponent $C_*>0$. Consider a common assignment as in Definition~\ref{def:assignment}. Each piece $I$ has dyadic parameters $0<\nu_I\le P^{C_*}$ and $1\le K_I\le S_I$. Suppose there are constants $C_0,M\ge1$ such that, for every phase $F$,
\[
\begin{aligned}
\#\{j:j\text{ is assigned to }I\}&\le C_0\Lambda S_I,\\
\#\{I\in\mathscr P(F):(\nu_I,S_I,K_I)=(\nu,S,K)\}&\le M(1+|\log\nu|+\log S)^{m_0}.
\end{aligned}
\]
Let $\mathcal A\ge0$ be a constant. For every fixed $j$, assume that the assigned sets in normalized coefficient coordinates satisfy
\[
\operatorname{vol}E_j(\nu,S,K)\le P^{C_*}\mathcal A\,\nu S^{-b}K^{-\tau}.
\]
The constants $C_0,M,\mathcal A$ are independent of $F,j,\nu,S,K$. They may depend on fixed physical parameters, and their dependence is retained in the following bound: for every fixed nonnegative sequence $b_j$,
\[
\int_{V_k}\left(\sum_j b_jv_j\mathbf1_{\mathcal A_j(F)}\right)^p\,dF
\lesssim P^C(C_0M\Lambda)^{p-1}\mathcal A\sum_j b_j^pv_j^{p-\chi_k}.
\]
Here $C$ and the implicit constant depend only on the fixed exponents and $m_0$. Both sums over $j$ run over the original fixed index family. On the left, $\mathbf1_{\mathcal A_j(F)}$ is one exactly when $j$ is assigned to a piece with positive rate. With $a_I(F)=\sum_{j\text{ assigned to }I}b_jv_j$, the common assignment gives $\sum_j b_jv_j\mathbf1_{\mathcal A_j(F)}=\sum_{I\in\mathscr P(F)}a_I(F)$. There is no lower cutoff on the positive rates $\nu_I$.
\end{proposition}
\begin{proof}
For each piece $I$, set
\begin{equation}\label{low:6:r3:holder-weight}
w_I=\nu_I^{-a}S_I^{\eta(p-1)}.
\end{equation}
The assumed number of pieces in each dyadic range gives
\begin{equation}\label{low:6:r3:inverse-sum}
\sum_{I\in\mathscr P(F)}w_I^{-1/(p-1)}
\le M\sum_{0<\nu\le P^{C_*}}\sum_{S\ge1}\sum_{1\le K\le S}\nu^{a/(p-1)}S^{-\eta}(1+|\log\nu|+\log S)^{m_0}
\lesssim MP^C.
\end{equation}
All parameter sums are dyadic. There are $O(1+\log S)$ possible values of $K$. The positive exponents $a/(p-1)$ and $\eta$ make the sums converge even with the displayed logarithmic factors.

Apply the two H\"older inequalities in the proof of Proposition~\ref{an:common-incidence}, using the weights~\eqref{low:6:r3:holder-weight} and the labels assigned to each piece. The inverse-weight factor is at most $(MP^C)^{p-1}$. The count within a piece contributes $(C_0\Lambda S_I)^{p-1}$, so its product with $w_I$ is $(C_0\Lambda)^{p-1}\nu_I^{-a}S_I^q$. For a fixed $j$, put $J=\mathcal L_{T_j}F$ and let $\nu_j(J),S_j(J)$ be the parameters of its assigned piece. Set both parameters to one when $j$ is inadmissible, and write $\mathbf1_j(J)=\mathbf1_{\mathcal A_j(\mathcal L_{T_j}^{-1}J)}$. Then
\begin{equation}\label{low:6:r3:anchored-sum}
\int_{V_k}\nu_j(J)^{-a}S_j(J)^q\mathbf1_j(J)\,dJ
\le P^{C_*}\mathcal A\sum_{0<\nu\le P^{C_*}}\nu^{1-a}\sum_{S\ge1}S^{q-b}\sum_{1\le K\le S}K^{-\tau}
\lesssim P^C\mathcal A.
\end{equation}
These sums converge because $1-a>0$, $b-q>0$, and $\tau>0$. For the contribution of $j$, the change back to $F$ gives $dF=v_j^{-\chi_k}dJ$. Multiplying by $(b_jv_j)^p$ and summing over $j$ proves the displayed incidence bound. Nonnegativity and monotone convergence justify countably many pieces.
\end{proof}

\subsection{Geometric preparation}\label{app:shared-geometry}

\begin{theorem}[Rank four in degrees four and five]\label{geo:rank-four-low}
Let $k\in\{4,5\}$, $J\in V_k$, $\ell(w)=p+Bw$ and $\widehat J=(J,\ell)$. Assume $\|\widehat J\|\le P$ and $\|[u\cdot\nabla J]\|^2+|Bu|^2\ge\kappa^2|u|^2$ for all $u\in\mathbb R^3$, where $P\ge2$ and $\kappa\ge P^{-1}$. Suppose $\operatorname{rank}R_{\widehat J}=4$.

\noindent\textup{(i)} Then $B=0$ and
\begin{equation}\label{low:4:rc:rankfour-form}
\begin{gathered}
J=a q^2+bq+\lambda y+\mu z\pmod{\mathbb R},\qquad q=x-Q(y,z),\\
Q(0)=0,\quad\deg Q\le2,\quad a\ne0,\quad\det D^2Q\ne0.
\end{gathered}
\end{equation}

\noindent\textup{(ii)} Conversely every phase of this form has an injective phase translation differential and, with $B=0$, upper-affine orbit rank four.

\noindent\textup{(iii)} On its bounded stopped part, the parameters, $|a|^{-1}$, and $\|(D^2Q)^{-1}\|$ are bounded by powers of $P$.
\end{theorem}
\begin{proof}
Apply Theorem~\ref{geo:classification} with $k\in\{4,5\}$. The profile is nonlinear. A cubic invariant would force $\deg f\le\lfloor k/3\rfloor=1$ and is therefore excluded. The quadratic-invariant branch has $\deg f\le\lfloor k/2\rfloor=2$. Writing $f(q)=aq^2+bq$ modulo constants gives the displayed form with $q=x-Q(y,z)$ and $Q(0)=0$. The same theorem supplies the converse, the vanishing physical linear part, and the conditioned parameter bounds.
\end{proof}

\begin{corollary}[Quantitative entry into the surface region]\label{geo:rank-four-entry-low}
There is an absolute exponent $C_0$ with the following property. Fix any integer $L\ge1$. For a sufficiently large $C_L$, every degree-$k$ augmented jet, $k\in\{4,5\}$, satisfying these stopped-window bounds and
\[
\sigma_5(R_{\widehat J})\le P^{-C_L}
\]
has $\|B\|_{\mathrm{op}}\le\kappa/2$ and lies, in phase coefficient norm, within $P^{-L}$ of the square-quadric family
\begin{equation}\label{low:4:rc:surface-family}
\mathfrak K_{P^{C_0}}=\left\{\begin{aligned}
&a\bigl(x-\widetilde Q(y,z)\bigr)^2+\lambda y+\mu z\pmod{\mathbb R}:\\
&|a|+|a|^{-1}+\|\widetilde Q\|+\|(D^2\widetilde Q)^{-1}\|+|\lambda|+|\mu|\le P^{C_0}
\end{aligned}\right\},
\end{equation}
where $\deg\widetilde Q\le2$. It also satisfies the phase-only stop
\[
\|[u\cdot\nabla J]\|\ge(\sqrt3/2)\kappa |u|.
\]

\end{corollary}
\begin{proof}
Use Lemma~\ref{geo:rank-entry} with $k\in\{4,5\}$, with the prescribed complete defect and the classified family just identified. Complete its square by $\widetilde Q=Q-b/(2a)$ to obtain the stated square-quadric window, with a fixed polynomial enlargement. Augmented proximity also makes the physical linear part as small as required; subtracting its squared contribution from the stated translation stop gives the asserted phase-only lower bound.
\end{proof}

In particular $L$ may be chosen to meet the quantitative root and Hessian neighborhoods in the actual critical-surface construction.
\begin{lemma}[Relative algebraic preparation before amplitude selection]\label{geo:relative-preparation}
Fix an algebraic format, an integer $N$, and $\epsilon>0$.

\noindent\textup{(i)} Let $f$ be a real algebraic branch on an interval of length $W$. Apart from an identically zero branch, one can retain a union of boundedly many intervals, losing at most $\epsilon W$ in total, such that every retained interval has length at least $cW$ and
\[
|f^{(j)}(t)|\lesssim_{N,\epsilon}|f(t)|W^{-j}\quad(0\le j\le N),\qquad |f(t)|\sim_{N,\epsilon}|f(t')|
\]
throughout each retained interval. The constants and number of intervals depend only on the fixed format, $N,\epsilon$.

\noindent\textup{(ii)} The construction is uniform in additional real parameters and can be made semialgebraically. A fixed finite collection of branches can be prepared simultaneously. Identically zero and nonzero constant branches are retained separately.
\end{lemma}
\begin{proof}
After a bounded algebraic stratification, a nonconstant branch satisfies an irreducible equation $P(t,f)=0$ of bounded degree. The leading coefficient in the second variable, its discriminant, and $P(t,0)$ are nonzero polynomials. Their complex zeros are the exceptional values; their number is bounded solely by the degree. Remove intervals of radius $\eta W$ around their real parts and around the two original endpoints, where $\eta>0$ is a sufficiently small fixed multiple of $\epsilon$. Discard the remaining components shorter than a further fixed multiple of $\eta W$, and trim their endpoints. The removed length is at most $\epsilon W$ after adjusting these fixed constants. The retained intervals have length at least $cW$, and every point of them has a complex disk of radius $cW$ free of exceptional values.

On each such disk the selected algebraic branch extends holomorphically and without zeros. The function $f$ has bounded valence there: for each $w_0$, the equation $P(t,w_0)=0$ has at most $\deg_tP$ roots. It cannot vanish identically in $t$, since irreducibility would then make the branch the excluded constant $w_0$.

The distortion estimate for nonvanishing functions of bounded valence \cite[Theorem 3.1]{FriedlandYomdin}, applied after rescaling each disk, gives $|f(z)|\sim|f(t)|$ on its concentric half disk. Cauchy's estimates yield $|f^{(j)}(t)|\lesssim_{N,\epsilon}|f(t)|W^{-j}$. In particular $|f'/f|\lesssim_{N,\epsilon}W^{-1}$; integration along a retained interval, whose length is at most $W$, gives the claimed value comparability. The exceptional roots, their real and imaginary parts, and the ordered interval selections are fixed-format semialgebraic relations. Taking the union of the exceptional sets gives the simultaneous version.
\end{proof}

\begin{lemma}[Derivative preparation on a characteristic]\label{geo:characteristic-preparation}
Let $Z=\partial_t+d(y,t)\partial_y$, where $d$ is of fixed algebraic format on its regular branch. Let $\gamma(t)=(y(t),t)$ be a connected trajectory segment of $Z$ in a fixed-format regular semialgebraic branch domain. Let $f$ be an algebraic scalar of fixed format on that domain and suppose $|f\circ\gamma|\le M$ on a trajectory interval $I$ of length $H$.

\noindent\textup{(i)} For each fixed integer $N$ and $\epsilon>0$, outside a subset of $I$ of measure at most $\epsilon H$ one has
\[
\left|\frac{d^j}{dt^j}(f\circ\gamma)\right|\lesssim_{N,\epsilon} M H^{-j},\qquad 1\le j\le N.
\]
The constants depend on the fixed format and $N,\epsilon$, not on the coefficients or trajectory length. Finitely many such assertions can be imposed at the same good point.

\noindent\textup{(ii)} Moreover $\operatorname{Var}_I(f\circ\gamma)\lesssim M$.

\noindent\textup{(iii)} For an algebraic scalar $v$, the intersection of $\gamma$ with any band $\nu/2<|v|<\nu$ has a uniformly bounded number of connected intervals.
\end{lemma}
\begin{proof}
Each derivative is $(Z^jf)\circ\gamma$, with $Z^jf$ algebraic of fixed format: implicit differentiation uses boundedly many rational operations and divisions by the nonzero pivots. Its level curves have boundedly many graph arcs $y=\psi(t)$ and vertical arcs. On a graph arc the crossing derivative $d(\psi(t),t)-\psi'(t)$ has boundedly many sign intervals and isolated zeros. A trajectory crosses a nonzero sign interval at most once, since a second crossing of the same orientation forces an opposite crossing between them. An identically zero interval is an integral curve, so uniqueness excludes isolated crossings. Vertical arcs are met at most once; tangencies occur at the boundedly many isolated zeros, and coincident subarcs add no components. Thus scalar sublevels, superlevels and bands pull back to uniformly boundedly many intervals, also proving the assertion for $v$. The one-dimensional area formula \cite[Theorems 1--2]{low:area} bounds variation by the isolated-crossing count for almost every level times the scalar range, giving $\operatorname{Var}_I(f\circ\gamma)\lesssim M$.

Apply the one-variable finite-difference argument in the proof of Lemma~\ref{geo:preparation} to $f\circ\gamma$, using the uniform component count just proved. Taking thresholds $C_{N,\epsilon}MH^{-j}$ and summing over the finite derivative list gives the asserted exceptional-set bound.
\end{proof}

\begin{lemma}[Derivative good points and curved strips]\label{geo:good-points-strips}
Let $L\ge W\ge P^{-C}$ and let a fixed-format algebraic scalar $f$ be bounded by $M$ on a rectangle $R$ of dimensions $L,W$.

\noindent\textup{(i)} For any fixed derivative order and any prescribed exceptional fraction $\varepsilon>0$, a subset with relative area at least $1-\varepsilon$ satisfies
\[
 |\partial_u^i\partial_v^j f|\lesssim_{\varepsilon} M L^{-i}W^{-j}
\]
for the finite list of derivatives. The assertion also holds simultaneously for finitely many scalars.

\noindent\textup{(ii)} If $|f|\le M$ only on a fixed-format set $A\subset R$ with $0<q\le1$ and $|A|\ge q|R|$, one can choose $G\subset A$ such that
\[
|G|\ge c q^C|R|,\qquad |\partial_u^i\partial_v^j f|\le Cq^{-C}ML^{-i}W^{-j}\quad\text{on }G
\]
for the prescribed finite derivative list. All exponents depend only on the fixed derivative orders and format.

\noindent\textup{(iii)} Under the hypotheses of Lemma~\ref{geo:planar-strips}, the labels admit common contained strips with the area and diameter bounds stated there and centerline bounds $|\zeta^{(j)}|\le P^CL^{1-j}$ for $1\le j\le N$, where $N$ is the prescribed order. For a chosen parameter point $s_0$, $|\zeta(s)-\zeta(s_0)|\le P^CL$.
\end{lemma}
\begin{proof}
For a full rectangle, apply Lemma~\ref{geo:preparation} after rescaling its two sides. For a fixed-format subset of relative area $q$, Fubini retains a set of sections of measure at least a fixed power of $q$ on which the occupied length is at least a fixed power of $q$ times the side length. Each section has boundedly many interval components; discard the short components and apply the one-dimensional finite-difference argument on the remaining ones. Repeat the section argument in both coordinate directions, applying it on the set retained at the preceding stage to the derivative functions needed next. Take the intersection of the resulting nested good sets. These derivative functions and selected sets have bounded format. The number of stages is fixed, and each stage costs only a fixed power of $q$ in the retained area and derivative constants, giving the stated $cq^C$ area bound. The strip assertion is Lemma~\ref{geo:planar-strips}.
\end{proof}

\subsection{Regular sections and the physical cutoff}\label{app:shared-regular}

\begin{lemma}[Sections of the regular $d$ chart]\label{reg:d-sections}
Fix $k\ge2$ and a regular $d$ chart of Lemma~\ref{reg:regular-preparation}. Let $L_d$ be the augmented orbit matrix with the $D$ column omitted, assume $\sigma_{\min}(L_d)\ge P^{-C}$, and let $\mathcal U_d=\{g:d(g)=0\}$.

\noindent\textup{(i)} The admissibility condition can be assigned to a bounded number of semialgebraic sections $s_\nu(t)$ with $d(s_\nu(t))=t$, together with a bounded number of exceptional fibers. Every assigned state has the form $g=s_\nu(t)u$ with $u\in\mathcal U_d$ and $\|u\|+\|u^{-1}\|\le P^C$ in homogeneous affine matrices. The exceptional parameter values belong to the finite list supplied by the semialgebraic decomposition. The sections have uniformly bounded description complexity.

\noindent\textup{(ii)} The labels assigned over a parameter interval of length at most one have count at most $P^C\Lambda$.
\end{lemma}
\begin{proof}
Let $g(r)$ be a semialgebraic path in a connected component of a fixed $d$ fiber. Its body velocity $\xi=g^{-1}g'$ lies in the Lie algebra of $\mathcal U_d$, so $\widehat J'=L_d\xi$ and
\[
|\xi|\le P^C\|\widehat J'\|.
\]
Lemma~\ref{geo:paths} supplies a fixed-format path with augmented-jet variation at most $P^C$. Thus $\int|\xi|\le P^C$. In homogeneous affine matrices the strictly upper triangular algebra has nilpotency step three, so the relative matrix is a polynomial of degree at most three in iterated integrals of $\xi$. The relative map and its inverse are bounded by $P^C$. Lemma~\ref{geo:conditioning} bounds the number of connected fiber components.

Semialgebraic trivialization over a finite partition of the $t$ axis \cite[Theorem 5.46]{BPR} supplies one section in each fiber component. Assign ties in a fixed order. Put $v=-L_d^+(D\widehat J)$, with the left pseudoinverse of Definition~\ref{def:regular-coordinates}. Along a section the $D$ coefficient of its body velocity is one, while its other five coefficients satisfy
\begin{equation}\label{reg:d-section-rates}
\boldsymbol\beta=v+L_d^+\widehat J'.
\end{equation}
\equationalias{low:4:rd:least_squares}\equationalias{low:5:rd5:least_squares}
Thus $|s^{-1}s'|\le P^C(1+\|\widehat J'\|)$. On a unit interval its integral is bounded by $P^C$, using the same bounded-variation argument. The fixed-fiber comparison and family packing prove the unit-interval claim. Exceptional fibers are controlled by the first paragraph.
\end{proof}

\begin{lemma}[Bound for the middle singular value]\label{reg:middle-cutoff}
For any fixed $k\ge2$, let an original prepared section satisfy the bounded augmented window, the stop~\eqref{reg:stop}, and $\|\widehat J^{\prime}\|\le P^C/H$. Suppose its original body field has norm between $P^{-C}$ and $P^C$, and its physical singular values remain comparable, up to powers of $P$, to the original $d_i$. Then $H\le P^C/d_2$.
\end{lemma}
\begin{proof}
For an affine field $V(x)=Nx+b$ put $M_V=[N\ b]$ and normalize $\norm{M_V}_F=1$. On the compact bounded stopped jet set, the function $\max\{\norm{[VJ]}+\norm{V\ell},\sigma_2(M_V)\}$ is positive. Indeed, a zero would give an exact rank-one annihilator $V=L(x)u\cdot\nabla$ with $u\ne0$ and $L$ nonzero affine. The physical equation gives $Bu=0$. The phase equation says $L(u\cdot\nabla J)$ is constant. If $L$ is nonconstant, degree and the integral-domain property imply $u\cdot\nabla J=0$; if $L$ is constant, that derivative is constant. Either contradicts \eqref{reg:stop}. Lemma~\ref{geo:conditioning} therefore bounds the maximum below by $P^{-C}$. If $\norm{[VJ]}+\norm{V\ell}<\tfrac12P^{-C}$, then $\sigma_2(M_V)\ge P^{-C}$.

The original prepared body field has norm bounded above and below; its augmented action is $O(H^{-1})$. Moreover $B'=BN$, $p'=Bb$, and augmented-jet preparation gives $\norm{B'}+\abs{p'}\lesssim_P H^{-1}$. Thus $\norm{BM_V}\lesssim_P H^{-1}$. Since $M_V$ has two singular values bounded below, $B$ contracts a two-dimensional subspace by a factor $\lesssim_P H^{-1}$. The min--max principle gives $d_2\lesssim_P H^{-1}$. All constants here are bounded by powers of $P$, including the normalization of $M_V$. If $H$ is too short for the small-action step, its polynomial bound and the upper bound for $d_2$ give the same conclusion after increasing the constant. This argument uses the original prepared section, before any auxiliary anisotropic change of frame.
\end{proof}

\subsection{Normalized frames and projected packing}\label{atlas:normalized-frames}

\begin{lemma}[Finite normal fibers and transverse Jacobians]\label{atlas:normal-fibers}
Let $G$ be a polynomial of degree at most $D_G$, and write $h(x,y,z)=(x+ay+bz+c,y+dz+e,z+f)$ and $\widetilde G=G\circ h$. At a fixed tangential center $(e,f)$ impose $\widetilde G(0)=\widetilde G_y(0)=\widetilde G_z(0)=0$, with $\lambda=\widetilde G_x(0)\ne0$. Complete the normalization either by chart I, $\widetilde G_{yz}(0)=0$ with $\widetilde G_{yy}(0)\ne0$, or by chart II, $\widetilde G_{zz}(0)=0$ with $\widetilde G_{yz}(0)\ne0$.

\noindent\textup{(i)} There are at most $3D_G$ normalized frames at that center.

\noindent\textup{(ii)} The absolute Jacobian of the four defining equations in $(c,a,b,d)$ is respectively
\[
 |\lambda|^3|\widetilde G_{yy}(0)|,\qquad 2|\lambda|^3|\widetilde G_{yz}(0)|.
\]
\end{lemma}
\begin{proof}
First solve $G(c,e,f)=0$ with $G_x(c,e,f)\ne0$. There are at most $D_G$ such roots, including when the restriction to this line is identically zero, in which case none has the required derivative. At a regular root put $a_0=-G_y/\lambda$ and $b_0=-G_z/\lambda$. The first derivative equations force
\begin{equation}\label{atlas:normal-ab}
a=a_0,\qquad b=b_0+da_0.
\end{equation}
\equationalias{low:4:sf:eq:ab}\equationalias{low:5:qsf:eq:ab}
Let $H=(h_{ij})$ be the Hessian of $G$ on the tangent vectors $(a_0,1,0)$ and $(b_0,0,1)$. Chart I requires $h_{12}+dh_{11}=0$ and has at most one solution with its nonzero pivot. Chart II requires $h_{22}+2dh_{12}+d^2h_{11}=0$ and has at most two solutions with its nonzero pivot. Identically zero equations supply no permitted solutions. This gives the bound $3D_G$.

For the Jacobian, eliminate $c$ by the first equation, then $a,b$ by the two first derivative equations. Each diagonal derivative is $\lambda$. Along these three equations,~\eqref{atlas:normal-ab} holds; the last derivative in $d$ is $h_{11}$ in chart I and $2(h_{12}+dh_{11})$ in chart II. Multiplication gives the two asserted determinants. This is the ordinary normalization Jacobian, before any coefficient substitution.
\end{proof}

\begin{proposition}[Projected packing from regular normalization]\label{atlas:projected-packing}
Fix a phase $F$ and a polynomial $G$ of degree at most the fixed number $D_G$. Let $(g_i)_{i=1}^{N(F)}$ be a finite indexed family in $\mathcal U$. For each $i$, choose a normalized frame $h_i=g_i\circ k_i$ as in Lemma~\ref{atlas:normal-fibers}. The homogeneous affine matrices of $k_i$ and $k_i^{-1}$ have norms at most $P^C$. Put $\widetilde G_i=G\circ h_i$ and assume
\[
 \|\widetilde G_i\|_{\mathrm{coeff}}\le P^C,\qquad |\partial_x\widetilde G_i(0)|\ge P^{-C},
\]
where the coefficient norm includes every coefficient of $\widetilde G_i$. In chart I assume $|\partial_{yy}\widetilde G_i(0)|\ge P^{-C}$; in chart II assume $|\partial_{yz}\widetilde G_i(0)|\ge P^{-C}$. Let $\mathcal N_F^+$ be the fixed outer set of normalized frames, retaining the coefficient and physical restrictions on $F$ in Definition~\ref{def:admissibility}. Assume that the normalized solutions continue as $h_i u_i(s,t)\in\mathcal N_F^+$ for
\[
 |s|,|t|<r_P,\qquad r_P\ge P^{-C},
\]
where $u_i(s,t)\in\mathcal U$ is the local normalization correction in coordinates relative to $h_i$, continued from $u_i(0,0)=\mathrm{id}$, and has tangential translation coordinates $(s,t)$ and its affine matrix and inverse have norms at most $P^C$. Suppose the original indexed family satisfies the relative-neighborhood bound \eqref{an:local-packing}, with parameter $\Lambda$, in the triangular coordinates of Lemma~\ref{an:rq}. Then
\[
N(F)\le P^C\Lambda |E_F^+|,
\]
where
\[
 E_F^+=\{(e(h),f(h)):h\in\mathcal N_F^+\},
\]
and $e(h),f(h)$ are the last two translation coordinates in \eqref{geo:group}. The count $N(F)$ retains the multiplicity of the original indices.
\end{proposition}
\begin{proof}
The Jacobians in Lemma~\ref{atlas:normal-fibers}, together with the coefficient bounds, give an inverse-power radius for solving the normalization equations in coordinates relative to $h_i$. The continuation hypothesis places these solutions in $\mathcal N_F^+$ on the stated square. The coefficient calculation uses $G\circ h_i$, so its bounds are independent of the size of $h_i$ itself.

If $(e_i,f_i)$ and $d_i$ are its tangential center and shear, this square projects to
\[
 \mathcal P_i=(e_i,f_i)+\{(s+d_it,t):|s|,|t|<r_P\}\subset E_F^+.
\]
Its area is $4r_P^2$, since the tangential shear has determinant one. At a point common to several parallelograms, Lemma~\ref{atlas:normal-fibers} gives at most $3D_G$ possible continued frames. For each such frame, all original family frames differ by bounded right factors, obtained from $k_i,u_i$ and their inverses. The right factors and their inverses have bounds $P^C$, so \eqref{an:local-packing} bounds the overlap by $3D_GP^C\Lambda$. Therefore
\[
4r_P^2N(F)=\int_{E_F^+}\sum_{i=1}^{N(F)}\mathbf1_{\mathcal P_i}\,de\,df\le3D_GP^C\Lambda|E_F^+|.
\]
Dividing by $4r_P^2$ gives the claimed count.
\end{proof}

\begin{theorem}[Algebraic surface packing cutoff]\label{atlas:affine-area}
Let $\Sigma=\{(c(y,z),y,z):(y,z)\in E\}$ be a regular graph branch of a nonzero polynomial equation of degree at most $D$, with $|\det D^2c|\ge\kappa>0$. Let $B$ be invertible and suppose the physical image under $T(w)=b+Bw$ lies in a ball of radius $R$. Then
\[
|E|\lesssim_D \kappa^{-1/4}R^{3/2}|\det B|^{-1/2}.
\]
\end{theorem}
\begin{proof}
For a regular surface $M$, let $K$ denote Gaussian curvature and $dA$ Euclidean surface area measure. Define its absolute affine curvature area by
\[
\mathcal A(M)=\int_M |K|^{1/4}\,dA.
\]
The absolute curvature in this definition allows either Hessian signature; the required affine transformation law follows from the direct calculation below.

The regular area of a degree-$D$ algebraic surface in a radius-$R$ ball is $\lesssim_D R^2$ \cite[Proposition 2.5]{GuthArea}. For a generic normal, its preimages with nonzero Gaussian curvature are nonsingular zeros of $G=0$, $u\cdot\nabla G=0$, and $w\cdot\nabla G=0$, where $u,w$ span the normal's orthogonal plane. Their degrees are $D,D-1,D-1$, so \cite[Corollary 1.3]{low:roots} bounds their number by $D(D-1)^2$. Apply this on square-free components; zero-curvature points contribute zero Gauss Jacobian. The Gauss-map area formula \cite[Theorem 1]{low:area} gives $\int_M|K|\,dA\lesssim_D 1$. H\"older therefore gives
\[
\mathcal A(M)\le
\left(\int_M|K|\,dA\right)^{1/4}
\operatorname{Area}(M)^{3/4}\lesssim_D R^{3/2}.
\]
For a graph, $dA=\sqrt{1+|\nabla c|^2}\,dy\,dz$ and $K=\det D^2c/(1+|\nabla c|^2)^2$, so
\[
|K|^{1/4}dA=|\det D^2c|^{1/4}\,dy\,dz.
\]
If $v=|\det B|$ and $\alpha=|B^{-T}n|$ for the original unit normal, the cross-product identity $(Bu)\times(Bw)=(\det B)B^{-T}(u\times w)$ gives the area factor $v\alpha$. The transformed unit normal is $B^{-T}n/\alpha$ up to sign. For a local parametrization $\sigma$ of $M$, pairing it with $B\partial_{ij}\sigma$ shows that the second fundamental form is multiplied by $1/\alpha$ up to sign, while the determinant of the first fundamental form is multiplied by $(v\alpha)^2$. Therefore
\[
dA_B=v\alpha\,dA,\qquad |K_B|=|K|/(v^2\alpha^4).
\]
Thus $\mathcal A(BM)=v^{1/2}\mathcal A(M)$. The transformed defining polynomial has the same degree. Apply the physical bound to $T\Sigma$ and use $\mathcal A(\Sigma)\ge\kappa^{1/4}|E|$.
\end{proof}

\begin{lemma}[Physical-center Gauss incidence in fixed degree]\label{surf:physical-gauss-density}
Let $F\in V_k$ and let $J=[F\circ h]$ be an ordinary ($m=2$) Gauss-normalized endpoint, with $J_x=J_{xy}=J_{xz}=J_{xyz}=0$ at the origin, $\lambda=J_{xx}\ne0$, and $r=c_{yy}\ne0$.

\noindent\textup{(i)} Keeping the physical translation coordinates $(e,f)$ as point variables gives
\[
dF\,de\,df=C_k|\lambda|^4|r|\,dJ_{N_k-4}\,da\,db\,dc_0\,dd\,de\,df.
\]

\noindent\textup{(ii)} Put $\eta=e-df$. Assume the bounded complete windows and full-slice Jacobian and multiplicity hypotheses of Lemma~\ref{surf:null-density}, and the phase-only $H_0$ pivot bound at the fixed frame. For a slice map depending only on $J$, the conditional $H_0$ integral at fixed slice outputs and $(\eta,f)$ is at most $P^C$. In particular the fixed original frame with the null normalization uses its own $2|\lambda_0|^4|B_0|$ pivot, independently of the endpoint's $r$, where $\lambda_0$ and $B_0$ are the normal and mixed-curvature pivots at that fixed frame.
\end{lemma}
\begin{proof}
Apply chart I of Lemma~\ref{atlas:normal-fibers} to $G=F_x$. At fixed physical center its determinant is $|\lambda|^3|\widetilde G_{yy}(0)|=|\lambda|^4|r|$. In body columns these directions are $(A,C,T_x,D-aC)$; subtracting the $aC$ multiple preserves the determinant. The remaining coefficient coordinates are free slice coordinates and the complete upper-unipotent coefficient pullback has determinant one. This proves the joint incidence identity.

Now use precisely the unit-Jacobian substitution and triangular factorization $h=k g_0(\eta,f)$ in Lemma~\ref{surf:null-density}. The fixed frame condition lies on a single $H_0$ orbit after the intrinsic outputs and $(\eta,f)$ are fixed. Lemma~\ref{surf:haar} uses the original pivot from Lemma~\ref{geo:original-subgroup-pivots}; in a null normalization its value is the minor $2|\lambda_0|^4|B_0|$ computed in Lemma~\ref{surf:null-density}. Any additional endpoint rescaling contributes its separate Jacobian in the subsequent count.
\end{proof}

\begin{thebibliography}{99}
\bibitem{BGZZK} S. Basu, S. Guo, R. Zhang, and P. Zorin-Kranich, \emph{A stationary set method for estimating oscillatory integrals}, J. Eur. Math. Soc. \textbf{28} (2026), 2067--2100. \href{https://doi.org/10.4171/JEMS/1506}{doi:10.4171/JEMS/1506}.
\bibitem{BPR} S. Basu, R. Pollack, and M.-F. Roy, \emph{Algorithms in Real Algebraic Geometry}, second edition, Algorithms and Computation in Mathematics, vol. 10, Springer, Berlin, 2006. \href{https://doi.org/10.1007/3-540-33099-2}{doi:10.1007/3-540-33099-2}.
\bibitem{low:roots} M. Bafna, M. Sudan, S. Velusamy, and D. Xiang, \emph{Elementary analysis of isolated zeroes of a polynomial system}, arXiv:2102.00602, 2021. \url{https://arxiv.org/abs/2102.00602}.
\bibitem{low:area} R. L. Jerrard, \emph{Area and coarea}, MAT1501 lecture notes, University of Toronto, 7--17 October 2013. \url{https://www.math.toronto.edu/rjerrard/1501/notes/week5-6.pdf}.

\bibitem{Moc96} G. Mockenhaupt, \emph{Bounds in Lebesgue spaces of oscillatory integral operators}, Habilitationsschrift, Universit\"at Siegen, 1996.
\bibitem{LewisSendov} A. S. Lewis and H. S. Sendov, \emph{Twice differentiable spectral functions}, SIAM J. Matrix Anal. Appl. \textbf{23} (2001), 368--386. \href{https://doi.org/10.1137/S089547980036838X}{doi:10.1137/S089547980036838X}.
\bibitem{EtingofSL2} P. Etingof, \emph{Representations of Lie groups and Lie algebras}, Lecture 11, MIT 18.745, Lie Groups and Lie Algebras I, Fall 2020, Theorems 11.16 and 11.18. \href{https://ocw.mit.edu/courses/18-745-lie-groups-and-lie-algebras-i-fall-2020/mit18_745_f20_lec11.pdf}{MIT OpenCourseWare lecture notes}.
\bibitem{deBoor} C. de Boor, \emph{Divided differences}, Surv. Approx. Theory \textbf{1} (2005), 46--69. \url{https://surveys-in-approximation-theory.com/papers/2/2.pdf}.
\bibitem{FriedlandYomdin} O. Friedland and Y. Yomdin, \emph{$(s,p)$-Valent functions}, arXiv:1503.00325v1, 2015. \url{https://arxiv.org/abs/1503.00325v1}. Theorem 3.1 records the nonvanishing multivalent distortion theorem.
\bibitem{GuthArea} L. Guth, \emph{Unexpected applications of polynomials in combinatorics}, expository notes, Proposition 2.5. \url{https://math.mit.edu/~lguth/Exposition/erdossurvey.pdf}.
\bibitem{DLMF} F. W. J. Olver et al., eds., \emph{NIST Digital Library of Mathematical Functions}, equations (5.9.8)--(5.9.9). \url{https://dlmf.nist.gov/5.9}.
\bibitem{BedratyukCovariants} L. Bedratyuk, \emph{Bivariate Poincar\'e series for algebra of covariants of a binary form}, arXiv:1006.1974v1, 2010. \url{https://arxiv.org/abs/1006.1974v1}.
\bibitem{PaperII} \emph{Complete affine jets and sharp integrability in three variables II. Degrees two through thirteen}, companion paper, 2026.
\bibitem{AKC80} G. I. Arkhipov, A. A. Karatsuba, and V. N. Chubarikov, \emph{Trigonometric integrals}, Math. USSR-Izv. \textbf{15} (1980), no.~2, 211--239. \href{https://doi.org/10.1070/IM1980v015n02ABEH001212}{doi:10.1070/IM1980v015n02ABEH001212}.
\bibitem{Ikr97} I. A. Ikromov, \emph{On the convergence exponent of trigonometric integrals}, Tr. Mat. Inst. Steklova \textbf{218} (1997), 179--189; English translation, Proc. Steklov Inst. Math. \textbf{218} (1997), 175--185. \url{https://www.mathnet.ru/eng/tm958}.
\bibitem{Ikr22} I. A. Ikromov, \emph{Letter to the Editor: Addition to the Paper ``On the Convergence Exponent of Trigonometric Integrals''}, Proc. Steklov Inst. Math. \textbf{319} (2022), 307--310. \href{https://doi.org/10.1134/S0081543822050200}{doi:10.1134/S0081543822050200}.
\bibitem{Parsell05} S. T. Parsell, \emph{A generalization of Vinogradov's mean value theorem}, Proc. London Math. Soc. (3) \textbf{91} (2005), no.~1, 1--32. \href{https://doi.org/10.1112/S002461150501525X}{doi:10.1112/S002461150501525X}.
\bibitem{PPW13} S. T. Parsell, S. M. Prendiville, and T. D. Wooley, \emph{Near-optimal mean value estimates for multidimensional Weyl sums}, Geom. Funct. Anal. \textbf{23} (2013), no.~6, 1962--2024. \href{https://doi.org/10.1007/s00039-013-0242-7}{doi:10.1007/s00039-013-0242-7}.
\bibitem{BD16} J. Bourgain and C. Demeter, \emph{Mean value estimates for Weyl sums in two dimensions}, J. London Math. Soc. (2) \textbf{94} (2016), no.~3, 814--838. \href{https://doi.org/10.1112/jlms/jdw063}{doi:10.1112/jlms/jdw063}.
\bibitem{BDGuo17} J. Bourgain, C. Demeter, and S. Guo, \emph{Sharp bounds for the cubic Parsell--Vinogradov system in two dimensions}, Adv. Math. \textbf{320} (2017), 827--875. \href{https://doi.org/10.1016/j.aim.2017.09.008}{doi:10.1016/j.aim.2017.09.008}.
\bibitem{GZ19} S. Guo and R. Zhang, \emph{On integer solutions of Parsell--Vinogradov systems}, Invent. Math. \textbf{218} (2019), no.~1, 1--81. \href{https://doi.org/10.1007/s00222-019-00881-6}{doi:10.1007/s00222-019-00881-6}.
\bibitem{PaperSupplement} \emph{Supplement to Complete affine jets and sharp integrability in three variables I and II}, 28 September 2026 revision. Distributed with this paper as ancillary material.
\end{thebibliography}
\enlargethispage{20pt}

\end{document}
